\chapter{Differentiation of the Spectral Doctrine}
\label{chap:differentiation-spectral-doctrine}

\section{Orientation: how does the spectral doctrine differentiate?}
\label{sec:differentiation-orientation}

Chapter~\Ref{chap:spectral-algebraic-geometry} constructed the native
spectral contexts used throughout the manuscript and organized their stable
coefficients into the indexed symmetric-monoidal doctrine
\[
    \QCoh:
    \bigl(\operatorname{Sch}^{\mathrm{sp}}\bigr)^{\op}
    \longrightarrow
    \CAlg(\Pr^L_{\mathrm{st}}),
\]
and established the substitution, Zariski descent, relative-spectrum
comprehension, and open--closed and linear context calculus needed below.
Those structures are fixed input. The present chapter asks what differential
structure is already carried by this spectral doctrine:
\[
    \boxed{\text{How does the spectral doctrine differentiate?}}
\]

The chapter answers this question through two complementary constructions.
The first is fibrewise and linear: the stable coefficient fibre is identified
affinely with the tangent stabilization of commutative spectral algebra and is
then equipped with the cofree cocommutative-coalgebra differential
exponential. The second is geometric: the cotangent complex represents
first-order variation, its degree-one classes construct square-zero
thickenings, and finite composites of those thickenings generate an
infinitesimal test geometry. These constructions are independent in origin.
They meet on free finite-projective models, where the derivative obtained from
cofree infinitesimal translation is compared with the universal spectral
K\"ahler derivation. The geometric branch then produces the endpoint
adjunctions, infinitesimal lifting and deformation calculus, and, after
stabilization, an exact fracture reconstruction.

Both branches use the same preparatory differential core:
\[
\boxed{
\text{spectral doctrine}
\longrightarrow
\text{affine tangent linearization}
\longrightarrow
\text{cotangent-classified variation and chain rule}.
}
\]
The tangent identification supplies the stable semantic environment. The
relative and global cotangent classifiers then organize variation by
substitution and transitivity. This cotangent calculus directly drives the
geometric branch and later supplies the comparison interface with the linear
branch; it does not construct the cofree exponential.

The fibrewise linear branch is tracked by
\[
    \QCoh(X),
    \qquad
    !_X,
    \qquad
    \eta,
    \qquad
    \mathsf d.
\]
Cofree cocommutative coalgebras construct the exponential and its canonical
monoidal structure. Primitive infinitesimal coalgebras then produce the
infinitesimal augmentation, codereliction, and deriving transformation. Only
after these maps have been constructed do finite-projective models compare
their convolution derivative with the spectral universal derivation.

The geometric branch is tracked by
\[
    \mathcal H_{\inf},
    \qquad
    \Re\dashv\Im\dashv\&.
\]
Degree-one cotangent classes first produce structured square-zero contexts.
Forgetting their presentations and closing under finite composition gives the
infinitesimal test category and its Zariski topos. Once the endpoint
adjunctions have been constructed, their modal form is
\[
    \Re\dashv\Im\dashv\&.
\]
Cotangent mapping spectra then compute the fibres of the associated lifting
comparison and control the first-order deformation and local-equation
problems.

For navigation, the order of the two constructions is:
\[
\boxed{
\begin{aligned}
\text{linear branch:}\qquad
&
\text{cofree cocommutative coalgebra exponential}
\\
&\longrightarrow
\text{primitive infinitesimal augmentation}
\\
&\longrightarrow
\text{codereliction and deriving transformation}
\\
&\longrightarrow
\text{comparison with the spectral universal derivation},
\\[0.8em]
\text{geometric branch:}\qquad
&
\text{structured square-zero context extension}
\\
&\longrightarrow
\text{finite infinitesimal contexts and their Zariski topos}
\\
&\longrightarrow
\text{differential cohesion }(\Re\dashv\Im\dashv\&)
\\
&\longrightarrow
\text{formal \(\acute{e}\)talification and infinitesimal disks}
\\
&\longrightarrow
\text{cotangent lifting, deformation, and equation calculus}
\\
&\longrightarrow
\text{stable endpoint semantics and BNV reconstruction}.
\end{aligned}}
\]
The display records construction order and later recognition statements.
Neither branch presents the other, and the BNV terminology in the final line
is only a forward reference to a comparison made after the stable fracture
theorem has been proved internally.

The cotangent classifier interfaces with both differential views:
\[
\begin{tikzcd}[column sep=small,row sep=large]
& L_{X/S}\in\QCoh(X)
    \arrow[dl, "{\text{linear classification}}"']
    \arrow[dr, "{\text{degree-one classes}}"]
& \\
(\QCoh(X),!_X,\mathsf d)
    \arrow[dr, "{\text{linear branch}}"']
&&
(X\hookrightarrow X_\eta)
    \arrow[dl, "{\text{geometric branch}}"]
\\
& \text{two differential views of the spectral doctrine} &
\end{tikzcd}
\]
This diagram is mnemonic rather than a dependency graph. In particular, its
left-hand arrow marks the later classification and comparison interface; it
does not construct the cofree coalgebra modality from the cotangent complex.

The chapter therefore proceeds from affine tangent linearization to the
relative and global cotangent calculus, completes the cofree linear branch,
and then returns to degree-one cotangent classes to construct the geometric
branch. Only after the infinitesimal site, endpoint semantics, and cotangent
lifting calculus are in place do we stabilize the endpoint adjunctions.
Construction precedes recognition throughout.

All conventions of Chapter~\Ref{chap:spectral-algebraic-geometry} remain in
force. Coordinate algebras and spectral schemes remain connective, whereas
module categories, cotangent complexes, mapping spectra, suspensions, fibres,
cofibres, and stabilizations are formed in their full stable ambient
categories. Coherence supplied functorially is suppressed from notation.
Affine charts, refinements, local equations, and factorizations chosen only
inside proofs remain proof witnesses rather than additional structure.

\begin{convention}[Relative and absolute notation]
\label{conv:deformation-relative-notation}
For a morphism of commutative ring spectra \(A\to B\), its relative cotangent
complex is written
\[
    L_{B/A}\in\Mod_B.
\]
For a morphism of spectral schemes \(f:X\to S\), its relative cotangent
complex is written
\[
    L_{X/S}\in\QCoh(X).
\]
When the base is the sphere spectrum or \(\Spec(\mathbb S)\), respectively,
we abbreviate
\[
    L_B:=L_{B/\mathbb S},
    \qquad
    L_X:=L_{X/\Spec(\mathbb S)}.
\]
The space of \(A\)-linear derivations from \(B\) to a \(B\)-module \(M\) is
written
\[
    \operatorname{Der}_A(B,M).
\]
No notation introduced in this chapter changes the variance conventions of
Chapter~\Ref{chap:spectral-algebraic-geometry}.
\end{convention}

\begin{convention}[Under--over categories]
\label{conv:deformation-under-over}
Let \(\mathcal C\) be an \(\infty\)-category and fix a morphism
\(A\to B\) in \(\mathcal C\). We write
\[
    \mathcal C_{A//B}:=(\mathcal C_{A/})_{/B}.
\]
Thus an object is a factorization
\[
    A\longrightarrow C\longrightarrow B
\]
of the fixed morphism, with its canonical coherent comparison. In
particular, \(\CAlg(\Sp)_{A//B}\) denotes commutative spectral algebras under
\(A\) and over \(B\), and \(\CAlg(\Sp)_{B//B}\) denotes retractive
commutative \(B\)-algebras. This notation is an ordinary slice construction
and introduces no additional structure.
\end{convention}

\begin{convention}[Theorem domains are properties]
\label{conv:deformation-theorem-domains}
Connectivity, finite local freeness, local finite presentation, nilpotence,
and representability by spectral schemes are predicates specifying theorem
domains in this chapter. They are not auxiliary fields of structure attached
to an object.
\end{convention}


\begin{convention}[Hooked arrows and geometric inclusions]
\label{conv:deformation-hooked-arrows}
A hooked arrow is used graphically for a specified geometric inclusion, in
particular for open immersions and, when convenient, for closed or
infinitesimal inclusions. Unless monomorphism is stated separately, the glyph
\(\hookrightarrow\) does not assert that the morphism is a monomorphism in the
ambient \(\infty\)-category. Thus a derived closed or square-zero immersion
may be displayed with a hook without suppressing its higher derived
self-intersection information.
\end{convention}

Foundational tangent and cotangent input is taken from
\cite{LurieHigherAlgebra,LurieDAGIV,LurieSAG}, while cofree spectral
cocommutative coalgebras are used as in
\cite{BachmannBurklund2026Coalgebras}. The differential-storage recognition
uses \cite{BluteCockettSeely2006Differential,Lemay2021Coderelictions},
with the linear/nonlinear and differential-lambda-calculus background of
\cite{Benton1995LNL,EhrhardRegnier2003Differential}. The endpoint
configuration is compared with differential cohesion only after its adjoint
quadruple has been constructed; at that point we use the notation
\[
    \Re\dashv\Im\dashv\&
\]
for reduction, infinitesimal shape, and infinitesimal flat. Likewise, the
stabilized endpoint semantics is first analyzed by the exact-adjunction
fracture theorem proved below and only afterward compared with
\cite[Sections~2--3]{BunkeNikolausVoelkl} and
\cite{HopkinsSinger2005}.

\section{Stabilization and the linear semantics of a model}
\label{sec:tangent-semantics}

Chapter~\Ref{chap:spectral-algebraic-geometry} attached the stable coefficient
fibre \(\Mod_B\) to an affine context \(\Spec(B)\). To differentiate the
doctrine, we first identify this fibre intrinsically as the linearization of
the nonlinear coordinate category at \(B\). This section establishes that
tangent identification and records the split square-zero object supplied by
the zeroth-space functor, which will make derivation terms visible in the next
section.

\subsection{The tangent category of commutative spectral algebra}

\begin{theorem}[Tangent category at a commutative ring spectrum]
\label{thm:tangent-category-calg}
Let \(B\in\CAlg(\Sp)\). There is a canonical equivalence of stable
\(\infty\)-categories
\[
    \operatorname{Stab}\bigl(\CAlg(\Sp)_{/B}\bigr)
    \simeq
    \Mod_B.
\]
Here \(\operatorname{Stab}\bigl(\CAlg(\Sp)_{/B}\bigr)\) denotes spectrum
objects in the pointed slice, equivalently the stabilization whose zeroth-space
functor lands in \(\bigl(\CAlg(\Sp)_{/B}\bigr)_*\).
Under this equivalence, the zeroth-space functor
\[
    \Omega^\infty
    :
    \operatorname{Stab}\bigl(\CAlg(\Sp)_{/B}\bigr)
    \longrightarrow
    \bigl(\CAlg(\Sp)_{/B}\bigr)_*
\]
identifies a \(B\)-module \(M\) with the split square-zero augmented algebra
\[
    B
    \longrightarrow
    B\oplus M
    \longrightarrow
    B.
\]
\end{theorem}

\begin{proof}
\leavevmode
\begin{enumerate}[label=\textup{(\arabic*)}]
    \item \textbf{Identify the resulting object.}
    The tangent theorem for commutative \(\mathbb E_\infty\)-algebras identifies the stabilization of the
    slice of commutative algebras at \(B\) with the category of modules over \(B\);
    see \cite{LurieHigherAlgebra,LurieDAGIV}.

    \item \textbf{Use terminality.}
    A pointed object of \(\CAlg(\Sp)_{/B}\) is a commutative algebra \(C\to B\) together with a section
    from the terminal object \(B\to B\), hence a diagram
    \[
        B\longrightarrow C\longrightarrow B
    \]
    whose composite is equivalent to \(\id_B\).

    \item \textbf{Transport the structure.}
    Under the tangent equivalence, the zeroth-space functor sends \(M\) to the split
    square-zero algebra whose underlying \(B\)-module is \(B\oplus M\) and whose
    multiplication on the \(M\)-summand is coherently null.
\end{enumerate}
\end{proof}

\begin{remark}[Linear semantics at a model]
\label{rem:linear-semantics-at-model}
The theorem identifies the stable coefficient fibre of an affine spectral
context with the linearization of the nonlinear model category at its
coordinate algebra:
\[
    \CAlg(\Sp)_{/B}
    \longrightarrow
    \operatorname{Stab}\bigl(\CAlg(\Sp)_{/B}\bigr)
    \simeq
    \Mod_B.
\]
Thus the module doctrine of Chapter~\Ref{chap:spectral-algebraic-geometry}
simultaneously supplies the tangent semantic fibres needed for deformation
theory.
\end{remark}

\subsection{Split square-zero algebra}

\begin{construction}[Split square-zero algebra]
\label{con:split-square-zero-algebra}
For \(B\in\CAlg(\Sp)\) and \(M\in\Mod_B\), write
\[
    B\oplus M
\]
for the commutative ring spectrum supplied by
\Ref{thm:tangent-category-calg}, with section and augmentation
\[
    B
    \xrightarrow{s}
    B\oplus M
    \xrightarrow{p}
    B,
    \qquad
    p\circ s\simeq\id_B.
\]
The underlying \(B\)-module is the direct sum \(B\oplus M\), and the product
of two elements in the coefficient summand is coherently zero.
\end{construction}

\begin{proposition}[Functoriality of split square-zero formation]
\label{prop:split-square-zero-functoriality}
For every commutative ring spectrum \(B\), split square-zero formation defines
a limit-preserving functor
\[
    \Mod_B
    \longrightarrow
    \CAlg(\Sp)_{B//B},
    \qquad
    M\longmapsto
    \bigl(B\to B\oplus M\to B\bigr).
\]
If \(B\) and \(M\) are connective, then \(B\oplus M\) is connective.
\end{proposition}

\begin{proof}
\leavevmode
\begin{enumerate}[label=\textup{(\arabic*)}]
    \item \textbf{Transport the structure.}
    Under \Ref{thm:tangent-category-calg}, the displayed functor is the zeroth-space functor of
    a stabilization.

    \item \textbf{Use the adjunction.}
    It is a right adjoint and therefore preserves limits.

    \item \textbf{Verify connectivity.}
    The underlying spectrum of the split square-zero algebra is \(B\oplus M\), so connectivity of
    both summands implies connectivity of the algebra.
\end{enumerate}
\end{proof}
The stable coefficient fibre has now been recovered from tangent
stabilization, and every coefficient determines a split square-zero augmented
algebra. What is still missing is the universal coefficient that classifies
sections of these augmentations.

\section{Derivation terms and cotangent classifiers}
\label{sec:derivation-cotangent}

Relative derivations are the terms represented by split square-zero
augmentations. This section first formulates those terms and then constructs
their representing object, the relative cotangent complex. The free-algebra
calculation at the end gives the basic model used repeatedly later.

\subsection{Relative derivation terms}

\begin{definition}[Spectral derivation]
\label{def:spectral-derivation}
Let \(A\to B\) be a morphism of commutative ring spectra and let
\(M\in\Mod_B\). The space of \(A\)-linear derivations from \(B\) to \(M\) is
\[
    \operatorname{Der}_A(B,M)
    :=
    \Map_{\CAlg(\Sp)_{A//B}}
    \bigl(B,B\oplus M\bigr).
\]
Thus an \(A\)-linear derivation is a section of the augmentation
\(B\oplus M\to B\) in commutative \(A\)-algebras.
\end{definition}

\begin{proposition}[Elementary functoriality of derivations]
\label{prop:derivation-elementary-functoriality}
Let \(A\to B\to C\) be morphisms of commutative ring spectra.
\begin{enumerate}[label=(\roman*)]
    \item A morphism of \(B\)-modules \(M\to M'\) induces
    \[
        \operatorname{Der}_A(B,M)
        \longrightarrow
        \operatorname{Der}_A(B,M').
    \]
    \item For every \(C\)-module \(N\), restriction along \(B\to C\) induces
    \[
        \operatorname{Der}_A(C,N)
        \longrightarrow
        \operatorname{Der}_A(B,N),
    \]
    where \(N\) is regarded on the right as a \(B\)-module.
\end{enumerate}
\end{proposition}

\begin{proof}
\leavevmode
\begin{enumerate}[label=\textup{(\arabic*)}]
    \item \textbf{Handle the first component.}
    The first map is postcomposition with the morphism of split square-zero algebras induced by
    \(M\to M'\).

    \item \textbf{Identify the resulting object.}
    For the second, compose an \(A\)-algebra section \(C\to C\oplus N\) with \(B\to C\) and use
    the canonical pullback identification
    \[
        B\times_C(C\oplus N)
        \simeq
        B\oplus N.
    \]

    \item \textbf{Verify naturality.}
    Naturality follows from functoriality of pullback and split square-zero formation.
\end{enumerate}
\end{proof}
Derivations are now identified as sections of split square-zero
augmentations. The next step replaces this family of mapping spaces by one
universal coefficient in the tangent fibre.

\subsection{The cotangent classifier}

\begin{theorem}[Existence of the relative cotangent classifier]
\label{thm:relative-cotangent-classifier}
Let \(A\to B\) be a morphism of commutative ring spectra. There exists a
\(B\)-module
\[
    L_{B/A}\in\Mod_B
\]
and a universal \(A\)-linear derivation such that, for every \(B\)-module
\(M\), there is a natural equivalence
\[
    \Map_{\Mod_B}(L_{B/A},M)
    \simeq
    \operatorname{Der}_A(B,M).
\]
The construction is functorial in commutative triangles of commutative ring
spectra.
\end{theorem}

\begin{proof}
\leavevmode
\begin{enumerate}[label=\textup{(\arabic*)}]
    \item \textbf{relative cotangent object of over.}
    The relative cotangent object of \(B\) over \(A\) is defined in the tangent
    bundle of \(\CAlg(\Sp)\).

    \item \textbf{Use the adjunction.}
    The tangent equivalence \Ref{thm:tangent-category-calg} identifies its fibre at \(B\)
    with \(\Mod_B\), and the stabilization adjunction gives the displayed representing
    equivalence.

    \item \textbf{Verify functoriality.}
    Functoriality is the functoriality of the relative cotangent object in the tangent bundle;
    see \cite{LurieHigherAlgebra,LurieDAGIV}.
\end{enumerate}
\end{proof}

\begin{corollary}[Cotangent classifier of an equivalence]
\label{cor:cotangent-equivalence-zero}
If \(A\to B\) is an equivalence, then
\[
    L_{B/A}\simeq0.
\]
\end{corollary}

\begin{proof}
\leavevmode
\begin{enumerate}[label=\textup{(\arabic*)}]
    \item \textbf{Identify contractibility.}
    Relative derivation spaces are contractible for every coefficient module, so their
    representing object is zero.
\end{enumerate}
\end{proof}

\begin{remark}[Classifier terminology]
\label{rem:cotangent-classifier-terminology}
The phrase \textbf{cotangent classifier} records the role of \(L_{B/A}\) in the
indexed semantics. It does not introduce a second object: throughout the
chapter the cotangent classifier is exactly the usual relative cotangent
complex.
\end{remark}
The cotangent classifier is defined abstractly by its representing property.
Its first concrete calculation is on free spectral coordinates, where the
classifier exposes the linear generators directly.

\subsection{Free models}

\begin{proposition}[Cotangent classifier of a free commutative algebra]
\label{prop:cotangent-free-algebra}
Let \(A\in\CAlg(\Sp)\), let \(M\in\Mod_A\), and put
\[
    B:=\operatorname{Sym}_A(M).
\]
Then there is a natural equivalence of \(B\)-modules
\[
    L_{B/A}
    \simeq
    B\otimes_A M.
\]
\end{proposition}

\begin{proof}
\leavevmode
\begin{enumerate}[label=\textup{(\arabic*)}]
    \item \textbf{Use the adjunction.}
    For every \(B\)-module \(N\), the free-algebra and extension-of-scalars
    adjunctions give
    \[
    \begin{aligned}
        \operatorname{Der}_A(B,N)
        &\simeq
        \Map_{\Mod_A}(M,N)
        \\
        &\simeq
        \Map_{\Mod_B}(B\otimes_A M,N).
    \end{aligned}
    \]

    \item \textbf{Apply Yoneda.}
    Compare with \Ref{thm:relative-cotangent-classifier} and apply Yoneda.
\end{enumerate}
\end{proof}

\begin{corollary}[Cotangent classifier of brave-new affine space]
\label{cor:cotangent-brave-new-affine-space}
For
\[
    A\{t_1,\ldots,t_n\}
    :=
    \operatorname{Sym}_A(A^{\oplus n}),
\]
one has
\[
    L_{A\{t_1,\ldots,t_n\}/A}
    \simeq
    A\{t_1,\ldots,t_n\}^{\oplus n}.
\]
Equivalently, the affine context \(\mathbb A^n_{\Spec A}\) of
\Ref{def:brave-new-affine-space} has free rank-\(n\) relative cotangent
classifier.
\end{corollary}
We now have the relative cotangent classifier and its free-model
calculation. For it to serve as the differential structure of an indexed
geometry, it must also behave coherently under change of context and under
composition.

\section{Substitution, the chain rule, and intrinsic cotangent estimates}
\label{sec:cotangent-substitution-transitivity}

This section supplies that calculus. Cotangent base change is the
substitution law, while transitivity is the chain rule. The subsequent
connectivity, Hurewicz, rigidity, and finiteness results record the intrinsic
estimates needed when cotangent data are later converted into geometric
lifting and equation problems.

\subsection{Substitution and transitivity}

\begin{theorem}[Substitution law for cotangent classifiers]
\label{thm:cotangent-base-change}
Let
\[
\begin{tikzcd}[column sep=large, row sep=large]
    A \arrow[r] \arrow[d]
        & B \arrow[d]
    \\
    A' \arrow[r]
        & B':=B\otimes_AA'
\end{tikzcd}
\]
be a pushout square of commutative ring spectra. Then there is a natural
equivalence of \(B'\)-modules
\[
    L_{B/A}\otimes_BB'
    \simeq
    L_{B'/A'}.
\]
\end{theorem}

\begin{proof}
\leavevmode
\begin{enumerate}[label=\textup{(\arabic*)}]
    \item \textbf{For every -module.}
    For every \(B'\)-module \(N\),
    \[
    \begin{aligned}
        \Map_{\Mod_{B'}}(L_{B/A}\otimes_BB',N)
        &\simeq
        \Map_{\Mod_B}(L_{B/A},N)
        \\
        &\simeq
        \operatorname{Der}_A(B,N)
        \\
        &\simeq
        \operatorname{Der}_{A'}(B',N)
        \\
        &\simeq
        \Map_{\Mod_{B'}}(L_{B'/A'},N).
    \end{aligned}
    \]

    \item \textbf{Apply the universal property.}
    The middle equivalence is the pushout universal property.

    \item \textbf{Apply Yoneda.}
    Yoneda gives the result.
\end{enumerate}
\end{proof}

\begin{corollary}[Substitution of derivation terms]
\label{cor:derivation-base-change}
In the situation of \Ref{thm:cotangent-base-change}, every \(B'\)-module \(N\)
admits a natural equivalence
\[
    \operatorname{Der}_{A'}(B',N)
    \simeq
    \operatorname{Der}_A(B,N).
\]
\end{corollary}

\begin{theorem}[Chain rule for cotangent classifiers]
\label{thm:cotangent-transitivity}
For composable morphisms
\[
    A\longrightarrow B\longrightarrow C
\]
of commutative ring spectra, there is a natural cofibre sequence of
\(C\)-modules
\[
    L_{B/A}\otimes_BC
    \longrightarrow
    L_{C/A}
    \longrightarrow
    L_{C/B}.
\]
\end{theorem}

\begin{proof}
\leavevmode
\begin{enumerate}[label=\textup{(\arabic*)}]
    \item \textbf{Fix the notation.}
    Let \(N\in\Mod_C\).

    \item \textbf{Compute the fibre.}
    Restriction of an \(A\)-linear derivation of \(C\) along \(B\to C\) gives a
    natural fibre sequence of derivation spaces
    \[
        \operatorname{Der}_B(C,N)
        \longrightarrow
        \operatorname{Der}_A(C,N)
        \longrightarrow
        \operatorname{Der}_A(B,N),
    \]
    where \(N\) is regarded in the last term as a \(B\)-module by restriction of
    scalars.

    \item \textbf{Compute the fibre.}
    The homotopy fibre over the zero restricted derivation consists exactly of those
    \(A\)-linear derivations of \(C\) which vanish on \(B\), that is, the
    \(B\)-linear derivations.

    \item \textbf{Use the adjunction.}
    By the cotangent representing property and extension--restriction adjunction, the three
    spaces are naturally
    \[
    \begin{aligned}
        \operatorname{Der}_B(C,N)
        &\simeq
        \Map_{\Mod_C}(L_{C/B},N),
        \\
        \operatorname{Der}_A(C,N)
        &\simeq
        \Map_{\Mod_C}(L_{C/A},N),
        \\
        \operatorname{Der}_A(B,N)
        &\simeq
        \Map_{\Mod_C}(L_{B/A}\otimes_BC,N).
    \end{aligned}
    \]

    \item \textbf{Transport the same comparison.}
    The same equivalences hold naturally after replacing \(N\) by every shift \(N[k]\).

    \item \textbf{Compute the mapping object.}
    Since the mapping spectrum in a stable \(\infty\)-category is determined by these shifted
    mapping spaces, the preceding fibre sequences assemble to a natural fibre sequence of
    mapping spectra
    \[
    \begin{aligned}
        \map_{\Mod_C}(L_{C/B},N)
        \longrightarrow
        \map_{\Mod_C}(L_{C/A},N)
        \longrightarrow
        \map_{\Mod_C}(L_{B/A}\otimes_BC,N).
    \end{aligned}
    \]

    \item \textbf{Identify the resulting comparison.}
    This holds naturally for every \(C\)-module \(N\).

    \item \textbf{Apply Yoneda.}
    Stable Yoneda therefore identifies
    \[
        L_{B/A}\otimes_BC
        \longrightarrow
        L_{C/A}
        \longrightarrow
        L_{C/B}
    \]
    as a cofibre sequence in \(\Mod_C\).
\end{enumerate}
\end{proof}

\begin{lemma}[Spectral localizations have trivial relative cotangent]
\label{lem:spectral-localization-cotangent-zero}
Let \(A\in\CAlg(\Sp)\) and let \(f\in\pi_0(A)\). Then
\[
    L_{A[f^{-1}]/A}\simeq0.
\]
The equivalence is natural under further localization.
\end{lemma}

\begin{proof}
\leavevmode
\begin{enumerate}[label=\textup{(\arabic*)}]
    \item \textbf{Fix the notation.}
    Put \(A_f:=A[f^{-1}]\).

    \item \textbf{Identify the resulting object.}
    For every \(A_f\)-module \(N\), the derivation classifier identifies
    \[
        \Map_{\Mod_{A_f}}(L_{A_f/A},N)
        \simeq
        \operatorname{Der}_A(A_f,N).
    \]

    \item \textbf{Compute the fibre.}
    The right-hand side is the homotopy fibre over \(\id_{A_f}\) of
    \[
        \Map_{\CAlg(\Sp)_{A/}}(A_f,A_f\oplus N)
        \longrightarrow
        \Map_{\CAlg(\Sp)_{A/}}(A_f,A_f),
    \]
    where the map is induced by the split augmentation \(A_f\oplus N\to A_f\).

    \item \textbf{Verify invertibility.}
    In both target algebras the image of \(f\) is invertible.

    \item \textbf{Apply the universal property.}
    The universal property of spectral localization therefore makes both displayed mapping
    spaces contractible, and the augmentation-induced map between them is an equivalence.

    \item \textbf{Conclude the step.}
    Hence
    \[
        \operatorname{Der}_A(A_f,N)\simeq *
    \]
    for every \(N\).

    \item \textbf{Apply Yoneda.}
    The representing property and Yoneda give \(L_{A_f/A}\simeq0\).

    \item \textbf{Apply the universal property.}
    Naturality is inherited from the functorial universal property of localization.
\end{enumerate}
\end{proof}

\begin{proposition}[Affine open immersions have trivial relative cotangent]
\label{prop:affine-open-cotangent-zero}
Let
\[
    j:\Spec(B)\hookrightarrow\Spec(A)
\]
be an affine spectral Zariski open immersion. Then
\[
    L_{B/A}\simeq0.
\]
\end{proposition}

\begin{proof}
\leavevmode
\begin{enumerate}[label=\textup{(\arabic*)}]
    \item \textbf{Construct the factorization.}
    Choose a finite family of principal spectral opens
    \[
        D_A^{\mathrm{sp}}(f_\alpha)=\Spec(A[f_\alpha^{-1}])
        \longrightarrow
        \Spec(A)
    \]
    which factor through \(\Spec(B)\) and whose pullbacks cover \(\Spec(B)\); such a family exists
    by the principal-basis theorem of Chapter~\Ref{chap:spectral-algebraic-geometry}.

    \item \textbf{Fix the notation.}
    Put
    \[
        B_\alpha:=B\otimes_AA[f_\alpha^{-1}].
    \]

    \item \textbf{Construct the factorization.}
    Because the principal open factors through \(\Spec(B)\), the canonical map
    \[
        B_\alpha\longrightarrow A[f_\alpha^{-1}]
    \]
    is an equivalence.

    \item \textbf{Apply cotangent transitivity.}
    By \Ref{lem:spectral-localization-cotangent-zero}, principal spectral localization has
    vanishing relative cotangent complex, so transitivity for \(A\to B\to B_\alpha\) gives
    \[
        L_{B/A}\otimes_BB_\alpha\simeq0.
    \]

    \item \textbf{Use joint conservativity.}
    The substitution functors \(-\otimes_BB_\alpha\) are jointly conservative for a finite affine Zariski
    cover by Chapter~\Ref{chap:spectral-algebraic-geometry}.

    \item \textbf{Conclude the step.}
    Therefore \(L_{B/A}\simeq0\).
\end{enumerate}
\end{proof}

\begin{remark}[Substitution and composition]
\label{rem:cotangent-doctrinal-interpretation}
Theorem~\Ref{thm:cotangent-base-change} is the substitution law for the
differentiated spectral doctrine. Theorem~\Ref{thm:cotangent-transitivity}
is its chain rule: variation along a composite context map is assembled from
variation in the first stage and variation relative to that stage. No further exchange
law is imposed; these morphisms and equivalences are the canonical ones
constructed by the representing property.
\end{remark}
Substitution and the chain rule describe how cotangent classifiers interact
with context change. We now compare the classifier with the underlying
relative homotopy of a map; this is the bridge to conormal and rigidity
arguments.

\subsection{Connectivity and relative Hurewicz comparison}

\begin{proposition}[Connectivity of the relative cotangent classifier]
\label{prop:cotangent-connective}
Let \(A\to B\) be a morphism of connective commutative ring spectra. Then
\[
    L_{B/A}\in(\Mod_B)_{\geq0}.
\]
\end{proposition}

\begin{proof}
\leavevmode
\begin{enumerate}[label=\textup{(\arabic*)}]
    \item \textbf{Verify connectivity.}
    This is the connective relative cotangent theorem for commutative ring spectra; see
    \cite{LurieHigherAlgebra}.

    \item \textbf{Verify connectivity.}
    Equivalently, resolve \(B\) over \(A\) by connective free commutative algebras,
    apply \Ref{prop:cotangent-free-algebra}, and use compatibility with geometric realization.
\end{enumerate}
\end{proof}

\begin{construction}[Canonical relative Hurewicz morphism]
\label{con:relative-hurewicz-morphism}
Let
\[
    A\longrightarrow B
\]
be a morphism of commutative ring spectra and put
\[
    Q:=\operatorname{cofib}(A\to B)
\]
in \(\Mod_A\). For a \(B\)-module \(N\), let
\[
    d:B\longrightarrow B\oplus N
\]
be an \(A\)-linear derivation and let
\[
    d_0:B\longrightarrow B\oplus N
\]
be the zero derivation. On underlying \(A\)-modules, the difference
\(d-d_0\) has zero component in the \(B\)-summand and therefore determines
an \(A\)-linear morphism
\[
    \delta_d:B\longrightarrow N.
\]
Because \(d\) and \(d_0\) agree after restriction to \(A\), the composite
\(A\to B\to N\) is canonically null and \(\delta_d\) factors through a
canonical morphism
\[
    \bar\delta_d:Q\longrightarrow N.
\]
Adjunction gives a \(B\)-linear morphism
\[
    B\otimes_AQ\longrightarrow N.
\]
The construction is natural in \(N\). Using
\[
    \operatorname{Der}_A(B,N)
    \simeq
    \Map_{\Mod_B}(L_{B/A},N),
\]
Yoneda therefore determines a canonical natural morphism
\[
    \epsilon_{B/A}:
    B\otimes_AQ
    \longrightarrow
    L_{B/A}.
\]
We call \(\epsilon_{B/A}\) the \textbf{relative Hurewicz morphism}. It is
characterized by the property that, for every classifying map
\(\eta:L_{B/A}\to N\), the adjoint of
\(\eta\epsilon_{B/A}\) is the factor
\(Q\to N\) of the underlying relative derivation.

To compare with the standard construction, let
\[
    L_B\longrightarrow L_{B/A}
\]
be the transitivity morphism. Its associated square-zero extension has
fibre \(L_{B/A}[-1]\), and the induced boundary
\(Q\to L_{B/A}\) is, by the preceding characterization, the adjoint of
\(\epsilon_{B/A}\). This is exactly the construction of
\cite[Construction~7.4.3.10]{LurieHigherAlgebra}. Thus the morphism above
is the standard canonical Hurewicz morphism.
\end{construction}

\begin{theorem}[Relative Hurewicz--conormal comparison]
\label{thm:relative-hurewicz-conormal}
Let \(f:A\to B\) be a morphism of connective commutative ring spectra which is
surjective on \(\pi_0\), and put
\[
    K:=\operatorname{fib}(A\to B).
\]
Then \(K\) is connective, \(L_{B/A}\) is \(1\)-connective, and shifting the
Hurewicz morphism of
\Ref{con:relative-hurewicz-morphism} gives a canonical natural morphism
\[
    B\otimes_AK
    \longrightarrow
    L_{B/A}[-1]
\]
which is an isomorphism on \(\pi_0\).
More generally, if
\[
    Q:=\operatorname{cofib}(A\to B)
\]
is \(n\)-connective for \(n\geq1\), then the fibre of
\[
    \epsilon_{B/A}:B\otimes_AQ\longrightarrow L_{B/A}
\]
is \(2n\)-connective. In particular, the Hurewicz morphism induces an
isomorphism on the first possible nonzero homotopy group.
\end{theorem}

\begin{proof}
\leavevmode
\begin{enumerate}[label=\textup{(\arabic*)}]
    \item \textbf{Verify connectivity.}
    Surjectivity of \(\pi_0(A)\to\pi_0(B)\) makes the fibre \(K\) connective and gives a canonical
    equivalence
    \[
        Q=\operatorname{cofib}(A\to B)\simeq K[1].
    \]

    \item \textbf{Verify connectivity.}
    The connectivity estimate for the canonical Hurewicz morphism is
    \cite[Theorem~7.4.3.12]{LurieHigherAlgebra}.

    \item \textbf{Verify connectivity.}
    Applied with \(n=1\), it makes \(L_{B/A}\) \(1\)-connective and the fibre of
    \(B\otimes_AQ\to L_{B/A}\) \(2\)-connective.

    \item \textbf{Derive the comparison.}
    Shifting by \([-1]\) gives the displayed conormal comparison and its
    \(\pi_0\)-isomorphism.

    \item \textbf{Transport the same comparison.}
    The same theorem for arbitrary \(n\geq1\) gives the final statement.
\end{enumerate}
\end{proof}

\begin{lemma}[Degree-zero quotient of the derived conormal module]
\label{lem:derived-conormal-classical-quotient}
In the situation of \Ref{thm:relative-hurewicz-conormal}, put
\[
    R:=\pi_0(A),
    \qquad
    S:=\pi_0(B),
    \qquad
    I:=\ker(R\to S).
\]
Then the relative fibre determines a canonical surjection
\[
    \pi_0(B\otimes_AK)
    \twoheadrightarrow
    I/I^2.
\]
Under the Hurewicz isomorphism of
\Ref{thm:relative-hurewicz-conormal}, this is a canonical quotient of
\(\pi_0(L_{B/A}[-1])\) onto the ordinary conormal module \(I/I^2\).
\end{lemma}

\begin{proof}
\leavevmode
\begin{enumerate}[label=\textup{(\arabic*)}]
    \item \textbf{Verify surjectivity.}
    The long exact homotopy sequence of \(K\to A\to B\) shows that the image of \(\pi_0(K)\to R\) is
    exactly \(I\), so there is a surjection
    \[
        \pi_0(K)\twoheadrightarrow I.
    \]

    \item \textbf{Verify connectivity.}
    Because \(B\) and \(K\) are connective, the connective tensor-product spectral
    sequence gives
    \[
        \pi_0(B\otimes_AK)
        \simeq
        S\otimes_R\pi_0(K).
    \]

    \item \textbf{Verify surjectivity.}
    Tensoring the preceding surjection with \(S=R/I\) gives
    \[
        S\otimes_R\pi_0(K)
        \twoheadrightarrow
        S\otimes_RI
        \simeq
        I/I^2.
    \]

    \item \textbf{Verify naturality.}
    Naturality follows from naturality of the fibre sequence and scalar extension.
\end{enumerate}
\end{proof}

\begin{corollary}[Cotangent rigidity]
\label{cor:cotangent-rigidity}
Let \(A\to B\) be a morphism of connective commutative ring spectra. If
\[
    \pi_0(A)\xrightarrow{\cong}\pi_0(B)
    \qquad\text{and}\qquad
    L_{B/A}\simeq0,
\]
then \(A\to B\) is an equivalence.
\end{corollary}

\begin{proof}
\leavevmode
\begin{enumerate}[label=\textup{(\arabic*)}]
    \item \textbf{Identify the resulting comparison.}
    This is the cotangent rigidity consequence of the relative Hurewicz theorem; see
    \cite[Corollary~7.4.3.4]{LurieHigherAlgebra}.

    \item \textbf{Identify the map.}
    The proof there applies the first-nonzero-relative-homotopy comparison supplied by the
    Hurewicz morphism.

    \item \textbf{Verify connectivity.}
    The hypotheses displayed here are exactly its connective theorem domain.
\end{enumerate}
\end{proof}
The Hurewicz comparison controls the first nonzero relative homotopy and yields
cotangent rigidity. A second intrinsic use of the same classifier is to
detect local finite presentation.

\subsection{Finiteness detected by cotangent data}

\begin{theorem}[Cotangent criterion for local finite presentation]
\label{thm:cotangent-finite-presentation}
Let \(A\to B\) be a morphism of connective commutative ring spectra.
\begin{enumerate}[label=(\roman*)]
    \item If \(B\) is locally of finite presentation over \(A\), then
    \(L_{B/A}\) is a perfect \(B\)-module.
    \item Conversely, if \(L_{B/A}\) is perfect and \(\pi_0(B)\) is finitely
    presented as a \(\pi_0(A)\)-algebra, then \(B\) is locally of finite
    presentation over \(A\).
\end{enumerate}
\end{theorem}

\begin{proof}
\leavevmode
\begin{enumerate}[label=\textup{(\arabic*)}]
    \item \textbf{Verify connectivity.}
    This is the connective spectral cotangent criterion for local finite presentation; see
    \cite{LurieDAGIV,LurieHigherAlgebra,LurieSAG}.

    \item \textbf{theorem is invoked here as foundational spectral.}
    The theorem is invoked here as foundational spectral algebra rather than reproved by a
    parallel cell theory.
\end{enumerate}
\end{proof}
These results give the affine cotangent interface needed later: substitution
and the chain rule control variation, the Hurewicz comparison controls the
first relative homotopy, and cotangent data detect the stated finiteness
property. The remaining problem is globalization, because the subsequent
geometric constructions require one cotangent object on an arbitrary spectral
scheme rather than a collection of affine choices.

\section{The intrinsic global cotangent doctrine}
\label{sec:global-cotangent-doctrine}

We globalize the affine classifier using only structures already constructed in
Chapter~\Ref{chap:spectral-algebraic-geometry}. Compatible affine cotangent
complexes are transported by substitution and glued by spectral Zariski
descent in \(\QCoh\). The structured ringed-topos cotangent complex is used
only afterward as a recognition of the resulting intrinsic object.

In the globalizations below, once affine data and comparison maps have been
shown compatible on common refinements, full faithfulness of spectral Zariski
descent is used to glue them uniquely. The affine covers and refinements used
for this verification remain proof witnesses.

\subsection{Construction by affine coefficient descent}

\begin{definition}[Compatible affine pair for a relative context]
\label{def:compatible-affine-pair-cotangent}
Let \(f:X\to S\) be a morphism of spectral schemes. A \textbf{compatible
affine pair} for \(f\) is a diagram
\[
\begin{tikzcd}[column sep=large,row sep=large]
    U=\Spec(B) \arrow[r,hook] \arrow[d]
        & X \arrow[d,"f"]
    \\
    V=\Spec(A) \arrow[r,hook]
        & S
\end{tikzcd}
\]
in which the horizontal arrows are affine open immersions. A refinement is a
commutative diagram of such pairs induced by affine open restriction in the
source and target.
\end{definition}

\begin{lemma}[Compatibility of affine cotangent classifiers]
\label{lem:affine-cotangent-compatible-pairs}
Let \((U,V)\) and \((U',V')\) be compatible affine pairs for
\(f:X\to S\), with
\[
    U'=\Spec(B')\longrightarrow U=\Spec(B),
    \qquad
    V'=\Spec(A')\longrightarrow V=\Spec(A)
\]
a refinement. Then there is a canonical equivalence
\[
    L_{B/A}\otimes_BB'
    \xrightarrow{\simeq}
    L_{B'/A'}.
\]
These equivalences are coherent under identities, composition, and iterated
refinement.
\end{lemma}

\begin{proof}
\leavevmode
\begin{enumerate}[label=\textup{(\arabic*)}]
    \item \textbf{square is a pushout.}
    The square
    \[
    \begin{tikzcd}[column sep=large,row sep=large]
        A \arrow[r] \arrow[d]
            & B \arrow[d]
        \\
        A' \arrow[r]
            & B\otimes_AA'
    \end{tikzcd}
    \]
    is a pushout.

    \item \textbf{Apply cotangent base change.}
    Cotangent base change gives
    \[
        L_{B/A}\otimes_B(B\otimes_AA')
        \simeq
        L_{(B\otimes_AA')/A'}.
    \]

    \item \textbf{Identify the map.}
    The further morphism \(
B\otimes_AA'\to B'
\) represents an affine open immersion, so its relative
    cotangent complex vanishes.

    \item \textbf{Apply cotangent transitivity.}
    Transitivity therefore identifies
    \[
        L_{(B\otimes_AA')/A'}\otimes_{B\otimes_AA'}B'
        \simeq
        L_{B'/A'}.
    \]

    \item \textbf{Derive the comparison.}
    Composing yields the displayed equivalence.

    \item \textbf{Apply cotangent base change.}
    All comparison maps are obtained from functorial cotangent base change and transitivity, so
    their higher compatibilities are inherited from those constructions.
\end{enumerate}
\end{proof}

\begin{construction}[Native global relative cotangent classifier]
\label{con:native-global-cotangent-classifier}
Let \(f:X\to S\) be a morphism of spectral schemes. On every compatible
affine pair
\[
    U=\Spec(B)\longrightarrow V=\Spec(A)
\]
place the quasi-coherent coefficient corresponding to \(L_{B/A}\). By
\Ref{lem:affine-cotangent-compatible-pairs}, these coefficients agree
canonically on every common affine refinement, with all higher compatibility.
These source charts form a Zariski basis of \(X\): given an affine open of
\(X\), pull back an affine open cover of \(S\) and refine the resulting
cover by affine principal opens of the source. The quasi-coherent descent
theorem of
Chapter~\Ref{chap:spectral-algebraic-geometry} produces a canonically
determined object
\[
    L_{X/S}\in\QCoh(X).
\]
We call it the \textbf{global relative cotangent classifier}.
\end{construction}

\begin{theorem}[Affine characterization and uniqueness]
\label{thm:global-cotangent-classifier}
For every morphism \(f:X\to S\), the object of
\Ref{con:native-global-cotangent-classifier} is characterized by the following
property: whenever
\[
    U=\Spec(B)\longrightarrow X
\]
is an affine open whose composite with \(S\) factors through an affine open
\(V=\Spec(A)\to S\), there is a canonical equivalence
\[
    L_{X/S}|_U
    \simeq
    L_{B/A},
\]
compatible with every further affine restriction. The space of descent gluings of this specified compatible affine cotangent
datum is contractible. Moreover
\[
    L_{X/S}\in\QCoh(X)_{\geq0}.
\]
\end{theorem}

\begin{proof}
\leavevmode
\begin{enumerate}[label=\textup{(\arabic*)}]
    \item \textbf{Set up the argument.}
    Existence is the descent construction above.

    \item \textbf{Use full faithfulness.}
    Full faithfulness of spectral Zariski descent makes the space of descent gluings of the
    fixed affine datum contractible and proves uniqueness up to its canonical contractible
    comparison space.

    \item \textbf{Verify connectivity.}
    Each affine restriction \(L_{B/A}\) is connective by \Ref{prop:cotangent-connective};
    connectivity in \(\QCoh(X)\) is detected on affine opens, so \(L_{X/S}\) is connective.
\end{enumerate}
\end{proof}

\begin{remark}[Recognition in the structured spectral model]
\label{rem:global-cotangent-intrinsic}
Under the comparison of
Chapter~\Ref{chap:spectral-algebraic-geometry} with structured spectral
schemes, the intrinsically descended object \(L_{X/S}\) identifies with the
sheaf-level cotangent complex
\[
    L_{\mathcal O_X/f^{-1}\mathcal O_S}
\]
of \cite{LurieSAG}. This comparison recognizes an object already constructed
in the native coefficient doctrine. No small ringed topos, affine atlas,
refinement, or independent descent datum is retained in the definition of
\(L_{X/S}\).
\end{remark}
The global object has now been constructed by affine coefficient descent. Its
first structural test is whether the affine substitution law survives that
gluing and hence defines genuine change of context globally.

\subsection{Global substitution and transitivity}

\begin{theorem}[Global cotangent substitution law]
\label{thm:global-cotangent-base-change}
Let
\[
\begin{tikzcd}[column sep=large, row sep=large]
    X' \arrow[r,"g'"] \arrow[d,"f'"']
        & X \arrow[d,"f"]
    \\
    S' \arrow[r,"g"']
        & S
\end{tikzcd}
\]
be Cartesian in spectral schemes. Then there is a canonical equivalence
\[
    L_{X'/S'}
    \simeq
    (g')^*L_{X/S}.
\]
The equivalence is coherent under identities, composition, and iterated
Cartesian base change.
\end{theorem}

\begin{proof}
\leavevmode
\begin{enumerate}[label=\textup{(\arabic*)}]
    \item \textbf{Choose the local data.}
    Choose a compatible affine pair
    \[
        U=\Spec(B)\longrightarrow V=\Spec(A)
    \]
    for \(X\to S\).

    \item \textbf{After pullback to , choose an affine.}
    After pullback to \(S'\), choose an affine open \(V'=\Spec(A')\to S'\) mapping to \(V\),
    form
    \[
        \bar B:=B\otimes_AA',
    \]
    and refine the corresponding source open by an affine open \(U'=\Spec(B')\to\Spec(\bar B)\).

    \item \textbf{Identify the map.}
    The latter map is an affine spectral Zariski open immersion.

    \item \textbf{affine comparison on this refinement.}
    The affine comparison on this refinement is therefore
    \[
    \begin{aligned}
        L_{B/A}\otimes_BB'
        &\simeq
        \bigl(L_{B/A}\otimes_B\bar B\bigr)
           \otimes_{\bar B}B'\\
        &\simeq
        L_{\bar B/A'}\otimes_{\bar B}B'\\
        &\simeq
        L_{B'/A'}.
    \end{aligned}
    \]

    \item \textbf{Apply cotangent base change.}
    The middle equivalence is cotangent base change and the last is transitivity together with
    \(L_{B'/\bar B}\simeq0\).

    \item \textbf{this is precisely the compatible-pair comparison of.}
    Equivalently, this is precisely the compatible-pair comparison of
    \Ref{lem:affine-cotangent-compatible-pairs}.

    \item \textbf{Verify compatibility.}
    These affine comparisons commute with every further compatible refinement by the coherence
    established in that lemma.

    \item \textbf{Use full faithfulness.}
    Full faithfulness of quasi-coherent Zariski descent therefore glues them uniquely to
    \[
        (g')^*L_{X/S}\xrightarrow{\simeq}L_{X'/S'}.
    \]

    \item \textbf{Apply base change.}
    Identity, composition, and iterated-base-change coherence are inherited from the affine
    cotangent base-change and transitivity constructions and are detected on the same affine
    basis.
\end{enumerate}
\end{proof}

\begin{theorem}[Global cotangent chain rule]
\label{thm:global-cotangent-transitivity}
For composable morphisms of spectral schemes
\[
    X\xrightarrow{f}Y\xrightarrow{g}S,
\]
there is a canonical cofibre sequence in \(\QCoh(X)\)
\[
    f^*L_{Y/S}
    \longrightarrow
    L_{X/S}
    \longrightarrow
    L_{X/Y}.
\]
\end{theorem}

\begin{proof}
\leavevmode
\begin{enumerate}[label=\textup{(\arabic*)}]
    \item \textbf{Restrict the construction.}
    Restrict to an affine open of \(X\), refine its image in \(Y\) and \(S\)
    by compatible affine opens, and apply \Ref{thm:cotangent-transitivity}.

    \item \textbf{Compute the fibre.}
    The affine cofibre sequences and their maps are compatible under further restriction.

    \item \textbf{Apply Zariski descent.}
    Since finite limits and colimits in the stable category \(\QCoh(X)\) are detected by affine
    Zariski descent, the local sequences glue to the asserted global cofibre sequence.
\end{enumerate}
\end{proof}

\begin{corollary}[Open context extensions have trivial relative tangent]
\label{cor:cotangent-open-zero}
If \(j:U\hookrightarrow X\) is an open immersion of spectral schemes, then
\[
    L_{U/X}\simeq0.
\]
\end{corollary}

\begin{proof}
\leavevmode
\begin{enumerate}[label=\textup{(\arabic*)}]
    \item \textbf{Set up the argument.}
    The assertion is local on source and target.

    \item \textbf{After restricting to compatible affine opens it.}
    After restricting to compatible affine opens it is exactly
    \Ref{prop:affine-open-cotangent-zero}.

    \item \textbf{Glue the local data.}
    The affine vanishings glue uniquely by \Ref{thm:global-cotangent-classifier}.
\end{enumerate}
\end{proof}
With global substitution and the chain rule in place, the global cotangent
object can be used exactly as in the affine case: mapping out of it defines
relative derivation terms on a spectral context.

\subsection{Global derivation terms}

\begin{definition}[Relative derivation term on a spectral context]
\label{def:global-derivation-term}
Let \(f:X\to S\) be a spectral scheme and let \(M\in\QCoh(X)\). Define the
space of relative derivation terms with coefficient \(M\) by
\[
    \operatorname{Der}_S(X;M)
    :=
    \Map_{\QCoh(X)}(L_{X/S},M).
\]
The universal element corresponding to \(\id_{L_{X/S}}\) is called the
universal relative \(S\)-derivation term.
\end{definition}

\begin{proposition}[Affine computation of global derivation terms]
\label{prop:global-derivation-affine-computation}
Let \(X\to S\) and \(M\) be as above. Choose an affine spectral Zariski atlas
\(U\to X\), let \(U_\bullet\to X\) be its \v{C}ech nerve, and on each
\(U_n\) choose an affine-open basis subordinate to compatible affine opens of
\(S\). Then
\[
    \operatorname{Der}_S(X;M)
    \simeq
    \operatorname*{Tot}_{[n]\in\Delta}
    \Map_{\QCoh(U_n)}
    \bigl(L_{X/S}|_{U_n},M|_{U_n}\bigr).
\]
Each mapping space on the right is in turn computed on the chosen compatible
affine-open basis by \Ref{thm:relative-cotangent-classifier}. The resulting
equivalence is independent of the computational atlas and all affine
refinements.
\end{proposition}

\begin{proof}
\leavevmode
\begin{enumerate}[label=\textup{(\arabic*)}]
    \item \textbf{Identify the resulting object.}
    Chapter~\Ref{chap:spectral-algebraic-geometry} identifies \(\QCoh(X)\) with the limit of its
    values on the ordinary \v{C}ech nerve of an affine Zariski atlas; no hypercompletion is
    used.

    \item \textbf{Compute the mapping object.}
    Mapping spaces in a limit of \(\infty\)-categories are limits of mapping spaces.

    \item \textbf{Conclude the step.}
    Therefore
    \[
    \begin{aligned}
        \Map_{\QCoh(X)}(L_{X/S},M)
        &\simeq
        \operatorname*{Tot}_{[n]\in\Delta}
        \Map_{\QCoh(U_n)}
        \bigl(L_{X/S}|_{U_n},M|_{U_n}\bigr).
    \end{aligned}
    \]

    \item \textbf{Compute the mapping object.}
    On every compatible affine open \(\Spec(B)\to\Spec(A)\), \Ref{thm:global-cotangent-classifier}
    identifies the restricted cotangent object with \(L_{B/A}\), and the affine mapping space is
    computed by \Ref{thm:relative-cotangent-classifier}.

    \item \textbf{Use full faithfulness.}
    Common affine refinements and full faithfulness of Zariski descent prove independence of all
    proof choices.
\end{enumerate}
\end{proof}
The preceding results can be compressed into an indexed statement. Rather
than treating each base-change equivalence separately, we regard the global
cotangent classifier as a section over the category of relative contexts.

\subsection{The cotangent classifier as a Cartesian section}

Let
\[
    \operatorname{Rel}^{\mathrm{sp}}_{\mathrm{Cart}}
    \subseteq
    \Fun(\Delta^1,\operatorname{Sch}^{\mathrm{sp}})
\]
be the wide subcategory whose objects are relative spectral contexts
\(X\to S\) and whose morphisms are Cartesian squares. Write
\[
    s:\operatorname{Rel}^{\mathrm{sp}}_{\mathrm{Cart}}
    \longrightarrow
    \operatorname{Sch}^{\mathrm{sp}},
    \qquad
    s(X\to S)=X
\]
for source evaluation. Pulling the Cartesian fibration
\(\int\QCoh\to\operatorname{Sch}^{\mathrm{sp}}\) of
Chapter~\Ref{chap:spectral-algebraic-geometry} back along \(s\) gives the
coefficient fibration over relative contexts. The theorem below says that the
cotangent classifier is a Cartesian section of precisely this fibration.

\begin{theorem}[Cartesian cotangent section of the differentiated doctrine]
\label{thm:cotangent-cartesian-section}
The assignment
\[
    (X\to S)\longmapsto L_{X/S}\in\QCoh(X)
\]
defines a canonical Cartesian section of the pulled-back coefficient
fibration over \(\operatorname{Rel}^{\mathrm{sp}}_{\mathrm{Cart}}\).
Equivalently, for every Cartesian square
\[
\begin{tikzcd}[column sep=large,row sep=large]
    X' \arrow[r,"g'"] \arrow[d]
        & X \arrow[d]
    \\
    S' \arrow[r,"g"']
        & S,
\end{tikzcd}
\]
the Cartesian transport of \(L_{X/S}\) is canonically
\[
    (g')^*L_{X/S}\simeq L_{X'/S'}.
\]
The identity, composition, and higher-pasting coherences are those of the
Cartesian transport and of global cotangent substitution.
\end{theorem}

\begin{proof}
\leavevmode
\begin{enumerate}[label=\textup{(\arabic*)}]
    \item \textbf{Restrict the construction.}
    The global cotangent object is defined by its compatible affine restrictions.

    \item \textbf{Identify the resulting object.}
    For every Cartesian square the equivalence of \Ref{thm:global-cotangent-base-change}
    identifies the pullback of those affine classifiers with the classifiers of the pulled-back
    relative context.

    \item \textbf{These equivalences.}
    These equivalences are coherent under identities and iterated Cartesian squares by that
    theorem.

    \item \textbf{They are therefore exactly the data asserting.}
    They are therefore exactly the data asserting that the section is Cartesian in the
    pulled-back coefficient fibration.
\end{enumerate}
\end{proof}

\begin{remark}[Differentiated substitution and chain rule]
\label{rem:differentiated-substitution-chain-rule}
The Cartesian-section theorem is the substitution law of the differentiated
spectral doctrine. The global cofibre sequence
\[
    f^*L_{Y/S}\longrightarrow L_{X/S}\longrightarrow L_{X/Y}
\]
of \Ref{thm:global-cotangent-transitivity} is its chain rule. Both are
constructed consequences of the cotangent representing property and spectral
Zariski descent; no independent differential coherence is imposed.
\end{remark}
The resulting interface is now fixed: affinely the stable coefficient fibre is
the tangent linearization, while globally \(L_{X/S}\) classifies variation,
is transported Cartesianly by substitution, and satisfies the chain rule. The
first branch can therefore remain entirely inside the stable coefficient
fibre. The cotangent classifier will re-enter only at the comparison with the
universal spectral derivation and then as the input to the geometric branch.

\section{Differential-linear structure of the stable coefficient fibres}
\label{sec:differential-linear-tangent-fibres}

The first branch remains entirely inside the stable coefficient fibre.
Affinely this fibre is literally the tangent stabilization identified in
\Ref{thm:tangent-category-calg}; for a general spectral scheme we use the
Zariski-descended fibre \(\QCoh(X)\) and make no additional global
stabilization claim. We now construct on every such stable coefficient fibre
the cofree cocommutative-coalgebra exponential, its primitive infinitesimal
structure, and the resulting codereliction and deriving transformation. The
branch concludes by comparing that derivative with the spectral universal
derivation and by recording its actual oplax behavior under substitution.

\subsection{Cofree cocommutative coalgebras}

\begin{definition}[Cocommutative coalgebras in a stable coefficient fibre]
\label{def:tangent-cocommutative-coalgebras}
For a spectral context \(X\), put
\[
    \mathcal L_X:=\QCoh(X)
\]
and define
\[
    \operatorname{CoCom}(\mathcal L_X)
    :=
    \CAlg(\mathcal L_X^{\op})^{\op}.
\]
Write
\[
    U_X:\operatorname{CoCom}(\mathcal L_X)\longrightarrow\mathcal L_X
\]
for the forgetful functor.
\end{definition}

\begin{theorem}[Cofree coalgebra and the linear exponential]
\label{thm:cofree-coalgebra-linear-exponential}
For every spectral context \(X\), the forgetful functor \(U_X\) admits a
cofree right adjoint
\[
    C_X:\mathcal L_X\longrightarrow\operatorname{CoCom}(\mathcal L_X).
\]
Define
\[
    !_X:=U_XC_X:\mathcal L_X\longrightarrow\mathcal L_X.
\]
Then \(!_X\) carries the canonical comonad structure induced by
\(U_X\dashv C_X\). For every \(M\in\mathcal L_X\), the object \(!_XM\)
carries canonical cocommutative comonoid maps
\[
    \Delta_M:!_XM\longrightarrow !_XM\otimes !_XM,
    \qquad
    e_M:!_XM\longrightarrow\mathcal O_X.
\]
\end{theorem}

\begin{proof}
\leavevmode
\begin{enumerate}[label=\textup{(\arabic*)}]
    \item \textbf{Use preservation.}
    The category \(\mathcal L_X=\QCoh(X)\) is presentable, stable, and symmetric monoidal, and its tensor
    product is accessible and preserves colimits separately.

    \item \textbf{existence theorem for cofree spectral cocommutative coalgebras.}
    The existence theorem for cofree spectral cocommutative coalgebras therefore applies; see
    \cite{BachmannBurklund2026Coalgebras}.

    \item \textbf{Use the product structure.}
    The same source identifies \(\operatorname{CoCom}(\mathcal L_X)\) as Cartesian symmetric monoidal and the forgetful
    functor \(U_X\) as symmetric monoidal, so the underlying object of a product of
    coalgebras is the tensor product of the underlying coefficients.

    \item \textbf{Use the adjunction.}
    The adjunction \(U_X\dashv C_X\) produces the comonad \(U_XC_X\).

    \item \textbf{Use the stated property.}
    Since \(C_X(M)\) is a cocommutative coalgebra, its underlying object carries the stated
    diagonal and counit.
\end{enumerate}
\end{proof}

Throughout this section we repeatedly compare morphisms into a cofree
coalgebra. We use the following consequence of the adjunction without
reintroducing it each time: two such morphisms are canonically compared once
their cofree adjuncts are canonically compared, and the corresponding higher
coherence is transported through the same mapping-space equivalence. Thus
later proofs will expose the adjunct that controls a comparison rather than
repeat the general uniqueness argument.

\begin{proposition}[Canonical monoidal structure of the cofree exponential]
\label{prop:cofree-monoidal-coalgebra-modality}
For \(M,N\in\mathcal L_X\), there is a canonical coalgebra morphism
\[
    \widetilde m_{M,N}:
    C_X(M)\times C_X(N)
    \longrightarrow
    C_X(M\otimes N)
\]
whose cofree adjunct is
\[
    !_XM\otimes !_XN
    \xrightarrow{\ \varepsilon_M\otimes\varepsilon_N\ }
    M\otimes N.
\]
Write
\[
    m_{M,N}:=U_X\widetilde m_{M,N}:
    !_XM\otimes !_XN\longrightarrow !_X(M\otimes N).
\]
The identity of the tensor unit similarly has a canonical cofree lift whose
underlying map is
\[
    m_{\mathcal O_X}:
    \mathcal O_X\longrightarrow !_X\mathcal O_X.
\]
These maps give \(C_X\) its canonical lax symmetric-monoidal structure and
make \(!_X=U_XC_X\) a lax symmetric-monoidal comonad. The comonad counit and
comultiplication are monoidal, and every \(m_{M,N}\) and
\(m_{\mathcal O_X}\) is a morphism of cocommutative coalgebras.
Consequently, after passage to the homotopy category, the maps
\[
    (!_X,\delta,\varepsilon,\Delta,e,m,m_{\mathcal O})
\]
form the ordinary monoidal coalgebra-modality structure associated to the
cofree cocommutative coalgebra.
\end{proposition}

\begin{proof}
\leavevmode
\begin{enumerate}[label=\textup{(\arabic*)}]
    \item \textbf{product coalgebra.}
    The product coalgebra
    \[
        C_X(M)\times C_X(N)
    \]
    has underlying coefficient
    \[
        !_XM\otimes !_XN.
    \]

    \item \textbf{Use the cofree universal property.}
    The cofree mapping-space equivalence therefore gives a canonical coalgebra morphism
    \[
        \widetilde m_{M,N}
    \]
    corresponding to the specified map
    \[
        \varepsilon_M\otimes\varepsilon_N:
        !_XM\otimes !_XN
        \longrightarrow
        M\otimes N.
    \]

    \item \textbf{Establish the next comparison.}
    Equivalently,
    \[
        \varepsilon_{M\otimes N}\,m_{M,N}
        \simeq
        \varepsilon_M\otimes\varepsilon_N.
    \]

    \item \textbf{Use the adjunction.}
    For the tensor unit, the terminal cocommutative coalgebra has underlying object \(\mathcal O_X\);
    applying the same adjunction to \(\id_{\mathcal O_X}\) gives \(m_{\mathcal O_X}\).

    \item \textbf{Associativity, unitality, and symmetry require no additional.}
    Associativity, unitality, and symmetry require no additional choices.

    \item \textbf{Use the cofree universal property.}
    For example, the two coalgebra maps
    \[
        C_X(M)\times C_X(N)\times C_X(P)
        \rightrightarrows
        C_X(M\otimes N\otimes P)
    \]
    obtained from the two bracketings are carried by the cofree mapping-space equivalence to the
    same specified adjunct
    \[
        \varepsilon_M\otimes\varepsilon_N\otimes\varepsilon_P.
    \]

    \item \textbf{Use the cofree universal property.}
    By the cofree comparison convention above, the common adjunct determines the canonical
    comparison between the two coalgebra maps and its higher coherence.

    \item \textbf{unit and symmetry comparisons.}
    The unit and symmetry comparisons are obtained identically.

    \item \textbf{Use the adjunction.}
    This is precisely the canonical lax symmetric-monoidal structure on the right adjoint of the
    symmetric monoidal functor \(U_X\).

    \item \textbf{counit is monoidal by the defining adjunct.}
    The counit is monoidal by the defining adjunct equation above.

    \item \textbf{For the comultiplication, recall that where.}
    For the comultiplication, recall that
    \[
        \delta_M
        =
        U_X\alpha_{C_X(M)},
    \]
    where \(\alpha\) is the unit of \(U_X\dashv C_X\).

    \item \textbf{Compare the two coalgebra.}
    Compare the two coalgebra maps from
    \[
        C_X(M)\times C_X(N)
    \]
    to
    \[
        C_X(!_X(M\otimes N))
    \]
    which express monoidality of \(\delta\).

    \item \textbf{Use the cofree universal property.}
    After composition with the cofree counit, both correspond to the same specified adjunct
    \(m_{M,N}\), using the triangle identity
    \[
        \varepsilon_{U_XD}U_X\alpha_D
        \simeq
        \id_{U_XD}.
    \]

    \item \textbf{Transport the same comparison.}
    The same cofree comparison principle gives the canonical monoidality comparison for
    \(\delta\), together with its inherited higher coherence.

    \item \textbf{Conclude the argument.}
    Finally,
    \[
        \widetilde m_{M,N}
    \]
    and the unit lift are morphisms in
    \[
        \operatorname{CoCom}(\mathcal L_X)
    \]
    by construction.

    \item \textbf{Use preservation.}
    Their underlying maps therefore preserve \(\Delta\) and \(e\).

    \item \textbf{These are exactly the remaining monoidal coalgebra-modality.}
    These are exactly the remaining monoidal coalgebra-modality compatibilities after passage to
    the homotopy category.
\end{enumerate}
\end{proof}

\begin{remark}[No naive symmetric-power formula]
\label{rem:no-naive-cofree-symmetric-power-formula}
The notation \(!_XM\) denotes the underlying object of the cofree
cocommutative coalgebra. Its definition throughout this chapter is the
universal property of \(C_X\). In spectra the cosifted-limit and Tate
phenomena governing cofree coalgebras obstruct a naive infinite-product
formula in general; see \cite{BachmannBurklund2026Coalgebras}.
\end{remark}


\begin{proposition}[Seely equivalences in the stable coefficient fibre]
\label{prop:tangent-seely-equivalences}
For \(M,N\in\mathcal L_X\), there are canonical equivalences
\[
    \chi_{M,N}:!_X(M\oplus N)
    \xrightarrow{\simeq}
    !_XM\otimes !_XN,
    \qquad
    \chi_0:!_X0\xrightarrow{\simeq}\mathcal O_X.
\]
They are coherent under identities, associativity, symmetry, and finite
iteration. Under the product description of
\(\operatorname{CoCom}(\mathcal L_X)\), the two components of the inverse
comparison
\[
    \chi^{-1}_{M,N}:!_XM\otimes !_XN\longrightarrow !_X(M\oplus N)
\]
after composition with the cofree counit are
\[
    \varepsilon_M\otimes e_N:
    !_XM\otimes !_XN\longrightarrow M,
    \qquad
    e_M\otimes\varepsilon_N:
    !_XM\otimes !_XN\longrightarrow N.
\]
After passage to the homotopy category these are the Seely isomorphisms of the
storage modality associated to the cofree cocommutative coalgebra.
\end{proposition}

\begin{proof}
\leavevmode
\begin{enumerate}[label=\textup{(\arabic*)}]
    \item \textbf{Use the product structure.}
    The category of cocommutative coalgebras is Cartesian on the theorem domain of
    \Ref{thm:cofree-coalgebra-linear-exponential}, and its Cartesian product is carried by
    \(U_X\) to the tensor product of underlying objects; see
    \cite{BachmannBurklund2026Coalgebras}.

    \item \textbf{Use the adjunction.}
    Since \(C_X\) is a right adjoint, it preserves products.

    \item \textbf{Use the product structure.}
    Finite products in the stable category \(\mathcal L_X\) are biproducts, so there is a canonical
    equivalence
    \[
        C_X(M\oplus N)
        \simeq
        C_X(M)\times C_X(N).
    \]

    \item \textbf{Use terminality.}
    Applying \(U_X\) gives \(\chi_{M,N}\), and preservation of the terminal object gives
    \(\chi_0\).

    \item \textbf{Identify the resulting object.}
    To identify the comparison intrinsically, recall that a product projection
    \[
        C_X(M)\times C_X(N)\longrightarrow C_X(M)
    \]
    has underlying map
    \[
        !_XM\otimes !_XN
        \xrightarrow{1\otimes e_N}
        !_XM\otimes\mathcal O_X
        \simeq
        !_XM,
    \]
    and similarly for the second projection.

    \item \textbf{Use the product structure.}
    Hence the product comparison \(!_{X}(M\oplus N)\to !_XM\otimes !_XN\) is exactly the canonical Seely comparison obtained
    from the coalgebra diagonal followed by the two biproduct projections.

    \item \textbf{Use the cofree universal property.}
    Conversely, the inverse product equivalence is the unique coalgebra map into \(C_X(M\oplus N)\)
    whose cofree adjunct has components
    \[
        \varepsilon_M\otimes e_N
        \qquad\text{and}\qquad
        e_M\otimes\varepsilon_N.
    \]

    \item \textbf{Use the adjunction.}
    The identity, associativity, symmetry, unit, and all higher-pasting coherences are inherited
    from product preservation by the right adjoint \(C_X\) and the Cartesian coherence of
    \(\operatorname{CoCom}(\mathcal L_X)\).
\end{enumerate}
\end{proof}


\begin{proposition}[Canonical identification of the monoidal storage maps]
\label{prop:cofree-storage-identification}
The lax symmetric-monoidal maps
\[
    m_{M,N}:!_XM\otimes !_XN\longrightarrow !_X(M\otimes N),
    \qquad
    m_{\mathcal O_X}:\mathcal O_X\longrightarrow !_X\mathcal O_X
\]
constructed in \Ref{prop:cofree-monoidal-coalgebra-modality} are canonically
the monoidal maps induced from the Seely equivalences of
\Ref{prop:tangent-seely-equivalences}. Explicitly, if
\(m^{\mathrm{Seely}}_{M,N}\) denotes the ordinary storage-modality composite
recalled in \cite[Section~2]{Lemay2021Coderelictions}
\[
\begin{aligned}
    !_XM\otimes !_XN
    &\xrightarrow{\chi^{-1}_{M,N}}
    !_X(M\oplus N)
    \xrightarrow{\delta_{M\oplus N}}
    !_X!_X(M\oplus N)
    \\
    &\xrightarrow{!_X(\chi_{M,N})}
    !_X(!_XM\otimes !_XN)
    \xrightarrow{!_X(\varepsilon_M\otimes\varepsilon_N)}
    !_X(M\otimes N),
\end{aligned}
\]
then there is a canonical equivalence
\[
    m^{\mathrm{Seely}}_{M,N}\simeq m_{M,N}.
\]
The corresponding tensor-unit maps are canonically equivalent as well.
\end{proposition}

\begin{proof}
\leavevmode
\begin{enumerate}[label=\textup{(\arabic*)}]
    \item \textbf{Compare the two constructions.}
    Both \(m^{\mathrm{Seely}}_{M,N}\) and \(m_{M,N}\) underlie coalgebra maps
    \[
        C_X(M)\times C_X(N)\longrightarrow C_X(M\otimes N).
    \]

    \item \textbf{Use the cofree universal property.}
    By the cofree mapping-space equivalence it is enough to compare their composites with the
    cofree counit.

    \item \textbf{Verify naturality.}
    Naturality of \(\varepsilon\), followed by the comonad counit identity, gives
    \[
    \begin{aligned}
        \varepsilon_{M\otimes N}m^{\mathrm{Seely}}_{M,N}
        &\simeq
        (\varepsilon_M\otimes\varepsilon_N)
        \varepsilon_{!_XM\otimes !_XN}
        !_X(\chi_{M,N})
        \delta_{M\oplus N}
        \chi^{-1}_{M,N}
        \\
        &\simeq
        (\varepsilon_M\otimes\varepsilon_N)
        \chi_{M,N}
        \varepsilon_{!_X(M\oplus N)}
        \delta_{M\oplus N}
        \chi^{-1}_{M,N}
        \\
        &\simeq
        \varepsilon_M\otimes\varepsilon_N.
    \end{aligned}
    \]

    \item \textbf{Use the cofree universal property.}
    This is exactly the specified cofree adjunct of \(m_{M,N}\).

    \item \textbf{Derive the comparison.}
    The cofree comparison principle therefore gives the claimed canonical comparison, with its
    higher coherence.

    \item \textbf{Use terminality.}
    The tensor-unit statement is the same argument applied to the terminal coalgebra.
\end{enumerate}
\end{proof}

\begin{remark}[The storage structure is constructed]
\label{rem:tangent-storage-constructed}
The exponential modality used below is the canonical storage package induced
by the adjunction \(U_X\dashv C_X\) and product preservation of the cofree
right adjoint. The comonad maps come from the adjunction, the Seely
equivalences come from product preservation, and
\Ref{prop:cofree-storage-identification} proves that the directly constructed
monoidal maps \(m_{M,N}\) and \(m_{\mathcal O_X}\) are exactly the
canonical monoidal storage maps induced from those Seely equivalences. The
corresponding identification of the additive bialgebra maps is proved below
after those maps have been constructed.
\end{remark}
Differentiation now requires a canonical primitive infinitesimal element that
can be transported into the cofree coalgebra.

\subsection{Primitive infinitesimal coalgebras and coherent infinitesimal augmentation}

\begin{lemma}[Opposite split square-zero construction]
\label{lem:opposite-split-square-zero}
Let \(\mathcal C^\otimes\) be a stable symmetric-monoidal
\(\infty\)-category whose tensor product is exact separately in each variable.
Then \((\mathcal C^{\op})^\otimes\) is again stable symmetric monoidal with
tensor product exact separately, and every \(M\in\mathcal C\), regarded as an
object of \(\mathcal C^{\op}\), determines functorially a split square-zero
commutative algebra
\[
    \mathbb 1_{\mathcal C}\oplus M
    \in
    \CAlg(\mathcal C^{\op})
\]
whose augmentation ideal is \(M\). Its underlying object is the biproduct
\(\mathbb 1_{\mathcal C}\oplus M\), the unit acts in the evident way, and
the product on the \(M\)-summand is coherently null.
\end{lemma}

\begin{proof}
\leavevmode
\begin{enumerate}[label=\textup{(\arabic*)}]
    \item \textbf{Use exactness.}
    Opposition preserves stability, and an exact functor between stable \(\infty\)-categories
    remains exact after passage to opposites.

    \item \textbf{Use exactness.}
    Hence the tensor product on \(\mathcal C^{\op}\) is exact separately.

    \item \textbf{Identify the resulting object.}
    In a stable symmetric-monoidal \(\infty\)-category, augmentation ideal and unitalization
    identify augmented commutative algebra objects with nonunital commutative algebra objects;
    see \cite[Proposition~5.4.4.10]{LurieHigherAlgebra}.

    \item \textbf{Equip functorially with the trivial nonunital commutative algebra.}
    Equip \(M\) functorially with the trivial nonunital commutative algebra structure:
    every multiplication involving at least two inputs from \(M\) is the zero morphism,
    with the coherence supplied by the zero object of the stable category.

    \item \textbf{Use the product structure.}
    Its unitalization has underlying biproduct
    \[
        \mathbb 1_{\mathcal C}\oplus M,
    \]
    the tensor unit acts by the canonical unit constraints, and the product on \(M\otimes M\) is
    coherently null.

    \item \textbf{Verify compatibility.}
    Applying this construction in \(\mathcal C^{\op}\) gives the stated split square-zero commutative
    algebra.

    \item \textbf{Verify functoriality.}
    Functoriality is inherited from trivial nonunital algebra formation and unitalization.

    \item \textbf{Identify the resulting comparison.}
    This finite algebraic construction does not require presentability of \(\mathcal C^{\op}\).
\end{enumerate}
\end{proof}

\begin{construction}[Primitive infinitesimal coalgebra]
\label{con:primitive-infinitesimal-coalgebra}
Let
\[
    \mathcal L_X:=\QCoh(X).
\]
The tensor product on \(\mathcal L_X\) is exact separately, so
\Ref{lem:opposite-split-square-zero} applies. For
\(M\in\mathcal L_X\), form in \(\mathcal L_X^{\op}\) the functorial split
square-zero commutative algebra with underlying object
\[
    \mathcal O_X\oplus M
\]
and square-zero augmentation ideal \(M\). Passing to the opposite category
defines a cocommutative \(\mathbb E_\infty\)-coalgebra
\[
    P_X(M)\in\operatorname{CoCom}(\mathcal L_X),
    \qquad
    U_XP_X(M)\simeq\mathcal O_X\oplus M.
\]
Write
\[
    \operatorname{pr}_0:U_XP_X(M)\longrightarrow\mathcal O_X,
    \qquad
    \operatorname{pr}_1:U_XP_X(M)\longrightarrow M
\]
for the counit and coefficient projection on underlying objects, and write
\(\iota_0,\iota_1\) for the two biproduct inclusions. The unit summand is
group-like and the coefficient summand is primitive. In particular, the
binary comultiplication has the familiar underlying form
\[
    \Delta(1)=1\otimes1,
    \qquad
    \Delta(m)=m\otimes1+1\otimes m.
\]
The full \(\mathbb E_\infty\)-coalgebra coherence is the one supplied by the
opposite split square-zero algebra construction; the displayed binary formula
is only its underlying first operation.
\end{construction}

\begin{proposition}[Square-zero and primitive infinitesimal structures]
\label{prop:square-zero-primitive-comparison}
The underlying stable object \(\mathcal O_X\oplus M\) supports two canonical
and complementary infinitesimal structures:
\[
\begin{array}{c|c}
\text{split square-zero algebra} & \text{primitive }\mathbb E_\infty\text{-coalgebra}
\\ \hline
M\cdot M\simeq0 &
\Delta(m)=m\otimes1+1\otimes m.
\end{array}
\]
Both structures are functorial in \(M\). The first is the zeroth-space image
of tangent linearization, while the second is obtained by applying the same
split square-zero construction in the opposite stable symmetric-monoidal
\(\infty\)-category.
\end{proposition}

\begin{proof}
\leavevmode
\begin{enumerate}[label=\textup{(\arabic*)}]
    \item \textbf{algebra structure.}
    The algebra structure is \Ref{con:split-square-zero-algebra}.

    \item \textbf{coalgebra structure.}
    The coalgebra structure is \Ref{con:primitive-infinitesimal-coalgebra}.

    \item \textbf{Compare the two constructions.}
    Both constructions are functorial split square-zero constructions, respectively in
    \(\mathcal L_X\) and \(\mathcal L_X^{\op}\), so no separate collection of arity-wise coherences is chosen.
\end{enumerate}
\end{proof}

\begin{construction}[Canonical coherent infinitesimal augmentation]
\label{con:tangent-infinitesimal-augmentation}
Let
\[
    \alpha_D:
    D\longrightarrow C_XU_XD
\]
denote the unit of the adjunction \(U_X\dashv C_X\) at a cocommutative
coalgebra \(D\). For \(M\in\mathcal L_X\), let
\[
    \widehat H^{\inf}_M
    :=
    \alpha_{P_X(M)}:
    P_X(M)\longrightarrow C_XU_XP_X(M)
\]
and write
\[
    H^{\inf}_M
    :=
    U_X\widehat H^{\inf}_M:
    \mathcal O_X\oplus M
    \longrightarrow
    !_X\bigl(\mathcal O_X\oplus M\bigr).
\]
Thus \(\widehat H^{\inf}_M\) is a morphism of cocommutative
\(\mathbb E_\infty\)-coalgebras, while its underlying map \(H^{\inf}_M\)
is the canonical coalgebra structure for the comonad \(!_X=U_XC_X\):
\[
    \varepsilon_{\mathcal O_X\oplus M}H^{\inf}_M
    \simeq
    \id_{\mathcal O_X\oplus M},
\]
and
\[
    \delta_{\mathcal O_X\oplus M}H^{\inf}_M
    \simeq
    !_X(H^{\inf}_M)H^{\inf}_M.
\]
All these equivalences and their higher compatibility are the unit and
triangle coherences of the adjunction.
\end{construction}

\begin{construction}[Primitive tensor comparison]
\label{con:primitive-tensor-comparison}
Let \(D\in\operatorname{CoCom}(\mathcal L_X)\) and \(N\in\mathcal L_X\).
There is a canonical coalgebra morphism
\[
    \vartheta_{D,N}:
    D\times P_X(N)
    \longrightarrow
    P_X(U_XD\otimes N).
\]
To construct it, pass to \(\mathcal L_X^{\op}\). The product on the left
becomes the tensor product of the commutative algebra corresponding to \(D\)
with the split square-zero algebra \(\mathcal O_X\oplus N\). Its summand
\(U_XD\otimes N\) is a square-zero ideal because the \(N\)-summand is
square-zero. Hence the split square-zero universal construction gives a
canonical commutative-algebra morphism
\[
    \mathcal O_X\oplus(U_XD\otimes N)
    \longrightarrow
    A_D\otimes(\mathcal O_X\oplus N),
\]
where \(A_D\in\CAlg(\mathcal L_X^{\op})\) is the commutative algebra
corresponding to \(D\). Passing to opposites gives \(\vartheta_{D,N}\).
Under the symmetric-monoidal equivalence
\[
    U_X(D\times P_X(N))
    \simeq
    U_XD\otimes(\mathcal O_X\oplus N),
\]
its underlying map has the two components
\[
    e_D\otimes\operatorname{pr}_0
    :U_XD\otimes(\mathcal O_X\oplus N)\longrightarrow\mathcal O_X
\]
and
\[
    1\otimes\operatorname{pr}_1
    :U_XD\otimes(\mathcal O_X\oplus N)\longrightarrow U_XD\otimes N.
\]
Naturality of the adjunction unit makes \(U_X\vartheta_{D,N}\) a morphism between
the canonical \(!_X\)-coalgebra structures induced from \(D\times P_X(N)\)
and \(P_X(U_XD\otimes N)\).
\end{construction}

\begin{proposition}[The adjunction unit is an infinitesimal augmentation]
\label{prop:tangent-infinitesimal-augmentation-identities}
The maps \(H^{\inf}_M\) of
\Ref{con:tangent-infinitesimal-augmentation} satisfy the three
infinitesimal-augmentation identities coherently. Explicitly:
\begin{enumerate}[label=(\roman*)]
    \item \(H^{\inf}_M\) is a \(!_X\)-coalgebra structure:
    \[
        \varepsilon H^{\inf}_M\simeq\id,
        \qquad
        \delta H^{\inf}_M
        \simeq
        !_X(H^{\inf}_M)H^{\inf}_M;
    \]

    \item it is compatible with the primitive cocommutative coalgebra structure:
    if \(\Lambda_M\) denotes the comultiplication of \(P_X(M)\), then
    \[
        \Delta_{\mathcal O_X\oplus M}H^{\inf}_M
        \simeq
        (H^{\inf}_M\otimes H^{\inf}_M)\Lambda_M,
        \qquad
        e_{\mathcal O_X\oplus M}H^{\inf}_M\simeq\operatorname{pr}_0;
    \]

    \item for \(M,N\in\mathcal L_X\), the primitive tensor comparison
    \[
        U_X\vartheta_{C_X(M),N}:
        !_XM\otimes(\mathcal O_X\oplus N)
        \longrightarrow
        \mathcal O_X\oplus(!_XM\otimes N)
    \]
    is a morphism of \(!_X\)-coalgebras, where the source coalgebra structure is
    the canonical tensor product structure
    \[
    \begin{aligned}
        !_XM\otimes(\mathcal O_X\oplus N)
        &\xrightarrow{\ \delta_M\otimes H^{\inf}_N\ }
        !_X(!_XM)\otimes !_X(\mathcal O_X\oplus N)
        \\
        &\xrightarrow{\ m_{!_XM,\,\mathcal O_X\oplus N}\ }
        !_X\bigl(!_XM\otimes(\mathcal O_X\oplus N)\bigr),
    \end{aligned}
    \]
    and the target coalgebra structure is \(H^{\inf}_{!_XM\otimes N}\).
\end{enumerate}
After passage to the homotopy category these are precisely the conditions
\textup{IA.1}--\textup{IA.3} of the infinitesimal augmentation of
\cite[Definition~10]{Lemay2021Coderelictions}.
\end{proposition}

\begin{proof}
\leavevmode
\begin{enumerate}[label=\textup{(\arabic*)}]
    \item \textbf{Use the adjunction.}
    Part \textup{(i)} is the comonad-coalgebra equation for the underlying adjunction-unit map
    \[
        U_X\alpha_{P_X(M)}=H^{\inf}_M.
    \]

    \item \textbf{Derive the comparison.}
    Part \textup{(ii)} follows because \(\alpha_{P_X(M)}:P_X(M)\to C_XU_XP_X(M)\) is a morphism in \(\operatorname{CoCom}(\mathcal L_X)\); applying
    \(U_X\) gives the two displayed coalgebra compatibilities.

    \item \textbf{For (iii), first note that for arbitrary coalgebras there.}
    For \textup{(iii)}, first note that for arbitrary coalgebras \(D,E\) there is a
    canonical equivalence
    \[
        U_X\alpha_{D\times E}
        \simeq
        m_{U_XD,U_XE}
        \bigl(U_X\alpha_D\otimes U_X\alpha_E\bigr).
    \]

    \item \textbf{compare the two coalgebra.}
    Indeed, compare the two coalgebra maps
    \[
        D\times E\rightrightarrows C_X(U_XD\otimes U_XE).
    \]

    \item \textbf{Use the cofree universal property.}
    After composition with the cofree counit, both have adjunct \(\id_{U_XD\otimes U_XE}\), by the triangle
    identities for \(U_X\dashv C_X\).

    \item \textbf{cofree comparison principle therefore.}
    The cofree comparison principle therefore supplies the displayed comparison and its higher
    coherence.

    \item \textbf{Apply this with and .}
    Apply this with \(D=C_X(M)\) and \(E=P_X(N)\).

    \item \textbf{Then so the canonical -coalgebra structure on the underlying product.}
    Then
    \[
        U_X\alpha_{C_X(M)}=\delta_M,
        \qquad
        U_X\alpha_{P_X(N)}=H^{\inf}_N,
    \]
    so the canonical \(!_X\)-coalgebra structure on the underlying product \(!_XM\otimes(\mathcal O_X\oplus N)\) is
    exactly the source structure displayed in \textup{(iii)}.

    \item \textbf{Verify naturality.}
    Naturality of the adjunction unit for the coalgebra morphism
    \[
        \vartheta_{C_X(M),N}:
        C_X(M)\times P_X(N)
        \longrightarrow
        P_X(!_XM\otimes N)
    \]
    gives
    \[
    \begin{aligned}
        !_X(U_X\vartheta_{C_X(M),N})\,
        U_X\alpha_{C_X(M)\times P_X(N)}
        \simeq
        H^{\inf}_{!_XM\otimes N}\,
        U_X\vartheta_{C_X(M),N}.
    \end{aligned}
    \]

    \item \textbf{Identify the resulting comparison.}
    This is exactly the required \(!_X\)-coalgebra-morphism comparison.

    \item \textbf{Verify naturality.}
    All higher compatibility comes from naturality of \(\alpha\), the monoidal coherences of
    \Ref{prop:cofree-monoidal-coalgebra-modality}, and functoriality of products in \(\operatorname{CoCom}(\mathcal L_X)\).
\end{enumerate}
\end{proof}

\begin{construction}[Codereliction]
\label{con:tangent-codereliction}
Let
\[
    \operatorname{pr}_1:U_XP_X(M)=\mathcal O_X\oplus M\longrightarrow M
\]
be the coefficient projection. Its cofree adjunct is the canonical coalgebra
morphism
\[
\begin{aligned}
    H_M:
    P_X(M)
    &\xrightarrow{\ \alpha_{P_X(M)}\ }
    C_X(\mathcal O_X\oplus M)
    \\
    &\xrightarrow{\ C_X(\operatorname{pr}_1)\ }
    C_X(M).
\end{aligned}
\]
Equivalently, \(H_M\) is characterized by
\[
    \varepsilon_M\,U_XH_M\simeq\operatorname{pr}_1,
\]
and functoriality of the cofree adjunction identifies its underlying map as
\[
    U_XH_M
    \simeq
    !_X(\operatorname{pr}_1)H^{\inf}_M.
\]
Restricting along the coefficient inclusion therefore gives
\[
    \eta_M:
    M
    \xrightarrow{\iota_1}
    \mathcal O_X\oplus M
    \xrightarrow{U_XH_M}
    !_XM.
\]
We call \(\eta_M\) the \textbf{codereliction map}. Its recognition as a
codereliction in the ordinary differential-storage sense is made after
passage to the homotopy category in
\Ref{thm:tangent-differential-category-recognition}.
\end{construction}

\begin{proposition}[Basic coherent codereliction identities]
\label{prop:basic-codereliction-identities}
Let
\[
    \varepsilon_M:!_XM\longrightarrow M
\]
be the comonad counit, let
\(e_M:!_XM\to\mathcal O_X\) be the coalgebra counit, and let
\[
    u_M:\mathcal O_X\longrightarrow !_XM
\]
be the group-like morphism obtained by composing
\(\iota_0:\mathcal O_X\to P_X(M)\) with \(U_XH_M\). Then there are canonical
coherent equivalences
\[
    \varepsilon_M\eta_M\simeq\id_M,
    \qquad
    e_M\eta_M\simeq0,
\]
and
\[
    \Delta_M\eta_M
    \simeq
    \eta_M\otimes u_M+u_M\otimes\eta_M.
\]
These are the linear, constant, and primitive Leibniz identities for the
codereliction.
\end{proposition}

\begin{proof}
\leavevmode
\begin{enumerate}[label=\textup{(\arabic*)}]
    \item \textbf{Restrict the construction.}
    The first identity follows from the defining adjunct equation \(\varepsilon_MU_XH_M\simeq\operatorname{pr}_1\) after restriction
    along \(\iota_1\).

    \item \textbf{Identify the map.}
    The map \(H_M\) is a morphism of cocommutative \(\mathbb E_\infty\)-coalgebras.

    \item \textbf{Use preservation.}
    It therefore preserves the coalgebra counit and the full coalgebra comultiplication.

    \item \textbf{On the coefficient summand of , the counit.}
    On the coefficient summand of \(P_X(M)\), the counit is zero and the binary
    comultiplication is primitive.

    \item \textbf{Restrict the construction.}
    Restriction along \(\iota_1\) gives the second and third displayed equivalences.

    \item \textbf{Their higher compatibility.}
    Their higher compatibility is inherited from the single coalgebra morphism \(H_M\),
    rather than imposed on the three formulas separately.
\end{enumerate}
\end{proof}
The primitive infinitesimal input can now be combined with the bialgebra
multiplication to produce infinitesimal translation and hence the deriving
transformation.

\subsection{Coalgebra-level infinitesimal translation and deriving transformation}

\begin{construction}[Bialgebra multiplication of the cofree exponential]
\label{con:tangent-bialgebra-multiplication}
The product coalgebra
\[
    C_X(M)\times C_X(M)
\]
has underlying object \(!_XM\otimes !_XM\). Define
\[
    b_M:
    !_XM\otimes !_XM
    \longrightarrow
    M
\]
by
\[
    b_M
    :=
    \varepsilon_M\otimes e_M+e_M\otimes\varepsilon_M.
\]
By the cofree adjunction there is a canonical coalgebra morphism
\[
    \widetilde\nabla_M:
    C_X(M)\times C_X(M)
    \longrightarrow
    C_X(M)
\]
whose adjunct is \(b_M\). Write
\[
    \nabla_M:=U_X\widetilde\nabla_M:
    !_XM\otimes !_XM\longrightarrow !_XM.
\]
Similarly, \(P_X(0)\), whose underlying object is \(\mathcal O_X\), is the
terminal cocommutative coalgebra. The zero morphism
\(\mathcal O_X\to M\) therefore has a cofree lift
\[
    \widetilde u_M:P_X(0)
    \longrightarrow C_X(M),
\]
whose underlying map is \(u_M:\mathcal O_X\to !_XM\). This agrees with the
map \(u_M\) defined from the group-like summand above, because both coalgebra
maps have zero as their cofree adjunct.
\end{construction}

\begin{proposition}[Canonical identification of the additive storage maps]
\label{prop:cofree-additive-storage-identification}
Under the Seely equivalence
\[
    \chi_{M,M}:!_X(M\oplus M)\xrightarrow{\simeq}!_XM\otimes !_XM,
\]
the coalgebra morphism obtained from
\(C_X(+):C_X(M\oplus M)\to C_X(M)\) has underlying map
\[
    \nabla_M:!_XM\otimes !_XM\longrightarrow !_XM
\]
of \Ref{con:tangent-bialgebra-multiplication}. Likewise the coalgebra map
induced by \(0\to M\) has underlying map
\[
    u_M:\mathcal O_X\longrightarrow !_XM.
\]
Hence \(\nabla_M\) and \(u_M\) are exactly the additive bialgebra maps
canonically induced by the Seely storage structure.
\end{proposition}

\begin{proof}
\leavevmode
\begin{enumerate}[label=\textup{(\arabic*)}]
    \item \textbf{Transport across.}
    Transport \(C_X(+)\) across
    \[
        C_X(M\oplus M)\simeq C_X(M)\times C_X(M).
    \]

    \item \textbf{Use the cofree universal property.}
    The resulting coalgebra map into \(C_X(M)\) has cofree adjunct
    \[
        (+)\,\varepsilon_{M\oplus M}\chi^{-1}_{M,M}.
    \]

    \item \textbf{Invoke the cited result.}
    By the component calculation of \Ref{prop:tangent-seely-equivalences}, its two summands are
    \[
        \varepsilon_M\otimes e_M
        \qquad\text{and}\qquad
        e_M\otimes\varepsilon_M.
    \]

    \item \textbf{Conclude the step.}
    Thus the adjunct is
    \[
        \varepsilon_M\otimes e_M+e_M\otimes\varepsilon_M,
    \]
    which is exactly the map \(b_M\) defining \(\widetilde\nabla_M\) in
    \Ref{con:tangent-bialgebra-multiplication}.

    \item \textbf{Derive the comparison.}
    Cofreeness therefore gives a canonical equivalence between the transported addition map and
    \(\widetilde\nabla_M\), hence between their underlying maps.

    \item \textbf{Use the cofree universal property.}
    The zero-induced unit is treated identically: both coalgebra maps have zero as their cofree
    adjunct.

    \item \textbf{Verify functoriality.}
    Coherence is inherited from the coherent commutative-group structure of biproducts and
    functoriality of \(C_X\).
\end{enumerate}
\end{proof}

\begin{proposition}[Coherent bicommutative bialgebra structure]
\label{prop:tangent-bialgebra-structure}
The maps \(\nabla_M,u_M,\Delta_M,e_M\) make \(!_XM\) a coherently
bicommutative bialgebra object in \(\mathcal L_X\). Under the Seely
equivalence
\[
    !_X(M\oplus M)\simeq !_XM\otimes !_XM,
\]
the multiplication \(\nabla_M\) is the morphism induced by addition
\(M\oplus M\to M\), and \(u_M\) is induced by \(0\to M\).
\end{proposition}

\begin{proof}
\leavevmode
\begin{enumerate}[label=\textup{(\arabic*)}]
    \item \textbf{Use the product structure.}
    The addition and zero maps make every object of the stable \(\infty\)-category \(\mathcal L_X\)
    a coherently commutative group object for the biproduct monoidal structure.

    \item \textbf{Applying and transporting through the product equivalences of.}
    Applying \(C_X\) and transporting through the product equivalences of
    \Ref{prop:tangent-seely-equivalences} produces coalgebra maps
    \[
        C_X(M)\times C_X(M)\longrightarrow C_X(M),
        \qquad
        P_X(0)\longrightarrow C_X(M).
    \]

    \item \textbf{Invoke the cited result.}
    By \Ref{prop:cofree-additive-storage-identification}, their underlying maps are precisely
    \(\nabla_M\) and \(u_M\).

    \item \textbf{Associativity, commutativity, and unitality.}
    Associativity, commutativity, and unitality are therefore the functorial images of the
    coherent addition and zero laws.

    \item \textbf{Use the stated property.}
    Since the maps are morphisms in \(\operatorname{CoCom}(\mathcal L_X)\), multiplication and unit preserve \(\Delta_M\)
    and \(e_M\).

    \item \textbf{Use the product structure.}
    These are the bicommutative bialgebra compatibilities, with all higher coherence inherited
    from the functorial product construction and the stable biproduct structure.
\end{enumerate}
\end{proof}

\begin{construction}[Coalgebra-level infinitesimal translation]
\label{con:tangent-coalgebra-infinitesimal-translation}
Define
\[
    a_M:
    !_XM\otimes(\mathcal O_X\oplus M)
    \longrightarrow
    M
\]
by
\[
    a_M
    :=
    e_M\otimes\operatorname{pr}_1+\varepsilon_M\otimes\operatorname{pr}_0.
\]
The source is the underlying object of the product coalgebra
\(C_X(M)\times P_X(M)\). The cofree adjunction therefore supplies a canonical
coalgebra morphism
\[
    \widetilde{\mathsf d}_M:
    C_X(M)\times P_X(M)
    \longrightarrow
    C_X(M)
\]
characterized by
\[
    \varepsilon_M\,U_X\widetilde{\mathsf d}_M
    \simeq
    a_M.
\]
We call \(\widetilde{\mathsf d}_M\) \textbf{coalgebra-level infinitesimal
translation}.
\end{construction}

\begin{proposition}[Infinitesimal translation from codereliction]
\label{prop:infinitesimal-translation-from-codereliction}
There is a canonical equivalence of coalgebra morphisms
\[
    \widetilde{\mathsf d}_M
    \simeq
    \widetilde\nabla_M\circ(1\times H_M).
\]
Consequently, after restricting the underlying map along
\[
    !_XM\otimes M
    \xrightarrow{1\otimes\iota_1}
    !_XM\otimes(\mathcal O_X\oplus M),
\]
one obtains
\[
    \mathsf d_M:
    !_XM\otimes M
    \xrightarrow{1\otimes\eta_M}
    !_XM\otimes !_XM
    \xrightarrow{\nabla_M}
    !_XM.
\]
\end{proposition}

\begin{proof}
\leavevmode
\begin{enumerate}[label=\textup{(\arabic*)}]
    \item \textbf{Compare the two constructions.}
    Both sides of the first displayed equivalence are coalgebra morphisms
    \[
        C_X(M)\times P_X(M)
        \longrightarrow
        C_X(M).
    \]

    \item \textbf{Use the cofree universal property.}
    By the cofree comparison principle, it is enough to compare their specified adjuncts.

    \item \textbf{Use the stated property.}
    Since \(H_M\) is a coalgebra morphism,
    \[
        e_MU_XH_M\simeq\operatorname{pr}_0,
        \qquad
        \varepsilon_MU_XH_M\simeq\operatorname{pr}_1.
    \]

    \item \textbf{adjunct of.}
    The adjunct of
    \[
        \widetilde\nabla_M(1\times H_M)
    \]
    is consequently
    \[
    \begin{aligned}
        &(\varepsilon_M\otimes e_M+e_M\otimes\varepsilon_M)
        (1\otimes U_XH_M)
        \\
        &\qquad\simeq
        \varepsilon_M\otimes\operatorname{pr}_0
        +
        e_M\otimes\operatorname{pr}_1
        \\
        &\qquad=
        a_M.
    \end{aligned}
    \]

    \item \textbf{Invoke the cited result.}
    By definition, \(a_M\) is also the specified adjunct of \(\widetilde{\mathsf d}_M\).

    \item \textbf{Identify the equivalence.}
    The cofree comparison principle therefore gives the canonical equivalence
    \[
        \widetilde{\mathsf d}_M
        \simeq
        \widetilde\nabla_M(1\times H_M).
    \]

    \item \textbf{Use the adjunction.}
    Its higher compatibility is inherited from the adjunction and the coalgebra maps entering
    the construction.

    \item \textbf{Restrict the construction.}
    Restricting the resulting underlying map along the primitive inclusion
    \[
        !_XM\otimes M
        \xrightarrow{1\otimes\iota_1}
        !_XM\otimes(\mathcal O_X\oplus M)
    \]
    gives
    \[
        !_XM\otimes M
        \xrightarrow{1\otimes\eta_M}
        !_XM\otimes !_XM
        \xrightarrow{\nabla_M}
        !_XM,
    \]
    which is the asserted formula for \(\mathsf d_M\).
\end{enumerate}
\end{proof}

\begin{construction}[Deriving transformation]
\label{con:tangent-deriving-transformation}
The map \(\mathsf d_M\) of
\Ref{prop:infinitesimal-translation-from-codereliction} is the
\textbf{deriving-transformation map}; its ordinary deriving-transformation
recognition is \Ref{thm:tangent-differential-category-recognition}. For a
coKleisli morphism
\[
    f:!_XM\longrightarrow N,
\]
define its differential by
\[
    Df:=f\circ\mathsf d_M:
    !_XM\otimes M\longrightarrow N.
\]
The coalgebra morphism \(\widetilde{\mathsf d}_M\), rather than merely its
restriction \(\mathsf d_M\), is retained as the higher-categorical source of
this operation.
\end{construction}

\begin{theorem}[Coherent differential-exponential structure of the stable coefficient fibre]
\label{thm:tangent-differential-category-recognition}
For every spectral context \(X\), the preceding constructions canonically
realize the differential-exponential structure associated to the cofree
cocommutative coalgebra at the ordinary homotopy-category recognition
interface. More precisely:
\begin{enumerate}[label=(\roman*)]
    \item the homotopy category \(h\mathcal L_X\), with
    \[
        (!_X,\delta,\varepsilon,\Delta,e,m,m_{\mathcal O})
    \]
    and the Seely isomorphisms of
    \Ref{prop:tangent-seely-equivalences}, is an additive storage category;

    \item the image of \(H^{\inf}\) in \(h\mathcal L_X\) is an
    infinitesimal augmentation;

    \item the image of the morphism
    \[
        \eta_M:
        M\xrightarrow{\iota_1}
        \mathcal O_X\oplus M
        \xrightarrow{H^{\inf}_M}
        !_X(\mathcal O_X\oplus M)
        \xrightarrow{!_X(\operatorname{pr}_1)}
        !_XM
    \]
    is therefore a codereliction in \(h\mathcal L_X\);

    \item the deriving transformation associated to this codereliction is the
    image of the already constructed morphism
    \[
        \mathsf d_M
        =
        !_XM\otimes M
        \xrightarrow{1\otimes\eta_M}
        !_XM\otimes !_XM
        \xrightarrow{\nabla_M}
        !_XM.
    \]
\end{enumerate}
Thus \(h\mathcal L_X\) is a differential storage category, hence in
particular a differential category in the sense of
Blute--Cockett--Seely. Before passage to the homotopy category, the primitive
coalgebra, \(H^{\inf}\), the codereliction map, the monoidal and bialgebra
maps, and the coalgebra-level infinitesimal translation are genuine morphisms
in the stated \(\infty\)-categories. The complete ordinary
differential-storage equations are invoked only at the homotopy-category
recognition interface; no separate \(\infty\)-categorical algebraic theory of
differential storage categories is assumed.
\end{theorem}

\begin{proof}
\leavevmode
\begin{enumerate}[label=\textup{(\arabic*)}]
    \item \textbf{Use the cofree universal property.}
    The cofree adjunction gives the monoidal coalgebra modality of
    \Ref{prop:cofree-monoidal-coalgebra-modality}, and \Ref{prop:tangent-seely-equivalences}
    gives canonical Seely isomorphisms.

    \item \textbf{Identify the map.}
    The identification theorems \Ref{prop:cofree-storage-identification} and
    \Ref{prop:cofree-additive-storage-identification} prove that the monoidal and additive
    bialgebra maps already constructed from cofreedom are exactly the storage-modality maps
    induced by those Seely isomorphisms.

    \item \textbf{Use the product structure.}
    Since \(\mathcal L_X\) is stable, \(h\mathcal L_X\) is additive with finite biproducts.

    \item \textbf{Conclude the step.}
    Hence the first assertion is the ordinary additive storage structure determined by these
    displayed spectral maps; no additional storage datum is chosen.

    \item \textbf{coherent identities of the cited result descend.}
    The coherent identities of \Ref{prop:tangent-infinitesimal-augmentation-identities} descend
    exactly to \textup{IA.1}--\textup{IA.3} of \cite[Definition~10]{Lemay2021Coderelictions}.

    \item \textbf{Identify the resulting object.}
    Proposition~13 of that source therefore identifies
    \[
        M\xrightarrow{\iota_1}\mathcal O_X\oplus M
        \xrightarrow{H^{\inf}_M}!_X(\mathcal O_X\oplus M)
        \xrightarrow{!_X(\operatorname{pr}_1)}!_XM
    \]
    as a codereliction.

    \item \textbf{Use the cofree universal property.}
    By \Ref{con:tangent-codereliction}, this is precisely the image of the map \(\eta_M\)
    already constructed as the cofree adjunct of the primitive projection; the recognition
    theorem introduces no second codereliction.

    \item \textbf{Conclude the argument.}
    Finally, \cite[Theorem~8]{Lemay2021Coderelictions} identifies coderelictions with deriving
    transformations in an additive storage category and gives
    \[
        \mathsf d_M=\nabla_M(1\otimes\eta_M).
    \]

    \item \textbf{Restrict the construction.}
    This is the image of \Ref{con:tangent-deriving-transformation}, since
    \Ref{prop:infinitesimal-translation-from-codereliction} constructs the same map as the
    primitive restriction of the coalgebra morphism \(\widetilde{\mathsf d}_M\).

    \item \textbf{Conclude the step.}
    Thus the ordinary constant, Leibniz, linear, chain, interchange, and monoidal equations
    follow from the cited recognition results.

    \item \textbf{Identify the map.}
    The specific higher maps and canonical homotopies constructed before passage to \(h\mathcal L_X\)
    remain the higher-categorical source of that ordinary structure.
\end{enumerate}
\end{proof}

\begin{remark}[Relation to differential linear syntax]
\label{rem:tangent-differential-linear-logic}
The primary differential structure of this chapter lives in the stable
symmetric-monoidal \(\infty\)-category \(\QCoh(X)\): the primitive
\(\mathbb E_\infty\)-coalgebra, infinitesimal augmentation, codereliction, and
coalgebra-level infinitesimal translation are constructed before any passage
to homotopy classes. The homotopy category is only a recognition interface
with the differential-category terminology of
\cite{BluteCockettSeely2006Differential,Lemay2021Coderelictions}. This places
the stable coefficient fibres in the linear/nonlinear tradition of
\cite{Benton1995LNL} and relates their coKleisli differential to the
differential lambda-calculus of \cite{EhrhardRegnier2003Differential}, without
discarding the higher categorical information of the spectral doctrine.
\end{remark}
It remains to compare this cofree derivative with cotangent-theoretic
differentiation on free finite-linear spectral models.

\subsection{Comparison with the spectral universal derivation}

\begin{definition}[Cofree function algebra]
\label{def:cofree-function-algebra}
Let \(B\) be connective and let \(M\) be a finite projective \(B\)-module.
Define
\[
    \mathcal O_!(M)
    :=
    \underline{\operatorname{Hom}}_B(!_BM,B).
\]
The cocommutative \(\mathbb E_\infty\)-coalgebra structure of \(C_B(M)\)
gives \(\mathcal O_!(M)\) its convolution commutative \(B\)-algebra structure.
The comonad counit \(\varepsilon_M:!_BM\to M\) induces a \(B\)-linear map
\[
    M^\vee\longrightarrow\mathcal O_!(M),
\]
and hence a canonical commutative \(B\)-algebra morphism
\[
    j_M:
    \operatorname{Sym}_B(M^\vee)
    \longrightarrow
    \mathcal O_!(M).
\]
\end{definition}

\begin{remark}[Stable algebraic nature of the cofree function algebra]
\label{rem:cofree-function-algebra-not-connective}
The convolution algebra \(\mathcal O_!(M)\) is formed as a commutative
algebra object in the unrestricted stable coefficient category
\(\Mod_B\); equivalently,
\[
    \mathcal O_!(M)
    \in
    \CAlg(\Mod_B)
    \simeq
    \CAlg(\Sp)_{B/},
\]
with its underlying \(B\)-module allowed to range over the unrestricted
stable coefficient category. The stability belongs to \(\Mod_B\); its
\(\infty\)-category of commutative algebra objects is generally nonstable.

The comparison in this subsection is therefore formulated in the unrestricted
commutative spectral-algebraic universe. A connective interpretation of
\(\mathcal O_!(M)\) as a spectral coordinate algebra or affine spectral
context requires a separate connectivity theorem. The comparison map \(j_M\)
below is an algebraic map in that unrestricted universe, in accordance with
the connective-coordinate versus stable-coefficient convention of this
manuscript.
\end{remark}

\begin{lemma}[Dual of the primitive infinitesimal coalgebra]
\label{lem:dual-primitive-square-zero}
Let \(M\) be finite projective over \(B\). Then \(P_B(M)=B\oplus M\) is
dualizable, and dualization carries its primitive cocommutative coalgebra
structure to the split square-zero commutative \(B\)-algebra
\[
    P_B(M)^\vee
    \simeq
    B\oplus M^\vee.
\]
For every \(B\)-module \(K\), there is a canonical equivalence
\[
    \underline{\operatorname{Hom}}_B(K\otimes_BP_B(M),B)
    \simeq
    \underline{\operatorname{Hom}}_B(K,B)
    \otimes_BP_B(M)^\vee.
\]
When \(K\) underlies a cocommutative coalgebra, this equivalence identifies the
convolution algebra on the left with the tensor-product commutative algebra on
the right.
\end{lemma}

\begin{proof}
\leavevmode
\begin{enumerate}[label=\textup{(\arabic*)}]
    \item \textbf{Finite projectivity.}
    Finite projectivity makes \(M\), and hence \(B\oplus M\), dualizable.

    \item \textbf{contravariant symmetric-monoidal duality functor on dualizable objects.}
    The contravariant symmetric-monoidal duality functor on dualizable objects carries
    cocommutative coalgebras to commutative algebras.

    \item \textbf{Derive the comparison.}
    Dualizing the primitive comultiplication gives multiplication which is the ordinary unit
    action on the \(M^\vee\)-summand and vanishes on \(M^\vee\otimes_BM^\vee\); this is exactly the split
    square-zero algebra.

    \item \textbf{For the internal-Hom formula, for every , duality and closedness.}
    For the internal-Hom formula, for every \(Z\in\Mod_B\), duality and closedness give
    \[
    \begin{aligned}
        \Map_{\Mod_B}\bigl(Z,
          \underline{\operatorname{Hom}}_B(K,B)\otimes_BP_B(M)^\vee\bigr)
        &\simeq
        \Map_{\Mod_B}\bigl(Z\otimes_BP_B(M),
          \underline{\operatorname{Hom}}_B(K,B)\bigr)
        \\
        &\simeq
        \Map_{\Mod_B}\bigl(Z\otimes_BP_B(M)\otimes_BK,B\bigr)
        \\
        &\simeq
        \Map_{\Mod_B}\bigl(Z,
          \underline{\operatorname{Hom}}_B(K\otimes_BP_B(M),B)\bigr).
    \end{aligned}
    \]

    \item \textbf{Apply Yoneda.}
    Yoneda gives the equivalence.

    \item \textbf{Compatibility with convolution.}
    Compatibility with convolution is the symmetric-monoidal compatibility of duality: tensor
    products of coalgebras are sent to tensor products of their dual convolution algebras
    whenever the relevant factor is dualizable.
\end{enumerate}
\end{proof}

\begin{construction}[Spectral convolution derivation]
\label{con:spectral-convolution-derivation}
Let
\[
    R:=\mathcal O_!(M).
\]
Precomposition with the coalgebra morphism
\[
    \widetilde{\mathsf d}_M:
    C_B(M)\times P_B(M)\longrightarrow C_B(M)
\]
gives a morphism of commutative \(B\)-algebras
\[
    R
    \longrightarrow
    \underline{\operatorname{Hom}}_B
    (!_BM\otimes_BP_B(M),B).
\]
By \Ref{lem:dual-primitive-square-zero}, identify the target canonically with
\[
    R\otimes_B(B\oplus M^\vee)
    \simeq
    R\oplus(R\otimes_BM^\vee).
\]
Denote the resulting commutative \(B\)-algebra morphism by
\[
    \widetilde D:
    R
    \longrightarrow
    R\oplus(R\otimes_BM^\vee).
\]
It is a section of the split square-zero augmentation
\[
    R\oplus(R\otimes_BM^\vee)\longrightarrow R.
\]
Indeed, restriction of
\(\widetilde{\mathsf d}_M\) along the group-like summand
\[
    C_B(M)
    \simeq
    C_B(M)\times B
    \xrightarrow{1\times\iota_0}
    C_B(M)\times P_B(M)
\]
is the identity coalgebra morphism: after composing with
\(\varepsilon_M\), its adjunct is \(\varepsilon_M\).

Thus
\[
    \widetilde D
    \in
    \operatorname{Der}_B
    \bigl(R,R\otimes_BM^\vee\bigr)
\]
in the sense of \Ref{def:spectral-derivation}. Write
\[
    D:
    R\longrightarrow R\otimes_BM^\vee
\]
for its coefficient component under the canonical decomposition
\[
    R\oplus(R\otimes_BM^\vee).
\]
Equivalently,
\[
    \widetilde D=(\id_R,D)
\]
on underlying \(B\)-modules.

Under the internal-Hom adjunction, the coefficient component \(D\) is
precisely precomposition with the deriving transformation
\[
    \mathsf d_M:
    !_BM\otimes_BM\longrightarrow !_BM.
\]
\end{construction}

\begin{theorem}[Compatibility with the spectral K\"ahler derivative]
\label{thm:dill-kahler-comparison}
Let \(M\) be finite projective over a connective commutative ring spectrum
\(B\), and put
\[
    A:=\operatorname{Sym}_B(M^\vee).
\]
Then
\[
    L_{A/B}\simeq A\otimes_BM^\vee.
\]
Let
\[
    \widetilde d_{\mathrm{univ}}:
    A
    \longrightarrow
    A\oplus(A\otimes_BM^\vee)
\]
be the universal \(B\)-linear derivation section, and write
\[
    d_{\mathrm{univ}}:
    A
    \longrightarrow
    A\otimes_BM^\vee
\]
for its coefficient component.

The coalgebra-level infinitesimal translation of
\Ref{con:tangent-coalgebra-infinitesimal-translation} induces the spectral
convolution derivation section
\[
    \widetilde D:
    \mathcal O_!(M)
    \longrightarrow
    \mathcal O_!(M)
    \oplus
    \bigl(\mathcal O_!(M)\otimes_BM^\vee\bigr)
\]
of \Ref{con:spectral-convolution-derivation}. Write
\[
    D:
    \mathcal O_!(M)
    \longrightarrow
    \mathcal O_!(M)\otimes_BM^\vee
\]
for its coefficient component. Then the square
\[
\begin{tikzcd}[column sep=huge,row sep=large]
    A \arrow[r,"j_M"] \arrow[d,"d_{\mathrm{univ}}"']
        & \mathcal O_!(M) \arrow[d,"D"]
    \\
    A\otimes_BM^\vee \arrow[r,"j_M\otimes1"']
        & \mathcal O_!(M)\otimes_BM^\vee
\end{tikzcd}
\]
commutes canonically.

Equivalently, the corresponding square of spectral derivation sections
\[
\begin{tikzcd}[column sep=huge,row sep=large]
    A \arrow[r,"j_M"] \arrow[d,"\widetilde d_{\mathrm{univ}}"']
        & \mathcal O_!(M) \arrow[d,"\widetilde D"]
    \\
    A\oplus(A\otimes_BM^\vee)
        \arrow[r,"{j_M\oplus(j_M\otimes1)}"']
        &
    \mathcal O_!(M)
        \oplus
        \bigl(\mathcal O_!(M)\otimes_BM^\vee\bigr)
\end{tikzcd}
\]
commutes canonically. Thus, on free finite-linear spectral models,
differential-linear differentiation is canonically compatible along \(j_M\)
with the universal spectral K\"ahler derivation.
\end{theorem}

\begin{proof}
\leavevmode
\begin{enumerate}[label=\textup{(\arabic*)}]
    \item \textbf{Apply the cotangent calculation.}
    The cotangent calculation is \Ref{prop:cotangent-free-algebra}, applied to \(M^\vee\).

    \item \textbf{Fix the notation.}
    Put
    \[
            R:=\mathcal O_!(M)
            =
            \underline{\operatorname{Hom}}_B(!_BM,B).
        \]

    \item \textbf{Derive the comparison.}
    Finite projectivity of \(M\) gives a canonical duality equivalence
    \[
            R\otimes_BM^\vee
            \simeq
            \underline{\operatorname{Hom}}_B
            (!_BM\otimes_BM,B).
        \]

    \item \textbf{Transport the structure.}
    Under this equivalence, the coefficient component
    \[
            D:R\longrightarrow R\otimes_BM^\vee
        \]
    of the convolution derivation section is precomposition with the deriving transformation
    \[
            \mathsf d_M:
            !_BM\otimes_BM\longrightarrow !_BM
        \]
    by \Ref{con:spectral-convolution-derivation}.

    \item \textbf{Fix the notation.}
    Let
    \[
            g:M^\vee\longrightarrow R
        \]
    be the generator map underlying
    \[
            j_M:
            \operatorname{Sym}_B(M^\vee)\longrightarrow R;
        \]
    it is induced by the comonad counit
    \[
            \varepsilon_M:!_BM\longrightarrow M.
        \]

    \item \textbf{Use the adjunction.}
    Under the displayed internal-Hom equivalence, the composite
    \[
            M^\vee
            \xrightarrow{g}
            R
            \xrightarrow{D}
            R\otimes_BM^\vee
        \]
    is adjoint to
    \[
        \begin{aligned}
            M^\vee\otimes_B!_BM\otimes_BM
            &\xrightarrow{\ 1\otimes\mathsf d_M\ }
            M^\vee\otimes_B!_BM
            \\
            &\xrightarrow{\ 1\otimes\varepsilon_M\ }
            M^\vee\otimes_BM
            \xrightarrow{\ \operatorname{ev}\ }
            B.
        \end{aligned}
        \]

    \item \textbf{Derive the comparison.}
    The defining adjunct equation for coalgebra-level infinitesimal translation, restricted to
    the primitive summand, gives the canonical homotopy
    \[
            \varepsilon_M\mathsf d_M
            \simeq
            e_M\otimes\id_M:
            !_BM\otimes_BM\longrightarrow M.
        \]

    \item \textbf{Use the adjunction.}
    Consequently the map induced on generators by \(Dg\) is canonically adjoint to
    \[
            !_BM\otimes_BM
            \xrightarrow{e_M\otimes\id_M}
            B\otimes_BM
            \simeq
            M.
        \]

    \item \textbf{Use the adjunction.}
    Since the convolution unit
    \[
            1_R:B\longrightarrow R
        \]
    is adjoint to the coalgebra counit
    \[
            e_M:!_BM\longrightarrow B,
        \]
    this is exactly the internal-Hom representative of the canonical map
    \[
            M^\vee
            \longrightarrow
            R\otimes_BM^\vee,
            \qquad
            \lambda\longmapsto1_R\otimes\lambda.
        \]

    \item \textbf{Identify the map.}
    The coefficient component of the universal derivation section
    \[
            \widetilde d_{\mathrm{univ}}:
            A
            \longrightarrow
            A\oplus(A\otimes_BM^\vee)
        \]
    is characterized by the generator map
    \[
            M^\vee
            \longrightarrow
            A\otimes_BM^\vee,
            \qquad
            \lambda\longmapsto1_A\otimes\lambda.
        \]

    \item \textbf{After extension of coefficients along this.}
    After extension of coefficients along
    \[
            j_M:A\longrightarrow R,
        \]
    this becomes
    \[
            M^\vee
            \longrightarrow
            R\otimes_BM^\vee,
            \qquad
            \lambda\longmapsto1_R\otimes\lambda.
        \]

    \item \textbf{Now regard as an -module through .}
    Now regard \(R\otimes_BM^\vee\) as an \(A\)-module through \(j_M\).

    \item \textbf{Apply the cotangent calculation.}
    The free cotangent calculation gives
    \[
            L_{A/B}
            \simeq
            A\otimes_BM^\vee,
        \]
    and therefore the cotangent classifier and extension--restriction adjunction give a natural
    equivalence
    \[
        \begin{aligned}
            \operatorname{Der}_B
            \bigl(A,R\otimes_BM^\vee\bigr)
            &\simeq
            \Map_{\Mod_A}
            \bigl(
                A\otimes_BM^\vee,
                R\otimes_BM^\vee
            \bigr)
            \\
            &\simeq
            \Map_{\Mod_B}
            \bigl(
                M^\vee,
                R\otimes_BM^\vee
            \bigr).
        \end{aligned}
        \]
    Under this equivalence, a \(B\)-linear derivation from the free algebra \(A=\operatorname{Sym}_B(M^\vee)\) is
    determined precisely by its value on the free generator \(M^\vee\).

    \item \textbf{Identify the map.}
    The two \(B\)-linear derivations
    \[
            D\circ j_M
            \qquad\text{and}\qquad
            (j_M\otimes1)\circ d_{\mathrm{univ}}
        \]
    have just been shown to correspond to the same generator map
    \[
            M^\vee
            \longrightarrow
            R\otimes_BM^\vee,
            \qquad
            \lambda\longmapsto1_R\otimes\lambda.
        \]

    \item \textbf{Identify the equivalence.}
    The identity path of this generator map therefore transports through the preceding
    equivalence of derivation spaces to a canonical homotopy
    \[
            D\circ j_M
            \simeq
            (j_M\otimes1)\circ d_{\mathrm{univ}}.
        \]

    \item \textbf{Verify compatibility.}
    This is the canonical commuting square of coefficient components displayed in the theorem.

    \item \textbf{Conclude the argument.}
    Finally, both square-zero derivation sections have the corresponding algebra map as their
    first component and the preceding derivation as their coefficient component.

    \item \textbf{Verify compatibility.}
    Hence the same canonical homotopy gives the commuting square
    \[
        \begin{tikzcd}[column sep=huge,row sep=large]
            A \arrow[r,"j_M"] \arrow[d,"\widetilde d_{\mathrm{univ}}"']
                & R \arrow[d,"\widetilde D"]
            \\
            A\oplus(A\otimes_BM^\vee)
                \arrow[r,"{j_M\oplus(j_M\otimes1)}"']
                &
            R\oplus(R\otimes_BM^\vee).
        \end{tikzcd}
        \]

    \item \textbf{No finite free retract, basis, or other presentation of.}
    No finite free retract, basis, or other presentation of \(M\) is used.
\end{enumerate}
\end{proof}
The final issue for the linear branch is substitution: pullback transports the
cofree differential structure through a canonical comparison, but that
comparison need not be invertible.

\subsection{Substitution compatibility of the linear exponential}

\begin{construction}[Oplax substitution comparison for cofree coalgebras]
\label{con:cofree-coalgebra-substitution-comparison}
Let \(f:Y\to X\) be a morphism of spectral schemes. Strong symmetric
monoidality of
\[
    f^*:\QCoh(X)\longrightarrow\QCoh(Y)
\]
induces a functor on cocommutative coalgebras. For
\(M\in\QCoh(X)\), the morphism
\[
    f^*(!_XM)\xrightarrow{f^*\varepsilon_M}f^*M
\]
therefore classifies, by the cofree universal property in \(\QCoh(Y)\), a
canonical coalgebra morphism
\[
    f^*C_X(M)\longrightarrow C_Y(f^*M).
\]
On underlying coefficients this is a natural comparison
\[
    \beta_{f,M}:
    f^*(!_XM)
    \longrightarrow
    !_Y(f^*M).
\]
\end{construction}

\begin{proposition}[Oplax substitution compatibility of the differential exponential]
\label{prop:cofree-substitution-coherence}
The comparisons \(\beta_{f,M}\) are coherent under identities, composition,
and higher homotopies of substitution. They satisfy the following canonical
compatibilities, with the strong symmetric-monoidal identifications for
\(f^*\) suppressed from notation. Superscripts \(X\) and \(Y\) distinguish
the corresponding comonad, monoidal, coalgebra, Seely, and bialgebra maps in
the two coefficient fibres.
\begin{enumerate}[label=(\roman*)]
    \item \textbf{Comonad.}
    \[
        \varepsilon^{Y}_{f^*M}\,\beta_{f,M}
        \simeq
        f^*(\varepsilon^{X}_M),
    \]
    and
    \[
    \begin{aligned}
        \delta^{Y}_{f^*M}\,\beta_{f,M}
        \simeq
        !_Y(\beta_{f,M})\,
        \beta_{f,!_XM}\,
        f^*(\delta^{X}_M).
    \end{aligned}
    \]

    \item \textbf{Monoidal structure.}
    For \(M,N\in\QCoh(X)\),
    \[
    \begin{aligned}
        \beta_{f,M\otimes N}\,f^*(m^X_{M,N})
        \simeq
        m^Y_{f^*M,f^*N}
        (\beta_{f,M}\otimes\beta_{f,N}),
    \end{aligned}
    \]
    and the analogous unit comparison holds for
    \(m_{\mathcal O_X}\) and \(m_{\mathcal O_Y}\).

    \item \textbf{Cocommutative coalgebra structure.}
    Since \(\beta_{f,M}\) underlies a coalgebra morphism,
    \[
        \Delta^Y_{f^*M}\,\beta_{f,M}
        \simeq
        (\beta_{f,M}\otimes\beta_{f,M})f^*(\Delta^X_M),
        \qquad
        e^Y_{f^*M}\,\beta_{f,M}
        \simeq
        f^*(e^X_M).
    \]

    \item \textbf{Seely and additive bialgebra structure.}
    The Seely comparisons satisfy
    \[
    \begin{aligned}
        \chi^Y_{f^*M,f^*N}\,\beta_{f,M\oplus N}
        \simeq
        (\beta_{f,M}\otimes\beta_{f,N})\,f^*(\chi^X_{M,N}),
    \end{aligned}
    \]
    and the corresponding unit comparison holds. Moreover
    \[
        \beta_{f,M}\,f^*(\nabla^X_M)
        \simeq
        \nabla^Y_{f^*M}(\beta_{f,M}\otimes\beta_{f,M}),
        \qquad
        \beta_{f,M}\,f^*(u^X_M)
        \simeq
        u^Y_{f^*M}.
    \]

    \item \textbf{Primitive infinitesimals and codereliction.}
    Strong symmetric monoidality and exactness identify
    \[
        f^*P_X(M)\simeq P_Y(f^*M),
    \]
    and under this identification
    \[
        \beta_{f,\mathcal O_X\oplus M}\,f^*(H^{\inf}_M)
        \simeq
        H^{\inf}_{f^*M}.
    \]
    Consequently
    \[
        \beta_{f,M}\,f^*(\eta_M)
        \simeq
        \eta_{f^*M}.
    \]

    \item \textbf{Infinitesimal translation and differentiation.}
    The coalgebra-level infinitesimal translations satisfy
    \[
    \begin{aligned}
        \beta_{f,M}\circ f^*\bigl(U_X\widetilde{\mathsf d}_M\bigr)
        \simeq
        U_Y\widetilde{\mathsf d}_{f^*M}
        \circ
        (\beta_{f,M}\otimes1),
    \end{aligned}
    \]
    as maps
    \[
        f^*(!_XM)\otimes P_Y(f^*M)
        \longrightarrow
        !_Y(f^*M).
    \]
    Restricting to the primitive summand gives
    \[
        \beta_{f,M}\,f^*(\mathsf d_M)
        \simeq
        \mathsf d_{f^*M}(\beta_{f,M}\otimes1_{f^*M}).
    \]
\end{enumerate}
Consequently the stable coefficient fibres carry an oplax indexed family of the
coherent differential-exponential structures constructed above; no
invertibility of \(\beta\) is used.
\end{proposition}

\begin{proof}
\leavevmode
\begin{enumerate}[label=\textup{(\arabic*)}]
    \item \textbf{Use the adjunction.}
    The comparison \(\beta_{f,M}\) is the underlying map of the coalgebra morphism adjoint to
    \(f^*\varepsilon^X_M\).

    \item \textbf{Its defining adjunct equation.}
    Its defining adjunct equation is exactly the first formula in \textup{(i)}.

    \item \textbf{For comultiplication, both sides of the second formula.}
    For comultiplication, both sides of the second formula underlie coalgebra maps
    \[
        f^*C_X(M)\longrightarrow C_Y(!_Yf^*M).
    \]

    \item \textbf{Verify naturality.}
    After composition with the cofree counit in the target, the left side reduces to \(\beta_{f,M}\)
    by the comonad counit law, while the right side reduces to the same map by naturality of
    \(\varepsilon^Y\), the defining adjunct equation for \(\beta_{f,!_XM}\), and the comonad counit law for
    \(!_X\).

    \item \textbf{Verify naturality.}
    Cofreeness therefore gives the comparison; identity, composition, and higher-pasting
    coherences follow from the same adjunction naturality.

    \item \textbf{For (ii), both sides.}
    For \textup{(ii)}, both sides are coalgebra maps from \(f^*C_X(M)\times f^*C_X(N)\) to the cofree coalgebra
    \(C_Y(f^*M\otimes f^*N)\).

    \item \textbf{Derive the comparison.}
    Composing with the target counit gives on both sides
    \[
        f^*(\varepsilon^X_M)\otimes f^*(\varepsilon^X_N).
    \]

    \item \textbf{Use the cofree universal property.}
    The cofree mapping-space equivalence therefore gives the monoidal comparison.

    \item \textbf{unit case.}
    The unit case is identical.

    \item \textbf{Part (iii).}
    Part \textup{(iii)} is immediate because the map defining \(\beta\) is a morphism in
    \(\operatorname{CoCom}(\QCoh(Y))\).

    \item \textbf{Use the product structure.}
    For \textup{(iv)}, the Seely comparison is induced by the Cartesian product comparison
    \[
        C_X(M\oplus N)\xrightarrow{\simeq}C_X(M)\times C_X(N).
    \]

    \item \textbf{Apply and compare with the corresponding product comparison in the.}
    Apply \(f^*\) and compare with the corresponding product comparison in the target fibre.

    \item \textbf{Apply the universal property.}
    Both maps to each product factor have the same cofree adjunct, so the product universal
    property gives the displayed compatibility with \(\chi\).

    \item \textbf{Use exactness.}
    Exactness of \(f^*\) preserves biproduct addition and zero.

    \item \textbf{Use the cofree universal property.}
    Since \(\widetilde\nabla_M\) and \(\widetilde u_M\) are the cofree coalgebra lifts of the maps induced by
    addition and zero, the same cofree-adjunct comparison gives compatibility with \(\nabla\)
    and \(u\).

    \item \textbf{For (v), the primitive coalgebra construction.}
    For \textup{(v)}, the primitive coalgebra construction is the opposite split square-zero
    algebra construction.

    \item \textbf{Use the product structure.}
    Strong symmetric monoidality and preservation of finite biproducts by \(f^*\) therefore
    give \(f^*P_X(M)\simeq P_Y(f^*M)\).

    \item \textbf{Compare the two coalgebra.}
    Compare the two coalgebra maps
    \[
        P_Y(f^*M)\rightrightarrows C_Y(\mathcal O_Y\oplus f^*M)
    \]
    given by \(\alpha_{P_Y(f^*M)}\) and by the pullback of \(\alpha_{P_X(M)}\) followed by the cofree comparison
    \(\beta\).

    \item \textbf{Use the cofree universal property.}
    After composition with the cofree counit, both have adjunct \(\id_{\mathcal O_Y\oplus f^*M}\).

    \item \textbf{Conclude the step.}
    Hence
    \[
        \beta_{f,\mathcal O_X\oplus M}f^*(H^{\inf}_M)
        \simeq
        H^{\inf}_{f^*M}.
    \]

    \item \textbf{Next compare the two coalgebra.}
    Next compare the two coalgebra maps
    \[
        P_Y(f^*M)\rightrightarrows C_Y(f^*M)
    \]
    given respectively by \(H_{f^*M}\) and by
    \[
        f^*P_X(M)\xrightarrow{f^*H_M}f^*C_X(M)
        \longrightarrow C_Y(f^*M).
    \]

    \item \textbf{Use the cofree universal property.}
    After composition with the cofree counit both have adjunct \(\operatorname{pr}_1\).

    \item \textbf{Conclude the step.}
    Hence they agree canonically.

    \item \textbf{Restrict the construction.}
    Restriction along the primitive inclusion gives
    \[
        \beta_{f,M}f^*(\eta_M)\simeq\eta_{f^*M}.
    \]

    \item \textbf{For (vi), both displayed composites.}
    For \textup{(vi)}, both displayed composites are coalgebra maps into the cofree coalgebra
    \(C_Y(f^*M)\).

    \item \textbf{Compose them with .}
    Compose them with \(\varepsilon^Y_{f^*M}\).

    \item \textbf{For the first composite, the defining properties of.}
    For the first composite, the defining properties of \(\beta\) and \(\widetilde{\mathsf d}_M\) give the
    pullback of
    \[
        e^X_M\otimes\operatorname{pr}_1+\varepsilon^X_M\otimes\operatorname{pr}_0.
    \]

    \item \textbf{Derive the comparison.}
    For the second, part \textup{(iii)} gives preservation of the coalgebra counit, while part
    \textup{(i)} gives \(\varepsilon^Y_{f^*M}\beta_{f,M}\simeq f^*\varepsilon^X_M\).

    \item \textbf{Its adjunct.}
    Its adjunct is therefore the same map.

    \item \textbf{Derive the comparison.}
    Cofreeness gives the canonical comparison, coherently in \(f\) and \(M\).

    \item \textbf{Restrict the construction.}
    Restriction along the primitive summand proves the final formula.
\end{enumerate}
\end{proof}

\begin{remark}[No indexed equivalence is asserted]
\label{rem:cofree-substitution-not-equivalence}
The canonical substitution structure proved here is oplax indexed, with the
comparison
\[
    \beta_{f,M}:f^*(!_XM)\longrightarrow !_Y(f^*M).
\]
An equivalence for this comparison would require an additional base-change
theorem for cofree spectral coalgebras. Cofree formation is controlled by
cosifted limits, and spectral Tate phenomena obstruct the naive interchange of
those limits with tensor products in general; see
\cite{BachmannBurklund2026Coalgebras}. The differential compatibility above
is formulated with the actual comparison \(\beta\).
\end{remark}
The linear branch therefore exports three facts: every coefficient fibre
carries the cofree differential-exponential structure, its derivative agrees
with the spectral universal derivation on the stated finite-projective models,
and substitution acts through the canonical oplax comparison. The geometric
branch does not depend on this construction; it returns to the cotangent
classifier and uses degree-one classes to build nonlinear contexts.

\section{Structured infinitesimal context extension}
\label{sec:structured-infinitesimal-context-extension}

A degree-one derivation term is converted functorially into a square-zero
thickening whose centre is the original context. The conceptual passage is

\[
\begin{tikzcd}[column sep=huge,row sep=large]
L_{X/S} \arrow[r,"\eta"]
    \arrow[d,phantom,"{\text{classifying coefficient}}" description]
& M[1]
    \arrow[d,phantom,"{\text{square-zero construction}}" description]
\\
X \arrow[r,hook,"i_\eta"'] & X_\eta .
\end{tikzcd}
\]

This is analogous in purpose to context formation, but it has different
variance from the affine comprehension of
Chapter~\Ref{chap:spectral-algebraic-geometry}. The section first constructs
the affine structured extension and distinguishes its three base-change
operations; it then proves the Zariski compatibility required to glue the
construction intrinsically over an arbitrary spectral context.

\begin{convention}[Structured and underlying infinitesimal arrows]
\label{conv:structured-underlying-infinitesimal}
A \textbf{structured} square-zero extension is determined by its degree-one
derivation together with the coefficient occurring in that derivation; its
square-zero algebra, augmentation, and pullback comparison are constructed
functorially. An \textbf{underlying} square-zero extension is an arrow in the
equivalence-closed essential image of that construction. Passing to the
underlying arrow forgets the coefficient and classifying derivation. A
coefficient, factorization, filtration, pullback model, local equation,
relative refinement, or representative chosen later inside a proof is a proof
witness and is not retained on the underlying arrow.
\end{convention}

\subsection{Affine structured square-zero extensions}

\begin{construction}[Structured square-zero extension classified by a derivation]
\label{con:structured-square-zero-extension}
Let \(A\to B\) be a morphism of commutative ring spectra, let
\(M\in\Mod_B\), and let
\[
    \eta:L_{B/A}\longrightarrow M[1].
\]
Under \Ref{thm:relative-cotangent-classifier}, \(\eta\) represents an
\(A\)-linear derivation
\[
    d_\eta:B\longrightarrow B\oplus M[1].
\]
Let
\[
    d_0:B\longrightarrow B\oplus M[1]
\]
be the zero derivation and let
\[
    p:B\oplus M[1]\longrightarrow B
\]
be the split augmentation. Put
\[
    \mathcal A:=\CAlg(\Sp)_{A/}.
\]
Because
\[
    p\,d_\eta\simeq\id_B\simeq p\,d_0,
\]
the maps \(d_\eta\) and \(d_0\) are morphisms
\[
    \id_B\rightrightarrows p
\]
in the slice \(\mathcal A_{/B}\). Define the structured extension arrow by
the pullback in that slice
\[
    (q_\eta:B_\eta\to B)
    :=
    \id_B\times_p\id_B.
\]
Since limits in a slice are created in the ambient category, its underlying
commutative \(A\)-algebra is
\[
    B_\eta
    \simeq
    B\times_{B\oplus M[1]}B,
\]
with specified pullback square
\[
\begin{tikzcd}[column sep=large, row sep=large]
    B_\eta \arrow[r,"q_\eta"] \arrow[d,"r_\eta"']
        & B \arrow[d,"d_\eta"]
    \\
    B \arrow[r,"d_0"']
        & B\oplus M[1].
\end{tikzcd}
\]
Applying the augmentation \(p\) to the pullback comparison gives a canonical
homotopy
\[
    q_\eta\simeq r_\eta.
\]
The classified square-zero arrow is \(q_\eta:B_\eta\to B\); the second
projection and the displayed homotopy are functorially determined by the
pullback. The derivation \(\eta\), its coefficient \(M\), and this
functorially constructed arrow form the
\textbf{structured square-zero extension of \(B\) by \(M\) over \(A\)}
classified by \(\eta\).
\end{construction}

\begin{proposition}[Fibre and connectivity of a structured extension]
\label{prop:structured-square-zero-fibre}
In \Ref{con:structured-square-zero-extension}, the homotopy fibre of
\[
    q_\eta:B_\eta\longrightarrow B
\]
is canonically equivalent to \(M\), regarded as a \(B_\eta\)-module by
restriction of scalars. The induced multiplication on that fibre is
coherently null. If \(B\) and \(M\) are connective, then \(B_\eta\) is
connective, \(\pi_0(q_\eta)\) is surjective, and its kernel is square-zero.
\end{proposition}

\begin{proof}
\leavevmode
\begin{enumerate}[label=\textup{(\arabic*)}]
    \item \textbf{forgetful functor creates limits, so the defining square.}
    The forgetful functor \(\CAlg(\Sp)\to\Sp\) creates limits, so the defining square is Cartesian on
    underlying spectra.

    \item \textbf{Transport the structure.}
    Under the canonical decomposition
    \[
        B\oplus M[1],
    \]
    the two sections have the form
    \[
        d_\eta=(\id_B,\delta_\eta),
        \qquad
        d_0=(\id_B,0),
    \]
    where \(\delta_\eta:B\to M[1]\) is the underlying relative derivation.

    \item \textbf{Compute the fibre.}
    Consequently the pullback is canonically equivalent, compatibly with \(q_\eta\), to the
    homotopy fibre of \(\delta_\eta\):
    \[
        B_\eta
        \simeq
        \operatorname{fib}(\delta_\eta)
        \longrightarrow B.
    \]

    \item \textbf{Derive the comparison.}
    Equivalently, pullback invariance of fibres applied to the horizontal maps in the defining
    square gives
    \[
        \operatorname{fib}(q_\eta)
        \simeq
        \operatorname{fib}(d_0)
        \simeq
        \Omega M[1]
        \simeq
        M.
    \]

    \item \textbf{Compute the fibre.}
    Thus there is a canonical fibre sequence
    \[
        M\longrightarrow B_\eta
        \xrightarrow{q_\eta}B
        \xrightarrow{\delta_\eta}M[1].
    \]

    \item \textbf{tangent-theoretic derivation-to-extension construction equips the fibre with zero.}
    The tangent-theoretic derivation-to-extension construction equips the fibre with zero
    nonunital commutative multiplication; equivalently
    \[
        M\otimes_{B_\eta}M\longrightarrow M
    \]
    is coherently null.

    \item \textbf{Compute the fibre.}
    If \(B\) and \(M\) are connective, the displayed fibre sequence makes
    \(B_\eta\) connective.

    \item \textbf{Verify surjectivity.}
    Its long exact homotopy sequence gives surjectivity on \(\pi_0\), and the kernel is the
    image of \(\pi_0(M)\).

    \item \textbf{null multiplication on the fibre.}
    The null multiplication on the fibre makes the product of any two elements of that image
    zero.
\end{enumerate}
\end{proof}

\begin{theorem}[Structured square-zero parameterization]
\label{thm:structured-square-zero-classification}
Let \(A\to B\) be a morphism of commutative ring spectra and let
\(M\in\Mod_B\). Let
\[
    \operatorname{SqZ}^{\mathrm{str}}_A(B;M)
\]
be the \(\infty\)-groupoid whose objects are derivation terms
\(\eta:L_{B/A}\to M[1]\), read through the associated square-zero arrows of
\Ref{con:structured-square-zero-extension}, and whose morphisms are induced by
homotopies of classifying derivations. The square-zero algebra, augmentation,
and pullback identification are outputs of the construction rather than
further fields of data. Then there are natural equivalences
\[
\begin{aligned}
    \operatorname{SqZ}^{\mathrm{str}}_A(B;M)
    &\simeq
    \operatorname{Der}_A(B,M[1])
    \\
    &\simeq
    \Map_{\Mod_B}(L_{B/A},M[1]).
\end{aligned}
\]
\end{theorem}

\begin{proof}
\leavevmode
\begin{enumerate}[label=\textup{(\arabic*)}]
    \item \textbf{Handle the second component.}
    The second equivalence is \Ref{thm:relative-cotangent-classifier}.

    \item \textbf{Handle the first component.}
    The first is the functorial pullback construction of
    \Ref{con:structured-square-zero-extension}.

    \item \textbf{Identify contractibility.}
    For a fixed derivation, the space of pullback objects equipped with an identification with
    the specified pullback is contractible.

    \item \textbf{Conclude the step.}
    Hence adjoining the constructed square-zero arrow adds no independent structure.
\end{enumerate}
\end{proof}

\begin{lemma}[Square-zero boundary and the relative Hurewicz morphism]
\label{lem:square-zero-boundary-hurewicz}
Let \(A\to B\) be a morphism of commutative ring spectra, let
\(M\in\Mod_B\), and let
\[
    \eta:L_{B/A}\longrightarrow M[1]
\]
classify the structured square-zero extension
\[
    q_\eta:B_\eta\longrightarrow B.
\]
Put
\[
    Q_\eta:=\operatorname{cofib}(A\to B_\eta),
    \qquad
    Q:=\operatorname{cofib}(A\to B).
\]
Octahedrality for
\(A\to B_\eta\to B\) supplies a cofibre sequence
\[
    Q_\eta
    \longrightarrow
    Q
    \xrightarrow{\beta_\eta}
    M[1].
\]
Then \(\beta_\eta\) is canonically the composite
\[
    Q
    \longrightarrow
    B\otimes_AQ
    \xrightarrow{\epsilon_{B/A}}
    L_{B/A}
    \xrightarrow{\eta}
    M[1].
\]
The identification is natural in the relative algebra, the coefficient, and
the classifying derivation.
\end{lemma}

\begin{proof}
\leavevmode
\begin{enumerate}[label=\textup{(\arabic*)}]
    \item \textbf{Fix the notation.}
    Let
    \[
        \delta_\eta:B\longrightarrow M[1]
    \]
    be the underlying \(A\)-linear morphism obtained by subtracting the zero derivation
    from the derivation section:
    \[
        d_\eta-d_0=(0,\delta_\eta).
    \]

    \item \textbf{Invoke the cited result.}
    By \Ref{prop:structured-square-zero-fibre}, the slice pullback construction is canonically
    the fibre construction for \(\delta_\eta\), compatibly with the classified projection
    \(q_\eta\).

    \item \textbf{Compute the fibre.}
    Hence there is a canonical fibre--cofibre sequence
    \[
        M
        \longrightarrow
        B_\eta
        \xrightarrow{q_\eta}
        B
        \xrightarrow{\delta_\eta}
        M[1].
    \]

    \item \textbf{Restrict the construction.}
    The map \(\delta_\eta\) is canonically null after restriction to \(A\), because
    \(d_\eta\) is an \(A\)-linear derivation.

    \item \textbf{Construct the factorization.}
    It therefore factors through
    \[
        Q=\operatorname{cofib}(A\to B)
        \longrightarrow
        M[1].
    \]

    \item \textbf{Compute the fibre.}
    Octahedrality for \(A\to B_\eta\xrightarrow{q_\eta}B\) identifies this factor with the boundary
    \[
        \beta_\eta:Q\to M[1]
    \]
    in the displayed cofibre sequence of the statement.

    \item \textbf{Construct the factorization.}
    By the characterization of \(\epsilon_{B/A}\) in \Ref{con:relative-hurewicz-morphism}, the factor
    of the underlying relative derivation \(\delta_\eta\) through \(Q\) is exactly the
    composite
    \[
        Q\longrightarrow
        B\otimes_AQ
        \xrightarrow{\epsilon_{B/A}}
        L_{B/A}
        \xrightarrow{\eta}
        M[1].
    \]

    \item \textbf{Conclude the step.}
    Thus this composite is \(\beta_\eta\).

    \item \textbf{Verify naturality.}
    Naturality is inherited from the functorial square-zero pullback, cofibre, and Hurewicz
    constructions.
\end{enumerate}
\end{proof}

\begin{proposition}[Relative-base change of structured square-zero extensions]
\label{prop:structured-square-zero-base-change}
Let \(A\to B\) and \(A\to A'\) be morphisms of commutative ring spectra, put
\[
    B':=B\otimes_AA',
    \qquad
    M':=M\otimes_BB',
\]
and let
\[
    \eta:L_{B/A}\longrightarrow M[1]
\]
classify a structured square-zero extension \(B_\eta\to B\). Let
\[
    \eta':L_{B'/A'}
    \simeq
    L_{B/A}\otimes_BB'
    \longrightarrow
    M'[1]
\]
be its cotangent base change. Then there is a canonical equivalence of
augmented \(A'\)-algebras
\[
    B_\eta\otimes_AA'
    \xrightarrow{\simeq}
    B'_{\eta'}
\]
compatible with the maps to \(B'\). These equivalences are coherent under
iterated base change.
\end{proposition}

\begin{proof}
\leavevmode
\begin{enumerate}[label=\textup{(\arabic*)}]
    \item \textbf{Fix the notation.}
    Write the defining structured square-zero pullback as
    \[
    \begin{tikzcd}[column sep=large,row sep=large]
        B_\eta \arrow[r] \arrow[d]
            & B \arrow[d,"d_\eta"]
        \\
        B \arrow[r,"d_0"']
            & B\oplus M[1].
    \end{tikzcd}
    \]

    \item \textbf{Extension of scalars from to.}
    Extension of scalars from \(A\) to \(A'\) is symmetric monoidal.

    \item \textbf{Apply cotangent base change.}
    Hence it carries the split square-zero algebra to the split square-zero algebra
    \[
        (B\oplus M[1])\otimes_AA'
        \simeq
        B'\oplus M'[1],
    \]
    and cotangent base change identifies the scalar extension of the section \(d_\eta\) with
    the section \(d_{\eta'}\).

    \item \textbf{zero section.}
    The zero section is carried to the zero section.

    \item \textbf{Apply base change.}
    Thus the base-changed square supplies a canonical morphism
    \[
        B_\eta\otimes_AA'
        \longrightarrow
        B'\times_{B'\oplus M'[1]}B'
        =B'_{\eta'}.
    \]

    \item \textbf{It remains to.}
    It remains to prove that this morphism is an equivalence; no preservation of arbitrary
    pullbacks by algebraic extension of scalars is being assumed.

    \item \textbf{forgetful functor from commutative algebras to module.}
    The forgetful functor from commutative algebras to module spectra creates limits.

    \item \textbf{Compute the fibre.}
    Consequently the defining pullback has an underlying fibre sequence
    \[
        B_\eta
        \longrightarrow
        B\oplus B
        \xrightarrow{\ d_\eta-d_0\ }
        B\oplus M[1].
    \]

    \item \textbf{Compute the fibre.}
    Tensoring this fibre sequence over \(A\) with \(A'\) is exact in module spectra
    and gives
    \[
        B_\eta\otimes_AA'
        \longrightarrow
        B'\oplus B'
        \xrightarrow{\ d_{\eta'}-d_0\ }
        B'\oplus M'[1].
    \]

    \item \textbf{fibre of the final arrow.}
    The fibre of the final arrow is exactly the underlying spectrum of \(B'_{\eta'}\).

    \item \textbf{Conclude the step.}
    Therefore the canonical algebra morphism above is an equivalence on underlying spectra and
    hence an equivalence of commutative \(A'\)-algebras.

    \item \textbf{Apply cotangent base change.}
    Every map used in the construction is the canonical scalar-extension, cotangent-base-change,
    split-square-zero, or finite-limit comparison.

    \item \textbf{Apply base change.}
    Their identity, composition, and higher-pasting coherences therefore give the stated
    coherence under iterated base change.
\end{enumerate}
\end{proof}



\begin{lemma}[Algebraic pullback of structured square-zero extensions]
\label{lem:structured-square-zero-cartesian-pullback}
Let \(A\to D\) be a morphism of commutative ring spectra and let
\[
    q_\eta:E\longrightarrow D
\]
be the structured square-zero extension classified by
\[
    \eta:L_{D/A}\longrightarrow M[1],
    \qquad M\in\Mod_D.
\]
For every morphism of commutative \(A\)-algebras
\[
    u:D'\longrightarrow D,
\]
put
\[
    E':=E\times_DD'.
\]
Then \(E'\to D'\) is canonically a structured square-zero extension with
coefficient \(M\) regarded as a \(D'\)-module by restriction of scalars. The
construction is coherent under identities, composition, and iterated
algebraic pullback. If \(D'\) and the restricted coefficient are connective,
then \(E'\) is connective.
\end{lemma}

\begin{proof}
\leavevmode
\begin{enumerate}[label=\textup{(\arabic*)}]
    \item \textbf{Fix the notation.}
    Put \(\mathcal A:=\CAlg(\Sp)_{A/}\).

    \item \textbf{Identify the map.}
    The structured extension is the finite limit
    \[
        q_\eta
        \simeq
        \id_D\times_p\id_D
    \]
    in the slice \(\mathcal A_{/D}\), where
    \[
        p:D\oplus M[1]\longrightarrow D
    \]
    is the split augmentation and the two maps \(\id_D\rightrightarrows p\) are the derivation section
    \(d_\eta\) and the zero section.

    \item \textbf{Apply base change.}
    Base change along \(u\) is the pullback functor
    \[
        u^*:\mathcal A_{/D}\longrightarrow\mathcal A_{/D'},
        \qquad
        (C\to D)\longmapsto(C\times_DD'\to D').
    \]

    \item \textbf{Use the adjunction.}
    It is right adjoint to postcomposition along \(u\), hence preserves limits.

    \item \textbf{Conclude the step.}
    Therefore
    \[
    \begin{aligned}
        u^*q_\eta
        &\simeq
        u^*(\id_D\times_p\id_D)
        \\
        &\simeq
        u^*\id_D\times_{u^*p}u^*\id_D
        \\
        &\simeq
        \id_{D'}\times_{u^*p}\id_{D'}.
    \end{aligned}
    \]

    \item \textbf{pullback of the split augmentation.}
    The pullback of the split augmentation is canonically
    \[
        u^*p:
        D'\oplus\operatorname{Res}^{D}_{D'}M[1]
        \longrightarrow D'.
    \]

    \item \textbf{Verify compatibility.}
    Indeed, limits of commutative ring spectra are created on underlying spectra, and the
    pullback of \(D\oplus M[1]\to D\) has underlying \(D'\)-module \(D'\oplus\operatorname{Res}^{D}_{D'}M[1]\), with the functorially
    pulled-back square-zero multiplication.

    \item \textbf{two pulled-back sections.}
    The two pulled-back sections are respectively the pulled-back derivation and the zero
    derivation.

    \item \textbf{Conclude the step.}
    Consequently
    \[
        E'
        \simeq
        D'\times_{D'\oplus\operatorname{Res}^{D}_{D'}M[1]}D',
    \]
    which is the structured square-zero extension classified by that pulled-back derivation.

    \item \textbf{coefficient fibre.}
    The coefficient fibre is the restricted \(D'\)-module \(M\).

    \item \textbf{Compute the fibre.}
    Connectivity follows from the resulting fibre sequence.

    \item \textbf{Apply base change.}
    Identity, composition, and all higher-pasting coherences are those of functorial slice base
    change and therefore require no additional compatibility datum.
\end{enumerate}
\end{proof}

\begin{theorem}[Geometric Cartesian base change of structured square-zero extensions]
\label{thm:structured-square-zero-geometric-base-change}
Let
\[
    q:E\longrightarrow B
\]
be a structured square-zero extension of connective commutative ring spectra
with connective coefficient \(M\in\Mod_B\). For every morphism of connective
commutative ring spectra
\[
    E\longrightarrow C,
\]
put
\[
    D:=B\otimes_EC,
    \qquad
    M_D:=M\otimes_BD.
\]
Then the pushout arrow
\[
    C\longrightarrow D
\]
is canonically a structured square-zero extension with coefficient \(M_D\).
Equivalently, the Cartesian base change of
\(\Spec(B)\to\Spec(E)\) along \(\Spec(C)\to\Spec(E)\) is a structured
square-zero infinitesimal substitution. These transported structures are
coherent under identities, composition, and iterated geometric Cartesian base
change.
\end{theorem}

\begin{proof}
\leavevmode
\begin{enumerate}[label=\textup{(\arabic*)}]
    \item \textbf{Forget any relative base that may.}
    Forget any relative base that may have been used to construct \(q\).

    \item \textbf{Verify connectivity.}
    A structured square-zero extension of connective commutative ring spectra is a point of the
    connective derivation category of \cite{LurieDAGIV}.

    \item \textbf{Verify connectivity.}
    In the notation of that source, the functor from connective derivations to connective
    algebra arrows factors through a left fibration over the category whose morphisms are
    pushout squares; see \cite[Theorem~3.25]{LurieDAGIV}.

    \item \textbf{Verify connectivity.}
    For the displayed pushout
    \[
    \begin{tikzcd}[column sep=large,row sep=large]
        E \arrow[r] \arrow[d] & B \arrow[d] \\
        C \arrow[r] & D=B\otimes_EC,
    \end{tikzcd}
    \]
    the ring spectrum \(D\) is connective.

    \item \textbf{Verify connectivity.}
    Moreover connective pushout gives
    \[
        \pi_0(D)
        \simeq
        \pi_0(B)\otimes_{\pi_0(E)}\pi_0(C),
    \]
    and \(\pi_0(E)\to\pi_0(B)\) is surjective by \Ref{prop:structured-square-zero-fibre}.

    \item \textbf{Verify surjectivity.}
    Hence \(\pi_0(C)\to\pi_0(D)\) is surjective.

    \item \textbf{Verify connectivity.}
    Thus the target arrow belongs to the connective, surjective-on-\(\pi_0\) arrow category
    and the displayed square is a morphism in the pushout-square subcategory used in
    \cite[Theorem~3.25]{LurieDAGIV}.

    \item \textbf{Left-fibration transport along this morphism therefore produces.}
    Left-fibration transport along this morphism therefore produces a structured square-zero
    presentation of \(C\to D\).

    \item \textbf{transported tangent coefficient.}
    The transported tangent coefficient is extension of scalars and is canonically
    \[
        M\otimes_BD=M_D.
    \]

    \item \textbf{left-fibration transport.}
    The left-fibration transport supplies the identity, composition, and iterated-pushout
    coherences.

    \item \textbf{If the square.}
    If the square is equipped with a map from an additional base algebra, that map is retained
    on the pushed-out arrow; when a relative classifying derivation is required, it is
    constructed from this base structure by \Ref{lem:underlying-square-zero-relative-refinement}
    after passing to the underlying square-zero arrow.
\end{enumerate}
\end{proof}

\begin{remark}[Three distinct square-zero base-change operations]
\label{rem:three-square-zero-base-changes}
The preceding results distinguish three operations which will be used below.
Relative-base change changes the base \(A\) of a derivation; algebraic pullback
forms a limit \(E\times_DD'\) in commutative algebras; geometric Cartesian
base change forms a pushout \(B\otimes_EC\) of coordinate algebras. Their
variance is different, and no one of the three is used as a substitute for
another.
\end{remark}

\begin{remark}[Structured data and underlying arrows]
\label{rem:structured-square-zero-no-hidden-presentation}
The parameterization theorem concerns structured extensions. Later we pass to
the equivalence-closed class of their underlying arrows. At that point no
chosen coefficient or classifying derivation is retained. Recognition of a
structured presentation from an arbitrary underlying arrow is a separate
theorem on whatever bounded domain it is available; it is not part of the
definition of an infinitesimal substitution.
\end{remark}
The affine construction is functorial in the coordinate algebra, but a global
spectral thickening also requires compatibility with restriction to affine
opens of the centre. We now prove that compatibility and use it to glue the
affine square-zero pieces. The gluing data are supplied by the construction
itself rather than retained independently.

\subsection{Native global structured infinitesimal context extension}

\begin{lemma}[Zariski restriction of structured square-zero extensions]
\label{lem:square-zero-zariski-restriction}
Let \(A\to B\) be connective, let \(M\in(\Mod_B)_{\geq0}\), and let
\[
    \eta:L_{B/A}\longrightarrow M[1]
\]
classify the structured extension \(B_\eta\to B\). Let
\[
    u:B\longrightarrow B'
\]
represent an affine spectral Zariski open immersion and put
\[
    M':=B'\otimes_BM,
    \qquad
    \eta':L_{B'/A}\simeq B'\otimes_BL_{B/A}
        \longrightarrow M'[1],
\]
where the first equivalence uses transitivity and
\(L_{B'/B}\simeq0\). Functoriality of split square-zero formation and of the
defining pullback produces a canonical morphism
\[
    u_\eta:B_\eta\longrightarrow B'_{\eta'}.
\]
This morphism represents an affine spectral Zariski open immersion. Moreover
the square of spectral schemes
\[
\begin{tikzcd}[column sep=large,row sep=large]
    \Spec(B') \arrow[r] \arrow[d]
        & \Spec(B) \arrow[d]
    \\
    \Spec(B'_{\eta'}) \arrow[r]
        & \Spec(B_\eta)
\end{tikzcd}
\]
is Cartesian. The construction and the Cartesian comparison are coherent
under identities and composition of affine open restrictions.
\end{lemma}

\begin{proof}
\leavevmode
\begin{enumerate}[label=\textup{(\arabic*)}]
    \item \textbf{First suppose.}
    First suppose that \(B'=B[f^{-1}]\) is principal.

    \item \textbf{Verify surjectivity.}
    The surjection
    \[
        \pi_0(B_\eta)\twoheadrightarrow\pi_0(B)
    \]
    allows a lift \(\widetilde f\in\pi_0(B_\eta)\) of \(f\).

    \item \textbf{Localize the structured extension at.}
    Localize the structured extension at \(\widetilde f\):
    \[
        B_\eta[\widetilde f^{-1}]
        \longrightarrow
        B[f^{-1}].
    \]

    \item \textbf{Verify compatibility.}
    View the defining pullback square of \(B_\eta\) as a diagram of commutative
    \(B_\eta\)-algebras through its two projections to \(B\) and the common map to
    \(B\oplus M[1]\).

    \item \textbf{Apply base change.}
    Base change this diagram along the principal localization
    \[
        B_\eta\longrightarrow B_\eta[\widetilde f^{-1}].
    \]

    \item \textbf{Apply base change.}
    On underlying module spectra principal localization is exact, so the base-changed diagram
    remains Cartesian; limits of commutative algebras are created on underlying spectra.

    \item \textbf{Apply base change.}
    The two base changes of \(B\) are canonically \(B[f^{-1}]\).

    \item \textbf{Verify connectivity.}
    Indeed, because \(M\) is connective, \(M[1]\) is \(1\)-connective, so the
    derivation section and the zero section
    \[
        B\rightrightarrows B\oplus M[1]
    \]
    induce the same map on \(\pi_0\).

    \item \textbf{Conclude the step.}
    Hence the two projections \(B_\eta\rightrightarrows B\) agree on \(\pi_0\), and both images of \(\widetilde f\)
    are \(f\).

    \item \textbf{Use spectral localization.}
    Likewise split square-zero formation commutes with localization, so the third vertex is
    \[
        B[f^{-1}]\oplus M[f^{-1}][1].
    \]

    \item \textbf{Use spectral localization.}
    Cotangent localization identifies the nonzero section with the derivation \(\eta'\).

    \item \textbf{Apply base change.}
    Hence the base-changed Cartesian square is precisely
    \[
        B[f^{-1}]
        \times_{B[f^{-1}]\oplus M[f^{-1}][1]}
        B[f^{-1}],
    \]
    and therefore
    \[
        B_\eta[\widetilde f^{-1}]
        \simeq
        B'_{\eta'}.
    \]

    \item \textbf{Use spectral localization.}
    Thus \(u_\eta\) is a principal localization.

    \item \textbf{Derive the comparison.}
    Base changing the corresponding open immersion along the centre \(\Spec(B)\to\Spec(B_\eta)\) gives
    \(\Spec(B[f^{-1}])\), proving the Cartesian assertion in the principal case.

    \item \textbf{proof did not retain .}
    The proof did not retain \(\widetilde f\).

    \item \textbf{Verify invertibility.}
    Indeed two lifts differ by an element of the square-zero kernel of \(\pi_0(B_\eta)\to\pi_0(B)\); after either
    one is inverted, the other differs from it by a nilpotent perturbation and is therefore
    invertible as well.

    \item \textbf{Conclude the step.}
    Thus the principal open subobject and the resulting comparison are independent of the proof
    choice.

    \item \textbf{For a general affine open immersion.}
    For a general affine open immersion \(\Spec(B')\to\Spec(B)\), choose a finite principal cover of its
    image.

    \item \textbf{Use spectral localization.}
    On every principal member and on every finite intersection the preceding construction
    identifies the corresponding restriction of \(B'_\eta\) with a principal localization of
    \(B_\eta\).

    \item \textbf{Verify functoriality.}
    These identifications agree on refinements by functoriality of cotangent base change and
    structured square-zero formation.

    \item \textbf{Apply Zariski descent.}
    Spectral Zariski descent then shows that \(\Spec(B'_{\eta'})\to\Spec(B_\eta)\) is an open immersion and that the
    displayed square is Cartesian.

    \item \textbf{Use the stated property.}
    Since the canonical morphism \(u_\eta\) existed before choosing the principal cover, the
    cover is only a proof witness.

    \item \textbf{Identity and composition coherence.}
    Identity and composition coherence are inherited from the functorial square-zero
    construction and may be checked on the principal basis.
\end{enumerate}
\end{proof}

\begin{lemma}[Compatible-pair restriction of structured square-zero extensions]
\label{lem:square-zero-compatible-pair-restriction}
Let \((U,V)\) and \((U',V')\) be compatible affine pairs for a morphism of
spectral schemes, with
\[
    U=\Spec(B),\qquad V=\Spec(A),\qquad
    U'=\Spec(B'),\qquad V'=\Spec(A')
\]
and with \((U',V')\to(U,V)\) a refinement. Let
\(M\in(\Mod_B)_{\geq0}\) and let
\[
    \eta:L_{B/A}\longrightarrow M[1].
\]
Write
\[
    \bar B:=B\otimes_AA',
    \qquad
    \bar M:=M\otimes_B\bar B,
    \qquad
    M':=M\otimes_BB'.
\]
Cotangent base change gives
\[
    \bar\eta:
    L_{\bar B/A'}
    \simeq
    L_{B/A}\otimes_B\bar B
    \longrightarrow
    \bar M[1].
\]
The map \(\bar B\to B'\) represents an affine spectral Zariski open
immersion, and transitivity gives
\[
    L_{B'/A'}
    \simeq
    L_{\bar B/A'}\otimes_{\bar B}B'.
\]
Let
\[
    \eta':L_{B'/A'}\longrightarrow M'[1]
\]
be the resulting restriction. Then there is a canonical morphism
\[
    B_\eta\longrightarrow B'_{\eta'}
\]
which represents an affine spectral Zariski open immersion, and the square
\[
\begin{tikzcd}[column sep=large,row sep=large]
    \Spec(B') \arrow[r] \arrow[d]
        & \Spec(B) \arrow[d]
    \\
    \Spec(B'_{\eta'}) \arrow[r]
        & \Spec(B_\eta)
\end{tikzcd}
\]
is Cartesian. These morphisms and Cartesian comparisons are coherent under
identities, composition, and iterated compatible-pair refinement.
\end{lemma}

\begin{proof}
\leavevmode
\begin{enumerate}[label=\textup{(\arabic*)}]
    \item \textbf{Apply base change.}
    The affine open \(V'\to V\) base changes to an affine open
    \[
        \bar U:=U\times_VV'=\Spec(\bar B)
        \longrightarrow U.
    \]

    \item \textbf{Construct the factorization.}
    Since \(U'\to U\) lies over \(V'\), it factors canonically through an affine open
    immersion
    \[
        U'\longrightarrow\bar U,
    \]
    so \(\bar B\to B'\) represents an affine spectral Zariski open immersion.

    \item \textbf{Apply base change.}
    Structured square-zero base change gives a canonical equivalence
    \[
        B_\eta\otimes_AA'
        \xrightarrow{\simeq}
        \bar B_{\bar\eta}.
    \]

    \item \textbf{Apply base change.}
    Consequently
    \[
        \Spec(\bar B_{\bar\eta})
        \longrightarrow
        \Spec(B_\eta)
    \]
    is the base change of the affine open \(V'\to V\), and hence is an affine open immersion.

    \item \textbf{Apply the cited result, now over the fixed base .}
    Apply \Ref{lem:square-zero-zariski-restriction}, now over the fixed base \(A'\), to the
    affine open \(\bar B\to B'\).

    \item \textbf{Derive the comparison.}
    It gives an affine open immersion
    \[
        \Spec(B'_{\eta'})
        \longrightarrow
        \Spec(\bar B_{\bar\eta})
    \]
    whose pullback to the centre is \(\Spec(B')\to\Spec(\bar B)\).

    \item \textbf{Derive the comparison.}
    Composing the two open immersions gives the asserted map to \(\Spec(B_\eta)\).

    \item \textbf{centre square.}
    The centre square is Cartesian by pullback pasting.

    \item \textbf{on coordinate algebras.}
    Equivalently, on coordinate algebras,
    \[
    \begin{aligned}
        B\otimes_{B_\eta}B'_{\eta'}
        &\simeq
        \bigl(B\otimes_{B_\eta}\bar B_{\bar\eta}\bigr)
        \otimes_{\bar B_{\bar\eta}}B'_{\eta'}
        \\
        &\simeq
        \bar B\otimes_{\bar B_{\bar\eta}}B'_{\eta'}
        \\
        &\simeq
        B'.
    \end{aligned}
    \]

    \item \textbf{Apply cotangent base change.}
    Every comparison used here is the canonical cotangent-base-change,
    structured-square-zero-base-change, or Zariski-restriction comparison.

    \item \textbf{Their identity, composition, and higher-pasting coherences therefore.}
    Their identity, composition, and higher-pasting coherences therefore give the last
    assertion.
\end{enumerate}
\end{proof}

\begin{construction}[Global structured infinitesimal context extension]
\label{con:global-structured-infinitesimal-extension}
Let \(f:X\to S\) be a morphism of spectral schemes, let
\(M\in\QCoh(X)_{\geq0}\), and let
\[
    \eta:L_{X/S}\longrightarrow M[1].
\]
For every compatible affine pair
\[
    U=\Spec(B)\longrightarrow V=\Spec(A)
\]
for \(X\to S\), restriction gives
\[
    \eta_U:L_{B/A}\longrightarrow M_U[1]
\]
and hence the affine structured square-zero extension
\[
    B_{\eta_U}
    :=
    B\times_{B\oplus M_U[1]}B.
\]
For every compatible affine-pair refinement
\((U',V')\to(U,V)\), the canonical morphism supplied by
\Ref{lem:square-zero-compatible-pair-restriction} identifies
\[
    \Spec(B'_{\eta_{U'}})
    \longrightarrow
    \Spec(B_{\eta_U})
\]
as the corresponding affine open restriction. These open maps are coherent
under identities, composition, and further compatible refinement. Thus the constructed affine thickenings form a
spectral Zariski gluing diagram. Define
\[
    X_\eta
\]
to be the spectral scheme obtained by gluing this diagram, and let
\[
    i_\eta:X\longrightarrow X_\eta
\]
be the morphism obtained by gluing the affine closed immersions
\[
    \Spec(B)\longrightarrow\Spec(B_{\eta_U}).
\]
Neither the affine cover nor any principal refinement used to verify the open
transition maps is retained in \(X_\eta\).
\end{construction}

\begin{proposition}[Geometric properties of global structured infinitesimal extension]
\label{prop:global-square-zero-geometric-properties}
In \Ref{con:global-structured-infinitesimal-extension}:
\begin{enumerate}[label=(\roman*)]
    \item \(X_\eta\) is a spectral scheme;
    \item \(i_\eta:X\to X_\eta\) is a closed immersion over \(S\);
    \item the underlying topological map is a homeomorphism;
    \item pullback along \(i_\eta\) induces an equivalence between the small
    spectral Zariski sites of \(X_\eta\) and \(X\);
    \item on every compatible affine pair
    \(U=\Spec(B)\to V=\Spec(A)\), the restriction is
    \[
        \Spec(B)\longrightarrow\Spec(B_{\eta_U});
    \]
    \item the construction commutes with arbitrary represented relative-base
    change: for every Cartesian square
    \[
    \begin{tikzcd}[column sep=large,row sep=large]
        X' \arrow[r,"g'"] \arrow[d]
            & X \arrow[d]
        \\
        S' \arrow[r]
            & S,
    \end{tikzcd}
    \]
    if \(\eta'\) is the cotangent base change of \(\eta\), then there is a
    canonical equivalence over \(S'\)
    \[
        X'_{\eta'}\xrightarrow{\simeq}X_\eta\times_SS',
    \]
    compatible with the centres and coherent under iterated relative-base
    change.
\end{enumerate}
\end{proposition}

\begin{proof}
\leavevmode
\begin{enumerate}[label=\textup{(\arabic*)}]
    \item \textbf{Apply Zariski descent.}
    The compatible-pair open-transition theorem
    \Ref{lem:square-zero-compatible-pair-restriction} makes the affine pieces \(\Spec(B_{\eta_U})\) into a
    genuine Zariski gluing diagram, so spectral Zariski descent constructs a spectral scheme and
    proves \textup{(i)} and \textup{(v)}.

    \item \textbf{Verify surjectivity.}
    By \Ref{prop:structured-square-zero-fibre}, every affine coordinate map
    \[
        B_{\eta_U}\longrightarrow B
    \]
    is surjective on \(\pi_0\) with square-zero kernel.

    \item \textbf{Conclude the step.}
    Hence every local geometric map \(\Spec(B)\to\Spec(B_{\eta_U})\) is a closed immersion, and closedness is local
    on the target.

    \item \textbf{Identify the resulting comparison.}
    This proves \textup{(ii)}.

    \item \textbf{Verify nilpotence.}
    A quotient of an ordinary ring by a nilpotent ideal induces a homeomorphism on prime
    spectra.

    \item \textbf{Conclude the step.}
    Thus
    \[
        |\Spec(B)|\xrightarrow{\cong}|\Spec(B_{\eta_U})|
    \]
    on every affine piece.

    \item \textbf{Glue the local data.}
    The homeomorphisms commute with the open transition maps of
    \Ref{lem:square-zero-compatible-pair-restriction}, so they glue to \textup{(iii)}.

    \item \textbf{Use spectral localization.}
    Principal spectral opens are detected by degree-zero localization.

    \item \textbf{Across a square-zero quotient, every principal open of the centre.}
    Across a square-zero quotient, every principal open of the centre has the unique
    corresponding principal open of the thickening, and conversely because the prime spectra
    agree.

    \item \textbf{Lemma the cited result.}
    Lemma~\Ref{lem:square-zero-compatible-pair-restriction} shows that these correspondences are
    compatible with intersections and covering families.

    \item \textbf{Principal opens form bases of both small.}
    Principal opens form bases of both small spectral Zariski sites, so we obtain the
    equivalence in \textup{(iv)}.

    \item \textbf{For (vi), restrict the Cartesian square to.}
    For \textup{(vi)}, restrict the Cartesian square to a compatible affine pair
    \[
        U=\Spec(B)\longrightarrow V=\Spec(A)
    \]
    for \(X\to S\), and choose an affine open \(V'=\Spec(A')\to S'\) mapping to \(V\).

    \item \textbf{Fix the notation.}
    Put
    \[
        \bar B:=B\otimes_AA'.
    \]

    \item \textbf{Apply cotangent base change.}
    Cotangent base change transports \(\eta_U\) to a class
    \[
        \bar\eta:L_{\bar B/A'}\longrightarrow
        (M_U\otimes_B\bar B)[1],
    \]
    and \Ref{prop:structured-square-zero-base-change} gives the canonical equivalence
    \[
        B_{\eta_U}\otimes_AA'
        \xrightarrow{\simeq}
        \bar B_{\bar\eta}.
    \]

    \item \textbf{Now let be any compatible affine refinement.}
    Now let \(U'=\Spec(B')\) be any compatible affine refinement of the pulled-back source over
    \(V'\).

    \item \textbf{Then represents an affine spectral Zariski open.}
    Then \(\bar B\to B'\) represents an affine spectral Zariski open immersion.

    \item \textbf{Restrict the construction.}
    By \Ref{lem:square-zero-compatible-pair-restriction}, the corresponding restriction of
    \(\Spec(\bar B_{\bar\eta})\) is canonically \(\Spec(B'_{\eta'})\), and its pullback to the centre is \(\Spec(B')\).

    \item \textbf{Restrict the construction.}
    Hence, on the compatible affine basis of \(X'\), the restrictions of \(X_\eta\times_SS'\) are
    exactly the affine thickenings used to construct \(X'_{\eta'}\).

    \item \textbf{Apply Zariski descent.}
    The restriction comparisons commute with further affine refinement by the coherence of
    \Ref{lem:square-zero-compatible-pair-restriction}; spectral Zariski descent therefore glues
    them uniquely to
    \[
        X'_{\eta'}\xrightarrow{\simeq}X_\eta\times_SS'.
    \]

    \item \textbf{Apply cotangent base change.}
    Identity, composition, and higher-pasting coherence are inherited from affine cotangent base
    change, structured square-zero scalar extension, open restriction, and descent.
\end{enumerate}
\end{proof}

\begin{corollary}[Fixed-site ringed-topos recognition]
\label{cor:global-square-zero-fixed-site-recognition}
Through the site equivalence of
\Ref{prop:global-square-zero-geometric-properties}\textup{(iv)}, the
structured spectral-scheme recognition of \(X_\eta\) has the same underlying
small Zariski \(\infty\)-topos as \(X\). On that identified topos its
structure sheaf is the finite pullback
\[
    \mathcal O_X^\eta
    \simeq
    \mathcal O_X
    \times_{\mathcal O_X\oplus M[1]}
    \mathcal O_X,
\]
where the two maps are the sheaf-level derivation represented by \(\eta\) and
the zero derivation.
\end{corollary}

\begin{proof}
\leavevmode
\begin{enumerate}[label=\textup{(\arabic*)}]
    \item \textbf{Identify the equivalence.}
    The site equivalence is already a theorem of the native construction.

    \item \textbf{Transport the structure.}
    Under the structured comparison, restrict to any affine open of the common site.

    \item \textbf{structure sheaf of the constructed thickening.}
    The structure sheaf of the constructed thickening is represented there by \(B_{\eta_U}\), which
    is exactly the displayed pullback of affine structure sheaves.

    \item \textbf{Glue the local data.}
    These affine identifications commute with compatible affine restriction by
    \Ref{lem:square-zero-compatible-pair-restriction} and hence glue.

    \item \textbf{Conclude the step.}
    Thus the fixed-topos formula recognizes the native construction; it does not define it.
\end{enumerate}
\end{proof}

\begin{definition}[Structured infinitesimal extension of a spectral context]
\label{def:structured-infinitesimal-extension-global}
Let \(f:X\to S\) and \(M\in\QCoh(X)_{\geq0}\). A \textbf{structured
infinitesimal extension of \(X\) over \(S\) by \(M\)} is a derivation term
\[
    \eta:L_{X/S}\longrightarrow M[1]
\]
read together with the thickening \(X\to X_\eta\) functorially constructed in
\Ref{con:global-structured-infinitesimal-extension}. The thickening and its affine
pullback comparisons are outputs of the construction, not additional fields
of data. Write
\[
    \operatorname{SqZ}^{\mathrm{str}}_S(X;M)
\]
for the resulting \(\infty\)-groupoid.
\end{definition}

\begin{theorem}[Global structured infinitesimal-extension parameterization]
\label{thm:global-square-zero-classification}
For every \(f:X\to S\) and \(M\in\QCoh(X)_{\geq0}\), the construction above
gives a natural equivalence
\[
    \operatorname{SqZ}^{\mathrm{str}}_S(X;M)
    \simeq
    \Map_{\QCoh(X)}(L_{X/S},M[1]).
\]
It is compatible with represented relative-base change.
\end{theorem}

\begin{proof}
\leavevmode
\begin{enumerate}[label=\textup{(\arabic*)}]
    \item \textbf{right-hand side.}
    The right-hand side is, by definition, the space of global relative derivation terms.

    \item \textbf{Apply Zariski descent.}
    Restriction to compatible affine pairs and quasi-coherent descent identify it with the limit
    of the affine derivation spaces
    \[
        \Map_{\Mod_B}(L_{B/A},M_U[1]).
    \]

    \item \textbf{Invoke the cited result.}
    By \Ref{thm:structured-square-zero-classification}, each affine space is the
    \(\infty\)-groupoid of the corresponding constructed structured square-zero extension.

    \item \textbf{Glue the local data.}
    Lemma~\Ref{lem:square-zero-compatible-pair-restriction} identifies the restriction maps on
    these groupoids with the open transition maps used to glue \(X_\eta\).

    \item \textbf{Apply Zariski descent.}
    Thus the descent limit is precisely the groupoid of global structured extensions constructed
    above.

    \item \textbf{Apply base change.}
    Relative-base-change compatibility is
    \Ref{prop:global-square-zero-geometric-properties}\textup{(vi)} together with
    \Ref{thm:global-cotangent-base-change}.
\end{enumerate}
\end{proof}
The global structured construction therefore produces square-zero context
extensions with the required base-change behavior. The infinitesimal test
geometry, however, must depend only on the resulting arrows, not on chosen
coefficients or classifying derivations. We therefore pass from structured
outputs to an intrinsic class of underlying arrows.

\section{Underlying square-zero substitutions and finite infinitesimal contexts}
\label{sec:underlying-square-zero-finite}

The first step forgets the presentation of a constructed square-zero arrow
while retaining the property of belonging to its equivalence-closed essential
image. The second closes this class under finite composition. This produces
the presentation-free infinitesimal substitutions that will form the objects
of the infinitesimal site.

\subsection{Underlying square-zero substitutions}

\begin{definition}[Underlying spectral square-zero extension]
\label{def:underlying-spectral-square-zero-extension}
A morphism
\[
    q:A'\longrightarrow A
\]
of connective commutative ring spectra is a \textbf{spectral square-zero
extension} if, as an object of
\(\Fun(\Delta^1,\CAlg(\Sp)^{\mathrm{cn}})\), it is equivalent to an arrow
\[
    q_\eta:A_\eta\longrightarrow A
\]
constructed by \Ref{con:structured-square-zero-extension} from a connective
\(A\)-module \(M\) and an absolute degree-one derivation
\[
    \eta:L_A\longrightarrow M[1].
\]
No structured presentation is retained on \(q\).
\end{definition}

\begin{lemma}[Relative and absolute presentations have the same underlying arrow]
\label{lem:relative-absolute-square-zero-underlying}
Let \(A_0\to A\) be a morphism of commutative ring spectra, let
\(M\in\Mod_A\), and let
\[
    \eta:L_{A/A_0}\longrightarrow M[1]
\]
be a relative degree-one derivation. Put
\[
    \eta^{\mathrm{abs}}:
    L_A\longrightarrow L_{A/A_0}\xrightarrow{\eta}M[1].
\]
Then the structured square-zero extension classified by \(\eta\) over
\(A_0\) and the structured square-zero extension classified by
\(\eta^{\mathrm{abs}}\) over \(\mathbb S\) have canonically the same
underlying algebra arrow
\[
    A_\eta\longrightarrow A.
\]
Consequently the equivalence-closed essential image in
\Ref{def:underlying-spectral-square-zero-extension} is unchanged if relative
structured presentations are allowed in its definition.
\end{lemma}

\begin{proof}
\leavevmode
\begin{enumerate}[label=\textup{(\arabic*)}]
    \item \textbf{Transport the structure.}
    Under cotangent classification, precomposition \(L_A\to L_{A/A_0}\) is exactly the operation of
    forgetting the \(A_0\)-linearity of a relative derivation.

    \item \textbf{Conclude the step.}
    Hence the relative class \(\eta\) and the absolute class \(\eta^{\mathrm{abs}}\) determine the same
    underlying commutative-algebra section
    \[
        d_\eta:A\longrightarrow A\oplus M[1];
    \]
    the former additionally carries the specified compatibility with the \(A_0\)-structure.

    \item \textbf{Identify the map.}
    The zero section is likewise the same underlying map.

    \item \textbf{Compare the two constructions.}
    Both square-zero algebras are therefore the same functorial pullback
    \[
        A\times_{A\oplus M[1]}A,
    \]
    and their projections to \(A\) agree.

    \item \textbf{Only the relative base structure.}
    Only the relative base structure has been forgotten.
\end{enumerate}
\end{proof}

\begin{definition}[Underlying square-zero infinitesimal substitution]
\label{def:spectral-square-zero-substitution}
A morphism of affine spectral contexts
\[
    \Spec(A)\longrightarrow\Spec(A')
\]
is a \textbf{square-zero infinitesimal substitution} if its opposite coordinate
map \(A'\to A\) is a spectral square-zero extension. A morphism
\(j:U\to T\) of arbitrary spectral schemes is a square-zero infinitesimal
substitution if it is affine and becomes one of the preceding affine arrows
after pullback along every affine open of \(T\).
\end{definition}

\begin{proposition}[Permanence of square-zero infinitesimal substitutions]
\label{prop:square-zero-substitution-permanence}
Square-zero infinitesimal substitutions are stable under equivalence and
arbitrary base change. Every such substitution is a closed immersion, induces
a homeomorphism on underlying topological spaces, and induces an equivalence
of small spectral Zariski sites.
\end{proposition}

\begin{proof}
\leavevmode
\begin{enumerate}[label=\textup{(\arabic*)}]
    \item \textbf{Identify the equivalence.}
    The assertions are invariant under equivalence in the arrow category, so for an affine arrow
    choose, only inside the proof, a structured representative.

    \item \textbf{Apply base change.}
    Geometric Cartesian base change of that representative is again structured by
    \Ref{thm:structured-square-zero-geometric-base-change}; forgetting the transported
    coefficient and derivation shows that the underlying class is stable under arbitrary base
    change.

    \item \textbf{Verify connectivity.}
    The fibre and connectivity theorem for a structured representative gives surjectivity on
    \(\pi_0\) with square-zero kernel, hence closedness and the homeomorphism on prime
    spectra.

    \item \textbf{Restrict the construction.}
    The Zariski restriction theorem then identifies the small spectral Zariski sites.

    \item \textbf{All claims descend through equivalence of arrows, and.}
    All claims descend through equivalence of arrows, and the global statement follows from the
    affine-local definition.
\end{enumerate}
\end{proof}

\begin{remark}[No square-zero presentation datum]
\label{rem:no-square-zero-presentation-datum}
An underlying square-zero substitution is a property of an already specified
morphism. A coefficient module or classifying derivation may be chosen inside
a proof after replacing the arrow by an equivalent constructed arrow, but such
a choice is not retained on the morphism. This distinction is mathematically
necessary: \cite[Warning~3.9]{LurieDAGIV} notes that an underlying square-zero
extension need not determine its coefficient and classifying derivation
uniquely, even up to equivalence. Recognition theorems such as
\cite[Theorem~3.17]{LurieDAGIV} recover structured presentations only on their
stated bounded theorem domains; no such recognition is built into the present
definition.
\end{remark}

\begin{lemma}[Relative refinement of a chosen square-zero presentation]
\label{lem:underlying-square-zero-relative-refinement}
Let
\[
    A\longrightarrow E\xrightarrow{q}B
\]
be morphisms of connective commutative ring spectra and suppose that the
underlying arrow \(q\) is a spectral square-zero extension. Choose, only
inside the proof, an absolute structured representative supplied by
\Ref{def:underlying-spectral-square-zero-extension}, with connective
\(B\)-module \(M\) and degree-one derivation
\[
    \eta:L_B\longrightarrow M[1].
\]
Then the transported map \(A\to E\) canonically determines a lift of
\(\eta\) through transitivity to a relative derivation
\[
    \eta_A:L_{B/A}\longrightarrow M[1],
\]
and the structured square-zero extension over \(A\) classified by
\(\eta_A\) is equivalent to \(E\to B\) as an arrow of commutative
\(A\)-algebras. The relative derivation is a proof witness associated to the
chosen structured presentation and is not retained on the underlying arrow.
\end{lemma}

\begin{proof}
\leavevmode
\begin{enumerate}[label=\textup{(\arabic*)}]
    \item \textbf{Replace , only inside the proof, by the.}
    Replace \(q\), only inside the proof, by the chosen structured presentation and
    transport the given map \(A\to E\) across that equivalence.

    \item \textbf{Conclude the step.}
    Thus
    \[
        q:E\longrightarrow B
    \]
    is literally the pullback in \(\CAlg(\Sp)_{/B}\) of the two morphisms
    \[
        d_\eta,d_0:
        \id_B\rightrightarrows
        (B\oplus M[1]\to B).
    \]

    \item \textbf{Identify the map.}
    The composite \(A\to E\to B\) makes \(A\to B\) an object of \(\CAlg(\Sp)_{/B}\), and the transported
    map \(A\to E\) is a point of
    \[
        \Map_{\CAlg(\Sp)_{/B}}(A,E).
    \]

    \item \textbf{Compute the mapping object.}
    Mapping into the defining pullback identifies this mapping space with the space of paths
    between the two composites
    \[
        A\longrightarrow B
        \xrightarrow{d_\eta,d_0}
        B\oplus M[1]
    \]
    over \(B\).

    \item \textbf{Restrict the construction.}
    Hence the specified map \(A\to E\) determines canonically a nullhomotopy of the restriction
    of the absolute derivation.

    \item \textbf{Transport the structure.}
    Under cotangent classification this is a nullhomotopy of
    \[
        L_A\otimes_AB
        \longrightarrow
        L_B
        \xrightarrow{\eta}
        M[1].
    \]

    \item \textbf{Apply cotangent transitivity.}
    Apply \(\map_{\Mod_B}(-,M[1])\) to the transitivity cofibre sequence
    \[
        L_A\otimes_AB
        \longrightarrow
        L_B
        \longrightarrow
        L_{B/A}.
    \]

    \item \textbf{Compute the fibre.}
    The chosen nullhomotopy is precisely the datum of a lift of \(\eta\) through this fibre
    sequence, and therefore determines a relative class
    \[
        \eta_A:L_{B/A}\longrightarrow M[1].
    \]

    \item \textbf{Restrict the construction.}
    Its associated relative derivation has the same underlying absolute derivation \(d_\eta\)
    and the same specified homotopy to the zero derivation after restriction along \(A\).

    \item \textbf{Identify contractibility.}
    Applying the slice pullback construction of \Ref{con:structured-square-zero-extension}
    therefore recovers \(E\to B\), together with its given commutative \(A\)-algebra
    structure, through a contractible space of comparisons.

    \item \textbf{No coefficient or relative class.}
    No coefficient or relative class is retained after the proof.
\end{enumerate}
\end{proof}
The one-stage class is now intrinsic to the arrow. Finite infinitesimal
contexts are obtained by closing this class under finite composition, without
retaining a factorization as part of the resulting morphism.

\subsection{Finite infinitesimal substitutions}

\begin{definition}[Finite spectral infinitesimal substitution]
\label{def:finite-spectral-infinitesimal-substitution}
Let \(\operatorname{Inf}^{\mathrm{fin}}_{\Sp}\) be the smallest class of
morphisms of spectral schemes which
\begin{enumerate}[label=(\roman*)]
    \item contains all equivalences and all square-zero infinitesimal
    substitutions;
    \item is closed under finite composition.
\end{enumerate}
A morphism in this class is a \textbf{finite spectral infinitesimal
substitution}. No finite factorization is retained as structure on the
morphism.
\end{definition}

\begin{proposition}[Permanence of finite infinitesimal substitutions]
\label{prop:finite-infinitesimal-permanence}
Finite spectral infinitesimal substitutions are affine and are stable under
equivalence, finite composition, and arbitrary base change.
\end{proposition}

\begin{proof}
\leavevmode
\begin{enumerate}[label=\textup{(\arabic*)}]
    \item \textbf{Set up the argument.}
    Equivalences are affine, square-zero infinitesimal substitutions are affine by
    \Ref{def:spectral-square-zero-substitution}, and affine morphisms are closed under finite
    composition.

    \item \textbf{Conclude the step.}
    Hence every finite infinitesimal substitution is affine.

    \item \textbf{Stability under equivalence and finite composition.}
    Stability under equivalence and finite composition is built into
    \Ref{def:finite-spectral-infinitesimal-substitution}.

    \item \textbf{Apply base change.}
    Let \(\mathcal B\) be the class of morphisms whose arbitrary base changes belong to
    \(\operatorname{Inf}^{\mathrm{fin}}_{\Sp}\).

    \item \textbf{Establish the next comparison.}
    Equivalences belong to \(\mathcal B\), square-zero substitutions belong to \(\mathcal B\) by
    \Ref{prop:square-zero-substitution-permanence}, and \(\mathcal B\) is closed under finite
    composition.

    \item \textbf{Apply base change.}
    Minimality gives base-change stability.
\end{enumerate}
\end{proof}

\begin{proposition}[Zariski support invariance]
\label{prop:finite-infinitesimal-zariski-invariance}
Let
\[
    j:U\longrightarrow T
\]
be a finite spectral infinitesimal substitution. Then:
\begin{enumerate}[label=(\roman*)]
    \item the induced map \(|U|\to|T|\) is a homeomorphism;
    \item pullback along \(j\) induces an equivalence between the posets of
    open spectral subschemes of \(T\) and \(U\);
    \item pullback induces an equivalence of their small spectral Zariski
    sites, preserving finite intersections, covering families, and \v{C}ech
    nerves.
\end{enumerate}
\end{proposition}

\begin{proof}
\leavevmode
\begin{enumerate}[label=\textup{(\arabic*)}]
    \item \textbf{Each assertion holds for equivalences and for.}
    Each assertion holds for equivalences and for square-zero substitutions by
    \Ref{prop:square-zero-substitution-permanence}.

    \item \textbf{Identify the map.}
    The class of morphisms satisfying each assertion is closed under composition.

    \item \textbf{Apply the minimality in the cited result.}
    Apply the minimality in \Ref{def:finite-spectral-infinitesimal-substitution}.

    \item \textbf{Use preservation.}
    Compatibility with \v{C}ech nerves follows because the site equivalence is induced by
    pullback and therefore preserves finite fibre products.
\end{enumerate}
\end{proof}

\begin{remark}[Finite means generated by finite composition]
\label{rem:finite-infinitesimal-no-order}
The adjective \textbf{finite} refers only to generation by finite composition.
No numerical order, length, filtration, or chosen sequence of square-zero
stages is assigned to a finite infinitesimal substitution. A factorization
may be chosen inside an induction, but the conclusion depends only on the
underlying arrow.
\end{remark}
Finite composition gives the formal class of infinitesimal substitutions. We
next show that this terminology is reflected intrinsically on degree-zero
coordinates: every such arrow is a nilpotent thickening there.

\subsection{Nilpotence of finite infinitesimal substitutions}

\begin{lemma}[Nilpotence under composition of surjections]
\label{lem:nilpotence-composite-surjections}
Let
\[
    R\xrightarrow{a}S\xrightarrow{b}T
\]
be surjections of ordinary commutative rings. Put
\[
    I:=\ker(a),\qquad J:=\ker(b),\qquad K:=\ker(ba).
\]
If \(I^m=0\) and \(J^n=0\), then
\[
    K^n\subseteq I
\]
and hence
\[
    K^{mn}=0.
\]
\end{lemma}

\begin{proof}
\leavevmode
\begin{enumerate}[label=\textup{(\arabic*)}]
    \item \textbf{Apply the statement objectwise.}
    Every element of \(K\) maps into \(J\) under \(a\).

    \item \textbf{Conclude the step.}
    Therefore every product of \(n\) elements of \(K\) maps to zero in \(S\),
    because \(J^n=0\).

    \item \textbf{Conclude the step.}
    Hence \(K^n\subseteq I\).

    \item \textbf{Derive the comparison.}
    Raising to the \(m\)-th power gives \(K^{mn}\subseteq I^m=0\).
\end{enumerate}
\end{proof}

\begin{proposition}[Finite infinitesimal substitutions are nilpotent thickenings on degree zero]
\label{prop:finite-infinitesimal-pi0-nilpotent}
Let
\[
    j:U\longrightarrow T
\]
be a finite spectral infinitesimal substitution. Affine-locally on \(T\),
write
\[
    U=\Spec(B),\qquad T=\Spec(A).
\]
Then \(\pi_0(A)\to\pi_0(B)\) is surjective with nilpotent kernel. In
particular, the adjective \textbf{infinitesimal} is intrinsic to the underlying
arrow and does not depend on a chosen square-zero factorization.
\end{proposition}

\begin{proof}
\leavevmode
\begin{enumerate}[label=\textup{(\arabic*)}]
    \item \textbf{Construct the factorization.}
    Choose, only inside the proof, a finite factorization
    \[
        U=U_0\longrightarrow U_1\longrightarrow\cdots\longrightarrow U_r=T
    \]
    into underlying square-zero substitutions.

    \item \textbf{Apply the statement objectwise.}
    Every stage is an affine morphism.

    \item \textbf{Use the stated property.}
    Since \(U_r=T\) is affine, \(U_{r-1}\) is affine; repeating this argument backwards shows
    inductively that every intermediate \(U_k\) is affine.

    \item \textbf{Fix the notation.}
    Write
    \[
        U_k=\Spec(A_k).
    \]

    \item \textbf{Verify surjectivity.}
    For every stage the coordinate map
    \[
        A_{k+1}\longrightarrow A_k
    \]
    is surjective on \(\pi_0\) with square-zero kernel by
    \Ref{prop:square-zero-substitution-permanence}.

    \item \textbf{Verify nilpotence.}
    Iterating \Ref{lem:nilpotence-composite-surjections} therefore shows that the kernel of
    \[
        \pi_0(A_r)\longrightarrow\pi_0(A_0)
    \]
    is nilpotent.

    \item \textbf{Verify surjectivity.}
    Since nilpotence and surjectivity are properties of the composite ring map, the conclusion
    is independent of the chosen factorization.

    \item \textbf{global statement.}
    The global statement is local on the affine target.
\end{enumerate}
\end{proof}
The nilpotence result makes the infinitesimal character of the underlying
arrow independent of any chosen square-zero factorization. We can therefore
compare the discrete sector with the familiar finite nilpotent pattern of
ordinary algebraic geometry.

\subsection{Discrete comparison with the classical finite infinitesimal pattern}

\begin{lemma}[Recognition of discrete square-zero quotients]
\label{lem:discrete-square-zero-recognition}
Let
\[
    R'\twoheadrightarrow R
\]
be a surjection of ordinary commutative rings with square-zero kernel \(J\).
Then the Eilenberg--Mac Lane morphism
\[
    HR'\longrightarrow HR
\]
is a spectral square-zero extension in the sense of
\Ref{def:underlying-spectral-square-zero-extension}.
\end{lemma}

\begin{proof}
\leavevmode
\begin{enumerate}[label=\textup{(\arabic*)}]
    \item \textbf{Verify surjectivity.}
    The morphism is a \(0\)-small extension in the sense of \cite[Definition~3.12 and
    Remark~3.16]{LurieDAGIV}: both rings are discrete, the map on degree zero is surjective, and
    the multiplication on its fibre \(HJ\) is null because \(J^2=0\).

    \item \textbf{Identify the resulting object.}
    The square-zero recognition theorem \cite[Theorem~3.17]{LurieDAGIV} identifies every
    \(0\)-small extension with a square-zero extension classified by a degree-one
    derivation.

    \item \textbf{Conclude the step.}
    Hence \(HR'\to HR\) belongs to the equivalence-closed essential image used in
    \Ref{def:underlying-spectral-square-zero-extension}.

    \item \textbf{No classifying derivation.}
    No classifying derivation is retained on the underlying arrow.
\end{enumerate}
\end{proof}

\begin{proposition}[Nilpotent discrete thickenings are finite square-zero composites]
\label{prop:discrete-nilpotent-finite-square-zero}
Let \(R\) be an ordinary commutative ring and let \(I\subseteq R\) be a
nilpotent ideal with \(I^N=0\). The discrete spectral thickening
\[
    \Spec H(R/I)\longrightarrow\Spec HR
\]
is a finite spectral infinitesimal substitution. More precisely, the
coordinate quotient factors as
\[
    R=R/I^N\longrightarrow R/I^{N-1}\longrightarrow\cdots
    \longrightarrow R/I^2\longrightarrow R/I,
\]
and the kernel of
\[
    R/I^{k+1}\longrightarrow R/I^k
\]
is \(I^k/I^{k+1}\), which is square-zero.
\end{proposition}

\begin{proof}
\leavevmode
\begin{enumerate}[label=\textup{(\arabic*)}]
    \item \textbf{For , one.}
    For \(k\geq1\), one has \(2k\geq k+1\), hence
    \[
        I^{2k}\subseteq I^{k+1}.
    \]

    \item \textbf{Conclude the step.}
    Thus multiplication on \(I^k/I^{k+1}\) is zero.

    \item \textbf{Invoke the cited result.}
    By \Ref{lem:discrete-square-zero-recognition}, Eilenberg--Mac Lane formation carries each
    quotient
    \[
        R/I^{k+1}\longrightarrow R/I^k
    \]
    to an underlying spectral square-zero extension.

    \item \textbf{Derive the comparison.}
    Finite composition gives the asserted infinitesimal substitution.
\end{enumerate}
\end{proof}

\begin{theorem}[Discrete finite-infinitesimal comparison]
\label{thm:discrete-finite-infinitesimal-comparison}
Under the fully faithful discrete enhancement of
Chapter~\Ref{chap:spectral-algebraic-geometry}, finite spectral infinitesimal
substitutions between discrete affine contexts are exactly the ordinary
nilpotent closed thickenings. Under this identification, codomain-pullback
Zariski coverings agree with the usual Zariski coverings of nilpotent
thickenings. Consequently the discrete affine sector of the spectral finite
infinitesimal site has the classical finite infinitesimal-site pattern.
\end{theorem}

\begin{proof}
\leavevmode
\begin{enumerate}[label=\textup{(\arabic*)}]
    \item \textbf{One direction.}
    One direction is \Ref{prop:discrete-nilpotent-finite-square-zero}.

    \item \textbf{Verify nilpotence.}
    Conversely, \Ref{prop:finite-infinitesimal-pi0-nilpotent} makes the degree-zero kernel of a
    finite spectral infinitesimal substitution nilpotent.

    \item \textbf{Identify the resulting object.}
    For discrete affine contexts the fully faithful enhancement identifies the spectral
    coordinate map with the corresponding ordinary quotient map.

    \item \textbf{Use preservation.}
    Zariski opens are preserved and reflected by the discrete comparison of
    Chapter~\Ref{chap:spectral-algebraic-geometry}; codomain-pullback covers therefore agree
    with their classical counterparts.
\end{enumerate}
\end{proof}
We now have a presentation-free class of finite infinitesimal substitutions,
stable under base change and invisible on Zariski support. The next task is to
organize these arrows as contexts for a second indexed semantics and to put a
locality theory on them.

\section{The category of infinitesimal substitutions and its Zariski topology}
\label{sec:infinitesimal-arrow-category}

Finite infinitesimal substitutions form a full subcategory of the arrow
category. A morphism is therefore genuinely two-dimensional:
\[
\begin{tikzcd}[column sep=huge,row sep=large]
U \arrow[r,"\alpha"] \arrow[d,"j"']
    & U' \arrow[d,"j'"]
\\
T \arrow[r,"\beta"']
    & T'.
\end{tikzcd}
\qquad
\operatorname{dom}(j)=U,
\quad
\operatorname{cod}(j)=T.
\]
The identity-arrow embedding supplies ordinary affine contexts inside this
category, while domain and codomain give its two endpoints. We first record
that arrow-category structure and then transport the spectral Zariski topology
through the codomain endpoint.

\subsection{The infinitesimal arrow category}

\begin{definition}[Affine finite-infinitesimal category]
\label{def:affine-finite-infinitesimal-category}
Let
\[
    \Aff_{\Sp}^{\inf}
    \subseteq
    \Fun(\Delta^1,\Aff_{\Sp})
\]
be the full subcategory spanned by finite spectral infinitesimal substitutions
between affine spectral schemes. Define
\[
    \operatorname{dom},\operatorname{cod}:
    \Aff_{\Sp}^{\inf}\longrightarrow\Aff_{\Sp}
\]
by
\[
    \operatorname{dom}(U\to T)=U,
    \qquad
    \operatorname{cod}(U\to T)=T,
\]
and let
\[
    i:\Aff_{\Sp}\longrightarrow\Aff_{\Sp}^{\inf},
    \qquad
    i(T)=\id_T.
\]
\end{definition}

\begin{proposition}[Endpoint adjunctions]
\label{prop:infinitesimal-endpoint-adjunctions}
There is a canonical adjoint triple
\[
    \operatorname{cod}\dashv i\dashv\operatorname{dom}.
\]
\end{proposition}

\begin{proof}
\leavevmode
\begin{enumerate}[label=\textup{(\arabic*)}]
    \item \textbf{Fix the notation.}
    Let \(j:U\to T\) be a finite infinitesimal substitution and \(X\in\Aff_{\Sp}\).

    \item \textbf{Verify compatibility.}
    A morphism
    \[
        j\longrightarrow\id_X
    \]
    in the arrow category is a commutative square
    \[
    \begin{tikzcd}[column sep=large, row sep=large]
        U \arrow[r] \arrow[d,"j"']
            & X \arrow[d,equal]
        \\
        T \arrow[r]
            & X,
    \end{tikzcd}
    \]
    and is determined by its lower map.

    \item \textbf{Conclude the step.}
    Hence
    \[
        \Map_{\Aff_{\Sp}^{\inf}}(j,iX)
        \simeq
        \Map_{\Aff_{\Sp}}(\operatorname{cod}(j),X).
    \]

    \item \textbf{a morphism.}
    Similarly, a morphism \(\id_X\to j\) is determined by its upper map \(X\to U\), giving
    \[
        \Map_{\Aff_{\Sp}^{\inf}}(iX,j)
        \simeq
        \Map_{\Aff_{\Sp}}(X,\operatorname{dom}(j)).
    \]

    \item \textbf{Verify naturality.}
    Naturality gives the two adjunctions.
\end{enumerate}
\end{proof}

\begin{proposition}[Finite products of infinitesimal substitutions]
\label{prop:infinitesimal-category-finite-products}
The category \(\Aff_{\Sp}^{\inf}\) is closed under finite products in the
arrow category, and
\[
    i(T\times S)
    \simeq
    i(T)\times i(S).
\]
\end{proposition}

\begin{proof}
\leavevmode
\begin{enumerate}[label=\textup{(\arabic*)}]
    \item \textbf{Use terminality.}
    The terminal arrow is the identity of \(\Spec(\mathbb S)\).

    \item \textbf{For and , factor the product arrow.}
    For \(U\to T\) and \(V\to S\), factor the product arrow as
    \[
        U\times V
        \longrightarrow
        T\times V
        \longrightarrow
        T\times S.
    \]

    \item \textbf{Apply base change.}
    The first stage is a base change of \(U\to T\) and the second a base change of \(V\to S\).

    \item \textbf{Compare the two constructions.}
    Both are finite infinitesimal substitutions by \Ref{prop:finite-infinitesimal-permanence},
    hence so is their composite.
\end{enumerate}
\end{proof}
These endpoint functors determine the appropriate locality convention.
Because finite infinitesimal substitutions preserve the small Zariski site, we
cover the codomain and pull the cover back to the domain. This makes both
endpoints continuous and later controls the Kan extensions.

\subsection{The codomain-pullback spectral Zariski topology}

\begin{definition}[Codomain-pullback Zariski basis cover]
\label{def:target-pullback-zariski-cover}
Let \(j:U\to T\) be an object of \(\Aff_{\Sp}^{\inf}\). A
\textbf{codomain-pullback Zariski basis cover} of \(j\) is a finite family
\[
    \{j_\alpha:U_\alpha\to T_\alpha\}_{\alpha\in A}
    \longrightarrow
    j
\]
such that \(\{T_\alpha\to T\}_{\alpha\in A}\) is a finite affine spectral
Zariski covering family and every square
\[
\begin{tikzcd}[column sep=large, row sep=large]
    U_\alpha \arrow[r] \arrow[d,"j_\alpha"']
        & U \arrow[d,"j"]
    \\
    T_\alpha \arrow[r]
        & T
\end{tikzcd}
\]
is Cartesian. Thus
\[
    U_\alpha\simeq U\times_TT_\alpha.
\]
\end{definition}

\begin{proposition}[The infinitesimal Zariski pretopology]
\label{prop:infinitesimal-zariski-pretopology}
Codomain-pullback Zariski basis covers contain identities and are stable under
pullback and finite composition. They therefore generate a Grothendieck
topology
\[
    \tau^{\inf}_{\mathrm{Zar}}
\]
on \(\Aff_{\Sp}^{\inf}\).
\end{proposition}

\begin{proof}
\leavevmode
\begin{enumerate}[label=\textup{(\arabic*)}]
    \item \textbf{Identity covers.}
    Identity covers are immediate.

    \item \textbf{Pulling a codomain-pullback cover of back along.}
    Pulling a codomain-pullback cover of \(j:U\to T\) back along a morphism \(j':U'\to T'\to j\) in the
    arrow category pulls the codomain cover back to a spectral Zariski cover of \(T'\).

    \item \textbf{Its domain.}
    Its domain is
    \[
        U'\times_{T'}(T_\alpha\times_TT')
        \simeq
        U'\times_TT_\alpha,
    \]
    so the resulting squares remain Cartesian.

    \item \textbf{For composition, compose the codomain Zariski covers.}
    For composition, compose the codomain Zariski covers and use iterated pullback for the
    domains.

    \item \textbf{spectral Zariski pretopology of Chapter the cited result.}
    The spectral Zariski pretopology of Chapter~\Ref{chap:spectral-algebraic-geometry} is stable
    under both operations.
\end{enumerate}
\end{proof}

\begin{proposition}[Endpoint continuity]
\label{prop:infinitesimal-endpoint-continuity}
The functors \(\operatorname{dom},\operatorname{cod}\) carry codomain-pullback Zariski basis covers to spectral
Zariski covering families, and \(i\) carries finite affine spectral Zariski
covers to codomain-pullback covers.
\end{proposition}

\begin{proof}
\leavevmode
\begin{enumerate}[label=\textup{(\arabic*)}]
    \item \textbf{For this.}
    For \(\operatorname{cod}\) this is the definition.

    \item \textbf{Apply base change.}
    For \(\operatorname{dom}\), the domain family is the base change
    \[
        \{U\times_TT_\alpha\to U\}_\alpha
    \]
    of a spectral Zariski cover.

    \item \textbf{Derive the comparison.}
    For \(i\), a cover \(\{T_\alpha\to T\}\) gives the codomain-pullback family \(\{\id_{T_\alpha}\to\id_T\}\).
\end{enumerate}
\end{proof}

\begin{lemma}[Affine-open invariance under finite infinitesimal substitutions]
\label{lem:finite-infinitesimal-affine-open-invariance}
Let
\[
    j:U\longrightarrow T
\]
be a finite spectral infinitesimal substitution between affine spectral
schemes. Pullback along \(j\) induces an equivalence
\[
    j^*:
    \operatorname{AffOp}(T)
    \xrightarrow{\simeq}
    \operatorname{AffOp}(U)
\]
between the categories of affine spectral Zariski open subobjects. Under this
equivalence finite intersections and finite affine covering families are
preserved.
\end{lemma}

\begin{proof}
\leavevmode
\begin{enumerate}[label=\textup{(\arabic*)}]
    \item \textbf{First suppose.}
    First suppose that \(j\) is represented by a structured square-zero extension
    \[
        \Spec(B)\longrightarrow\Spec(B_\eta).
    \]

    \item \textbf{For an affine open , the cited result.}
    For an affine open \(\Spec(B')\to\Spec(B)\), \Ref{lem:square-zero-zariski-restriction} constructs an
    affine open
    \[
        \Spec(B'_{\eta'})\longrightarrow\Spec(B_\eta)
    \]
    whose pullback to \(\Spec(B)\) is \(\Spec(B')\).

    \item \textbf{Restrict the construction.}
    The construction is compatible with identities, composition, and further affine restriction.

    \item \textbf{Verify surjectivity.}
    Thus pullback is essentially surjective on affine opens for a structured square-zero stage.

    \item \textbf{Identify the resulting object.}
    It is fully faithful because \Ref{prop:square-zero-substitution-permanence} identifies the
    small spectral Zariski sites of the source and target.

    \item \textbf{Verify surjectivity.}
    In particular, it identifies the mapping posets of open subobjects, and hence the full
    subcategories of affine open subobjects once essential surjectivity has been established.

    \item \textbf{Now let be an underlying square-zero substitution.}
    Now let \(j\) be an underlying square-zero substitution between affine spectral
    schemes.

    \item \textbf{Its opposite affine coordinate.}
    Its opposite affine coordinate map is, by definition, an underlying spectral square-zero
    extension, hence is equivalent in the arrow category to a structured square-zero extension.

    \item \textbf{Replace the arrow, only inside this proof.}
    Replace the arrow, only inside this proof, by such an equivalent structured representative.

    \item \textbf{Affineness of an open subobject and the Cartesian pullback property.}
    Affineness of an open subobject and the Cartesian pullback property are invariant under
    equivalence in the arrow category.

    \item \textbf{Conclude the step.}
    Therefore the preceding one-stage argument applies to every underlying square-zero
    substitution between affine spectral schemes.

    \item \textbf{No structured presentation.}
    No structured presentation is retained.

    \item \textbf{Conclude the argument.}
    Finally let
    \[
        j:U\longrightarrow T
    \]
    be a finite infinitesimal substitution between affine spectral schemes.

    \item \textbf{Construct the factorization.}
    Choose, only inside the proof, a factorization
    \[
        U=U_0
        \longrightarrow
        U_1
        \longrightarrow
        \cdots
        \longrightarrow
        U_r=T
    \]
    into underlying square-zero substitutions.

    \item \textbf{Apply the statement objectwise.}
    Every stage
    \[
        U_i\longrightarrow U_{i+1}
    \]
    is affine by the definition of a square-zero infinitesimal substitution.

    \item \textbf{Use the stated property.}
    Since \(U_r=T\) is affine, affineness of the morphism \(U_{r-1}\to U_r\) implies that \(U_{r-1}\)
    is affine.

    \item \textbf{Repeating the same argument backwards.}
    Repeating the same argument backwards shows inductively that every intermediate \(U_i\)
    is affine.

    \item \textbf{Construct the factorization.}
    Hence the one-stage affine-open equivalence proved above may be iterated along the entire
    chosen factorization.

    \item \textbf{Starting from an affine open of , this.}
    Starting from an affine open of \(U=U_0\), this iteration produces an affine open of
    \(T=U_r\) whose pullback is the original affine open.

    \item \textbf{Verify surjectivity.}
    Thus
    \[
        j^*:
        \operatorname{AffOp}(T)
        \longrightarrow
        \operatorname{AffOp}(U)
    \]
    is essentially surjective.

    \item \textbf{Use full faithfulness.}
    Full faithfulness follows intrinsically from
    \Ref{prop:finite-infinitesimal-zariski-invariance}, which identifies the small spectral
    Zariski sites of \(T\) and \(U\).

    \item \textbf{Construct the factorization.}
    The same intrinsic site equivalence shows that the iterated construction is independent of
    the chosen finite square-zero factorization: it is the unique open subobject of \(T\)
    corresponding to the specified open subobject of \(U\).

    \item \textbf{Use the stated property.}
    Since the site equivalence is induced by pullback, it preserves finite intersections and
    covering families.

    \item \textbf{Conclude the step.}
    Therefore \(j^*\) is the asserted equivalence of affine-open categories with the stated
    compatibility.
\end{enumerate}
\end{proof}

\begin{proposition}[Domain classification of infinitesimal Zariski covers]
\label{prop:source-classification-infinitesimal-zariski}
For every \(j:U\to T\) in \(\Aff_{\Sp}^{\inf}\), domain evaluation induces
an equivalence between the category of codomain-pullback Zariski basis covers of
\(j\) and the category of finite affine spectral Zariski covers of \(U\),
including refinement morphisms and finite intersections.
\end{proposition}

\begin{proof}
\leavevmode
\begin{enumerate}[label=\textup{(\arabic*)}]
    \item \textbf{codomain-pullback cover.}
    A codomain-pullback cover is determined by its codomain cover.

    \item \textbf{Invoke the cited result.}
    By \Ref{lem:finite-infinitesimal-affine-open-invariance}, pullback along \(j:U\to T\)
    identifies finite affine Zariski covers of \(T\) with finite affine Zariski covers of
    \(U\).

    \item \textbf{Transport the same comparison.}
    The same lemma, together with the intrinsic site equivalence of
    \Ref{prop:finite-infinitesimal-zariski-invariance}, shows that refinement morphisms and
    finite intersections are carried to the corresponding domain refinements and intersections.

    \item \textbf{Identify the resulting comparison.}
    This is exactly the claimed equivalence of cover categories.
\end{enumerate}
\end{proof}
The codomain-pullback pretopology is now established and is subcanonical on the
resulting arrow category. Passing to sheaves produces the ambient
infinitesimal topos in which the endpoint semantics will live.

\subsection{The infinitesimal Zariski topos}

\begin{definition}[Finite-infinitesimal spectral Zariski topos]
\label{def:finite-infinitesimal-zariski-topos}
Define
\[
    \operatorname{Stk}^{\mathrm{sp}}_{\inf,\mathrm{Zar}}
    :=
    \Sh_{\tau^{\inf}_{\mathrm{Zar}}}
    (\Aff_{\Sp}^{\inf}).
\]
As in Chapter~\Ref{chap:spectral-algebraic-geometry}, no hypercompletion is
understood unless explicitly stated. From this point onward write
\[
    \mathcal H_{\Sp}
    :=
    \operatorname{Stk}^{\mathrm{sp}}_{\mathrm{Zar}},
    \qquad
    \mathcal H_{\inf}
    :=
    \operatorname{Stk}^{\mathrm{sp}}_{\inf,\mathrm{Zar}}.
\]
These symbols are fixed for the remainder of the chapter.
\end{definition}

\begin{proposition}[Subcanonicity of the infinitesimal Zariski topology]
\label{prop:infinitesimal-zariski-subcanonical}
The topology \(\tau^{\inf}_{\mathrm{Zar}}\) is subcanonical.
\end{proposition}

\begin{proof}
\leavevmode
\begin{enumerate}[label=\textup{(\arabic*)}]
    \item \textbf{Fix the notation.}
    Let \(k:V\to S\) be an object of \(\Aff_{\Sp}^{\inf}\).

    \item \textbf{Compute the mapping object.}
    For \(j:U\to T\), the mapping space in the arrow category is the pullback
    \[
        \Map_{\Aff_{\Sp}^{\inf}}(j,k)
        \simeq
        \Map_{\Aff_{\Sp}}(U,V)
        \times_{\Map_{\Aff_{\Sp}}(U,S)}
        \Map_{\Aff_{\Sp}}(T,S).
    \]

    \item \textbf{For a codomain-pullback cover of , codomain evaluation.}
    For a codomain-pullback cover of \(j\), codomain evaluation is a spectral Zariski
    cover of \(T\), and domain evaluation is the corresponding spectral Zariski cover of
    \(U\).

    \item \textbf{Apply Zariski descent.}
    Representable functors on \(\Aff_{\Sp}\) satisfy Zariski descent by
    Chapter~\Ref{chap:spectral-algebraic-geometry}.

    \item \textbf{Apply Zariski descent.}
    Limits commute with limits in spaces, so the displayed pullback of representable mapping
    functors satisfies the same descent condition.
\end{enumerate}
\end{proof}
The absolute infinitesimal topos supplies the global test geometry. For a
fixed relative context \(X\to S\), we also need the corresponding affine basis
of infinitesimal thickenings incident with \(X\) and lying over \(S\); its
covering families are inherited from the global topology rather than chosen
anew.

\subsection{The relative spectral infinitesimal site}

\begin{lemma}[Affineness across finite infinitesimal substitutions]
\label{lem:finite-infinitesimal-affineness-invariance}
Let
\[
    j:U\longrightarrow T
\]
be a finite spectral infinitesimal substitution of spectral schemes. Then
\[
    U\text{ is affine}
    \quad\Longleftrightarrow\quad
    T\text{ is affine}.
\]
\end{lemma}

\begin{proof}
\leavevmode
\begin{enumerate}[label=\textup{(\arabic*)}]
    \item \textbf{If is affine, then is affine because.}
    If \(T\) is affine, then \(U\) is affine because every finite infinitesimal
    substitution is affine.

    \item \textbf{Prove the converse.}
    Conversely suppose that \(U\) is affine.

    \item \textbf{Derive the comparison.}
    Applying the \(0\)-truncated reconstruction of
    Chapter~\Ref{chap:spectral-algebraic-geometry} gives a morphism of ordinary schemes
    \[
        t_0U\longrightarrow t_0T.
    \]

    \item \textbf{Verify nilpotence.}
    By \Ref{prop:finite-infinitesimal-pi0-nilpotent}, this morphism is Zariski-locally a
    surjective quotient by a nilpotent ideal; hence it is an ordinary nilpotent closed
    thickening.

    \item \textbf{Affineness of ordinary schemes.}
    Affineness of ordinary schemes is invariant under nilpotent thickenings; we use this
    standard theorem only at the \(0\)-truncated comparison interface (compare
    \cite{LurieSAG}).

    \item \textbf{Use the stated property.}
    Since \(t_0U\) is affine, \(t_0T\) is affine.

    \item \textbf{affineness-detection statement of the -truncated reconstruction theorem of.}
    The affineness-detection statement of the \(0\)-truncated reconstruction theorem of
    Chapter~\Ref{chap:spectral-algebraic-geometry} then implies that \(T\) is affine.

    \item \textbf{Identify the resulting comparison.}
    This is the one use in the construction of the affine finite-infinitesimal basis where the
    downstream \(0\)-truncated comparison interface of
    Chapter~\Ref{chap:spectral-algebraic-geometry} is invoked.

    \item \textbf{native definitions of finite infinitesimal substitutions and their.}
    The native definitions of finite infinitesimal substitutions and their Zariski topology are
    formulated intrinsically in spectral geometry.

    \item \textbf{Construct the factorization.}
    The intrinsic nilpotence property of \(j\) supplies the required input, while any
    chosen finite square-zero factorization remains a proof witness.
\end{enumerate}
\end{proof}

\begin{definition}[Relative finite spectral infinitesimal site]
\label{def:relative-spectral-infinitesimal-site}
Let \(X\to S\) be a morphism of spectral schemes. Define
\[
    \operatorname{Inf}^{\mathrm{fin}}(X/S)
\]
to be the category whose objects are diagrams
\[
\begin{tikzcd}[column sep=large,row sep=large]
    U \arrow[r,hook] \arrow[d,hook]
        & X \arrow[d]
    \\
    T \arrow[r]
        & S
\end{tikzcd}
\]
in which \(U\hookrightarrow X\) is an affine spectral Zariski open immersion
and \(U\hookrightarrow T\) is a finite spectral infinitesimal substitution.
A morphism is a commutative morphism of such diagrams over \(X\to S\).
A covering family of \((U\hookrightarrow T)\) is a finite family obtained from
a finite affine spectral Zariski cover \(\{T_\alpha\to T\}\) by Cartesian
pullback:
\[
    U_\alpha:=U\times_TT_\alpha.
\]
We call the resulting site the \textbf{relative finite spectral infinitesimal
site of \(X/S\)}.
\end{definition}

\begin{proposition}[The relative site is inherited from the global infinitesimal geometry]
\label{prop:relative-infinitesimal-site-inherited}
The covering families of
\(\operatorname{Inf}^{\mathrm{fin}}(X/S)\) are precisely the restrictions of
the codomain-pullback Zariski covering families of
\(\Aff_{\Sp}^{\inf}\). They contain identities and are stable under
pullback and composition inside the relative category
\(\operatorname{Inf}^{\mathrm{fin}}(X/S)\). No independent infinitesimal
topology is chosen.
\end{proposition}

\begin{proof}
\leavevmode
\begin{enumerate}[label=\textup{(\arabic*)}]
    \item \textbf{For every object of the relative site.}
    For every object of the relative site, \(U\) is affine by definition and \(U\to T\)
    is finite infinitesimal.

    \item \textbf{Conclude the step.}
    Hence \(T\) is affine by \Ref{lem:finite-infinitesimal-affineness-invariance}.

    \item \textbf{Conclude the step.}
    Thus every object has an underlying object \(U\to T\) of \(\Aff_{\Sp}^{\inf}\), and its covering
    definition is exactly \Ref{def:target-pullback-zariski-cover}.

    \item \textbf{Fix the notation.}
    Let \((U'\to T')\to(U\to T)\) be a morphism in the relative category and let \(\{T_\alpha\to T\}\) be a covering
    family.

    \item \textbf{Restrict the construction.}
    The pullbacks
    \[
        T'_\alpha:=T'\times_TT_\alpha
    \]
    form a finite affine spectral Zariski cover of the affine \(T'\), and
    \[
        U'_\alpha
        :=U'\times_{T'}T'_\alpha
    \]
    is the corresponding affine open restriction of \(U'\).

    \item \textbf{Apply base change.}
    Base-change stability of finite infinitesimal substitutions makes \(U'_\alpha\to T'_\alpha\) an object of
    the same relative category.

    \item \textbf{Identify the resulting comparison.}
    This is exactly pullback stability of the global codomain-pullback pretopology; composition
    is inherited in the same way.

    \item \textbf{Identify the map.}
    The maps to \(X\to S\) are incidence data and do not alter the covering condition.
\end{enumerate}
\end{proof}

\begin{remark}[External base change and the affine basis]
\label{rem:relative-infinitesimal-site-external-base-change}
External base change for an arbitrary morphism of relative contexts
\((X'/S')\to(X/S)\) is handled on the affine basis by refining the pullback
of each affine open of \(X\) by affine spectral Zariski opens of \(X'\).
The construction is then performed on that affine refinement, as supplied by
ordinary Zariski locality.
\end{remark}

\begin{remark}[Classical infinitesimal-site comparison]
\label{rem:relative-infinitesimal-site-classical}
On discrete spectral contexts, \Ref{thm:discrete-finite-infinitesimal-comparison}
identifies the objects with nilpotent thickenings of affine open subcontexts
and identifies the covering families with their ordinary Zariski coverings.
Thus the construction is the spectral analogue of the finite infinitesimal
site pattern of \cite{Ogus1975InfinitesimalSite}; no ordinary scheme category
is used in its native definition.
\end{remark}
These constructions supply the infinitesimal test geometry: its arrow
category, Zariski topology, sheaf topos, and relative affine-basis sites. The
endpoint semantics can therefore be derived from this geometry rather than
postulated separately.

\section{Differential cohesion and the infinitesimal-shape calculus}
\label{sec:endpoint-semantics}

The identity-arrow embedding has both domain and codomain adjoints at the level
of affine infinitesimal contexts. This section promotes those endpoint
relations to an adjoint quadruple on sheaf topoi. Only after the quadruple has
been constructed do we recognize it as differential cohesion and package its
composites as
\[
    \Re\dashv\Im\dashv\&.
\]
The resulting infinitesimal-shape reflector then supplies the formal lifting
comparison and the infinitesimal disk relation used in the next section.

\subsection{Kan-extension formulas}

Let
\[
    i^*:
    \PSh(\Aff_{\Sp}^{\inf})
    \longrightarrow
    \PSh(\Aff_{\Sp})
\]
denote restriction along \(i^{\op}\). The endpoint adjunctions compute its
left and right Kan extensions pointwise, which gives the first three sheaf
adjoints.

\begin{proposition}[Endpoint Kan extensions]
\label{prop:endpoint-kan-extensions}
For every presheaf \(F\) on \(\Aff_{\Sp}\), there are canonical equivalences
\[
    \operatorname{Lan}_{i^{\op}}F
    \simeq
    \operatorname{cod}^*F,
    \qquad
    \operatorname{Ran}_{i^{\op}}F
    \simeq
    \operatorname{dom}^*F.
\]
Pointwise,
\[
    (\operatorname{cod}^*F)(U\to T)=F(T),
    \qquad
    (\operatorname{dom}^*F)(U\to T)=F(U).
\]
\end{proposition}

\begin{proof}
\leavevmode
\begin{enumerate}[label=\textup{(\arabic*)}]
    \item \textbf{Use the adjunction.}
    The adjunctions \(\operatorname{cod}\dashv i\dashv\operatorname{dom}\) of \Ref{prop:infinitesimal-endpoint-adjunctions} become, after
    passage to opposite categories, the adjunctions computing left and right Kan extension along
    \(i^{\op}\).

    \item \textbf{Use the adjunction.}
    A Kan extension along a functor possessing the relevant adjoint is computed by
    precomposition with that adjoint.
\end{enumerate}
\end{proof}

\begin{proposition}[The first three sheaf adjoints]
\label{prop:first-three-infinitesimal-sheaf-adjoints}
The endpoint formulas restrict to sheaves and give adjunctions
\[
    i_!
    \dashv
    i^*
    \dashv
    i_*
    :
    \operatorname{Stk}^{\mathrm{sp}}_{\mathrm{Zar}}
    \rightleftarrows
    \operatorname{Stk}^{\mathrm{sp}}_{\inf,\mathrm{Zar}},
\]
with
\[
    i_!=\operatorname{cod}^*,
    \qquad
    i_*=\operatorname{dom}^*.
\]
Both \(i_!\) and \(i_*\) are fully faithful, and \(i_!\) preserves all
limits.
\end{proposition}

\begin{proof}
\leavevmode
\begin{enumerate}[label=\textup{(\arabic*)}]
    \item \textbf{Apply Zariski descent.}
    On a codomain-pullback basis cover, the descent diagram for \(\operatorname{cod}^*F\) is exactly the
    Zariski descent diagram for the codomain cover.

    \item \textbf{descent diagram for.}
    The descent diagram for \(\operatorname{dom}^*F\) is the one for the domain cover, which is Zariski by
    \Ref{prop:infinitesimal-endpoint-continuity}.

    \item \textbf{Conclude the step.}
    Thus both precomposition functors carry Zariski sheaves to infinitesimal Zariski sheaves.

    \item \textbf{Restrict the construction.}
    Conversely, restriction along \(i\) carries codomain-pullback covers of identity
    arrows to ordinary Zariski covers, so \(i^*\) preserves sheaves.

    \item \textbf{Use the adjunction.}
    The presheaf adjunctions therefore restrict to sheaves.

    \item \textbf{Use the stated property.}
    Since
    \[
        \operatorname{cod}\,i\simeq\id_{\Aff_{\Sp}}
        \simeq
        \operatorname{dom}\,i,
    \]
    we have
    \[
        i^*i_!\simeq\id,
        \qquad
        i^*i_*\simeq\id,
    \]
    so \(i_!\) and \(i_*\) are fully faithful.

    \item \textbf{Conclude the argument.}
    Finally \(i_!=\operatorname{cod}^*\) is precomposition and therefore preserves all limits.
\end{enumerate}
\end{proof}
Three endpoint adjoints have been obtained. The fourth exists once the domain
embedding \(i_*\) is shown to preserve sheaf colimits, so the remaining task is
to compare domain pullback with Zariski sheafification.

\subsection{Domain sheafification and the fourth adjoint}

\begin{lemma}[Domain pullback preserves Zariski local equivalences]
\label{lem:source-pullback-local-equivalences}
Let
\[
    u:P\longrightarrow Q
\]
be a spectral Zariski local equivalence of presheaves on \(\Aff_{\Sp}\).
Then
\[
    \operatorname{dom}^*u:\operatorname{dom}^*P\longrightarrow \operatorname{dom}^*Q
\]
is a \(\tau^{\inf}_{\mathrm{Zar}}\)-local equivalence.
\end{lemma}

\begin{proof}
\leavevmode
\begin{enumerate}[label=\textup{(\arabic*)}]
    \item \textbf{Identify the equivalence.}
    The local-equivalence class of spectral Zariski sheafification is the strongly saturated
    class generated by covering-sieve monomorphisms \(R\hookrightarrow h_V\).

    \item \textbf{Fix the notation.}
    Let \(\mathcal K\) be the class of morphisms whose domain pullback is infinitesimally Zariski
    local.

    \item \textbf{Use preservation.}
    Precomposition with \(\operatorname{dom}\) preserves colimits, equivalences, compositions, and retract
    diagrams, while the infinitesimal local-equivalence class is strongly saturated.

    \item \textbf{It is enough to treat a covering-sieve.}
    It is enough to treat a covering-sieve generator.

    \item \textbf{Fix , an object , and a.}
    Fix \(R\hookrightarrow h_V\), an object \(j:U\to T\), and a section
    \[
        x\in(\operatorname{dom}^*h_V)(j)=\Map(U,V).
    \]

    \item \textbf{Pull back the covering sieve along .}
    Pull back the covering sieve along \(x\).

    \item \textbf{It contains a finite affine Zariski basis.}
    It contains a finite affine Zariski basis cover \(\{U_\alpha\to U\}\) by the
    finite-principal-refinement theorem of Chapter~\Ref{chap:spectral-algebraic-geometry}.

    \item \textbf{Invoke the cited result.}
    By \Ref{prop:source-classification-infinitesimal-zariski}, this domain cover is the domain
    family of a codomain-pullback basis cover
    \[
        \{j_\alpha:U_\alpha\to T_\alpha\}\longrightarrow j.
    \]

    \item \textbf{On each , the restricted section.}
    On each \(j_\alpha\), the restricted section lies in the pulled-back sieve.

    \item \textbf{Conclude the step.}
    Thus every local section of \(\operatorname{dom}^*h_V\) is infinitesimally Zariski-locally in the image of
    \(\operatorname{dom}^*R\).

    \item \textbf{Establish the next comparison.}
    Sheafification makes the resulting monomorphism both a monomorphism and an effective
    epimorphism, hence an equivalence.

    \item \textbf{Strong saturation.}
    Strong saturation proves the lemma.
\end{enumerate}
\end{proof}

\begin{theorem}[Domain pullback commutes with sheafification]
\label{thm:source-pullback-sheafification}
Let
\[
    a_{\mathrm{Zar}}^{\mathrm{sp}}
    :
    \PSh(\Aff_{\Sp})
    \longrightarrow
    \operatorname{Stk}^{\mathrm{sp}}_{\mathrm{Zar}}
\]
and
\[
    a_{\inf,\mathrm{Zar}}^{\mathrm{sp}}
    :
    \PSh(\Aff_{\Sp}^{\inf})
    \longrightarrow
    \operatorname{Stk}^{\mathrm{sp}}_{\inf,\mathrm{Zar}}
\]
be the two sheafification functors. There is a canonical equivalence
\[
    a_{\inf,\mathrm{Zar}}^{\mathrm{sp}}\operatorname{dom}^*
    \simeq
    \operatorname{dom}^*a_{\mathrm{Zar}}^{\mathrm{sp}}.
\]
Consequently \(i_*=\operatorname{dom}^*\) preserves all small colimits.
\end{theorem}

\begin{proof}
\leavevmode
\begin{enumerate}[label=\textup{(\arabic*)}]
    \item \textbf{For a presheaf , its sheafification unit.}
    For a presheaf \(P\), its sheafification unit \(P\to a_{\mathrm{Zar}}^{\mathrm{sp}}P\) is a Zariski local
    equivalence.

    \item \textbf{Invoke the cited result.}
    By \Ref{lem:source-pullback-local-equivalences}, its domain pullback is an infinitesimal
    Zariski local equivalence.

    \item \textbf{target is already a sheaf by the.}
    The target \(\operatorname{dom}^*a_{\mathrm{Zar}}^{\mathrm{sp}}P\) is already a sheaf by
    \Ref{prop:first-three-infinitesimal-sheaf-adjoints}; hence it realizes the infinitesimal
    sheafification of \(\operatorname{dom}^*P\).

    \item \textbf{Colimits of sheaves.}
    Colimits of sheaves are obtained by sheafifying presheaf colimits.

    \item \textbf{Use the stated property.}
    Since precomposition preserves presheaf colimits and domain pullback commutes with
    sheafification, \(\operatorname{dom}^*\) preserves sheaf colimits.
\end{enumerate}
\end{proof}

\begin{theorem}[The infinitesimal adjoint string]
\label{thm:infinitesimal-adjoint-string}
The functor \(i_*\) admits a right adjoint \(i^!\), giving an adjoint string
\[
    i_!
    \dashv
    i^*
    \dashv
    i_*
    \dashv
    i^!
    :
    \operatorname{Stk}^{\mathrm{sp}}_{\mathrm{Zar}}
    \rightleftarrows
    \operatorname{Stk}^{\mathrm{sp}}_{\inf,\mathrm{Zar}}.
\]
For \(E\in\operatorname{Stk}^{\mathrm{sp}}_{\inf,\mathrm{Zar}}\) and an
affine spectral context \(T\),
\[
    (i^!E)(T)
    \simeq
    \Map_{\operatorname{Stk}^{\mathrm{sp}}_{\inf,\mathrm{Zar}}}
    (i_*h_T,E).
\]
\end{theorem}

\begin{proof}
\leavevmode
\begin{enumerate}[label=\textup{(\arabic*)}]
    \item \textbf{Use the adjunction.}
    The first three adjoints are \Ref{prop:first-three-infinitesimal-sheaf-adjoints}.

    \item \textbf{Use preservation.}
    The domain embedding \(i_*\) preserves all small colimits by
    \Ref{thm:source-pullback-sheafification}; the two sheaf categories are presentable.

    \item \textbf{Use the adjunction.}
    The presentable adjoint functor theorem therefore supplies a right adjoint \(i^!\).

    \item \textbf{Apply Yoneda.}
    The pointwise formula follows from Yoneda and the adjunction \(i_*\dashv i^!\).
\end{enumerate}
\end{proof}
The endpoint quadruple has now been constructed. Only at this point do we
introduce the differential-cohesive terminology and package its composites as
the three standard modalities.

\subsection{Differential-cohesive recognition and modalities}

\begin{theorem}[Differential-cohesive recognition]
\label{thm:differential-cohesive-recognition}
The canonically constructed adjoint quadruple
\[
    i_!\dashv i^*\dashv i_*\dashv i^!
    :
    \mathcal H_{\Sp}
    \rightleftarrows
    \mathcal H_{\inf},
\]
with \(i_!\) and \(i_*\) fully faithful and \(i_!\) left exact, is the
\(\infty\)-categorical adjoint configuration called \textbf{differential
cohesion}; compare \cite{KhavkineSchreiber2017}. This is a
recognition theorem for the endpoint semantics already constructed above and
introduces no additional cohesion datum.
\end{theorem}

\begin{proof}
\leavevmode
\begin{enumerate}[label=\textup{(\arabic*)}]
    \item \textbf{Use full faithfulness.}
    The adjoint quadruple and full faithfulness are
    \Ref{prop:first-three-infinitesimal-sheaf-adjoints} and
    \Ref{thm:infinitesimal-adjoint-string}.

    \item \textbf{Use the stated property.}
    Since
    \[
        i_!=\operatorname{cod}^*,
    \]
    it is precomposition and therefore preserves all limits, in particular finite limits.

    \item \textbf{Use the adjunction.}
    These are the structural properties used in the cited differential-cohesive adjoint pattern.

    \item \textbf{Use the adjunction.}
    Every functor and adjunction has already been constructed from the infinitesimal arrow
    category and its Zariski topology, so the terminology adds no data.
\end{enumerate}
\end{proof}

\begin{remark}[Qualification of the cohesion terminology]
\label{rem:differential-cohesion-qualified}
We use \textbf{differential cohesion} only for the stated
\(\infty\)-categorical adjoint configuration constructed in this chapter, in
the sense of the corresponding endpoint-adjunction definition in
\cite[Definition~2.15]{KhavkineSchreiber2017}. The terminology imports no
additional synthetic differential-geometric axiom. In particular, no
Kock--Lawvere principle, microlinearity hypothesis, specified infinitesimal
object, formal-manifold condition, local-surjectivity axiom, or unqualified
Lawvere--Menni cohesion over spaces is assumed unless it is separately
constructed or stated as a theorem hypothesis. We also avoid the phrase
\textbf{infinitesimally cohesive}, which has a distinct use in derived
deformation theory.
\end{remark}

\begin{definition}[Differential-cohesive modalities]
\label{def:differential-cohesive-endpoint-constructions}
Define endofunctors of \(\mathcal H_{\inf}\) by the standard notation
\[
    \Re:=i_!i^*,
    \qquad
    \Im:=i_*i^*,
    \qquad
    \&:=i_*i^!.
\]
They are called, respectively, the \textbf{reduction modality}, the
\textbf{infinitesimal shape modality}, and the \textbf{infinitesimal flat
modality}. An object is \textbf{reduced} when the counit
\(\Re E\to E\) is an equivalence and \textbf{coreduced} when the unit
\(E\to\Im E\) is an equivalence.
\end{definition}

\begin{remark}[Reduction is relative to the infinitesimal endpoint]
\label{rem:differential-cohesive-reduction-meaning}
Here \textbf{reduction} means the differential-cohesive operation determined
by the endpoint inclusion \(i\): it removes the finite-infinitesimal direction
encoded by an object of \(\mathcal H_{\inf}\). This operation is distinct
from scheme-theoretic reduction \(X\mapsto X_{\mathrm{red}}\) and from the
\(0\)-truncation \(t_0\) of
Chapter~\Ref{chap:spectral-algebraic-geometry}. For a representable finite
infinitesimal arrow \(j:U\to T\), the calculation below gives
\[
    \Re(h_j)\simeq h_{\id_U}.
\]
\end{remark}

\begin{theorem}[Differential-cohesive modal triple]
\label{thm:endpoint-cohesive-triple}
There is a canonical adjoint triple
\[
    \Re\dashv\Im\dashv\&.
\]
The functors \(\Re\) and \(\&\) are idempotent comonads and \(\Im\) is an
idempotent monad. The modality \(\Re\) preserves finite products, while
\(\Im\) is a left-exact reflector onto the essential image of \(i_*\). There
are canonical transformations
\[
    \Re E\longrightarrow E\longrightarrow\Im E.
\]
For a finite infinitesimal context \(j:U\to T\),
\[
    (\Re E)(j)\simeq E(\id_T),
    \qquad
    (\Im E)(j)\simeq E(\id_U),
\]
and
\[
    (\&E)(j)
    \simeq
    \Map_{\mathcal H_{\inf}}(i_*h_U,E).
\]
The composite \(\Re E\to\Im E\) evaluates at \(j\) as restriction along the
infinitesimal substitution
\[
    E(\id_T)\longrightarrow E(\id_U).
\]
\end{theorem}

\begin{proof}
\leavevmode
\begin{enumerate}[label=\textup{(\arabic*)}]
    \item \textbf{Use the adjunction.}
    For \(E,F\in\mathcal H_{\inf}\), the adjoint quadruple gives
    \[
    \begin{aligned}
        \Map(\Re E,F)
        &=\Map(i_!i^*E,F)
        \\
        &\simeq\Map(i^*E,i^*F)
        \\
        &\simeq\Map(E,i_*i^*F)
        =\Map(E,\Im F),
    \end{aligned}
    \]
    so \(\Re\dashv\Im\).

    \item \textbf{Establish the next comparison.}
    Similarly \(\Im\dashv\&\).

    \item \textbf{Use full faithfulness.}
    Full faithfulness of \(i_!\) and \(i_*\) gives the corresponding idempotence.

    \item \textbf{Use preservation.}
    The functor \(\Re=i_!i^*\) preserves finite products because both factors do; indeed
    \(i_!\) preserves all limits here.

    \item \textbf{Use the adjunction.}
    The unit \(E\to i_*i^*E\) exhibits the essential image of the fully faithful right adjoint
    \(i_*\) as reflective, and both \(i^*\) and \(i_*\) preserve limits, so
    \(\Im\) is left exact.

    \item \textbf{pointwise formulas.}
    The pointwise formulas are the endpoint Kan-extension formulas of
    \Ref{prop:endpoint-kan-extensions}.

    \item \textbf{last assertion.}
    The last assertion is evaluation of the canonical composite on
    \[
        \id_U\longrightarrow j\longrightarrow\id_T.
    \]
\end{enumerate}
\end{proof}

\begin{remark}[The modalities are derived, not imposed]
\label{rem:endpoint-cohesive-no-axiom}
The three modalities are composites of the constructed endpoint adjoints.
Their adjunctions, idempotence, and natural transformations are consequences of
that quadruple; none is input structure on an infinitesimal sheaf.
\end{remark}

\begin{proposition}[Reduction of a representable infinitesimal context]
\label{prop:reduction-representable-infinitesimal-context}
Let \(j:U\to T\) be an object of \(\Aff_{\Sp}^{\inf}\), and let \(h_j\) be
its representable infinitesimal sheaf. Then
\[
    i^*h_j\simeq h_U
\]
and consequently
\[
    \Re(h_j)\simeq h_{\id_U}.
\]
Thus reduction sends a representable infinitesimal context to its
source/reduction.
\end{proposition}

\begin{proof}
\leavevmode
\begin{enumerate}[label=\textup{(\arabic*)}]
    \item \textbf{Apply Yoneda.}
    For an affine spectral context \(V\), Yoneda and the endpoint adjunction give
    \[
    \begin{aligned}
        (i^*h_j)(V)
        &=\Map_{\Aff_{\Sp}^{\inf}}(\id_V,j)
        \\
        &\simeq\Map_{\Aff_{\Sp}}(V,U)
        =h_U(V).
    \end{aligned}
    \]

    \item \textbf{Conclude the step.}
    Hence \(i^*h_j\simeq h_U\).

    \item \textbf{Derive the comparison.}
    Applying \(i_!\) gives the second equivalence.
\end{enumerate}
\end{proof}

\begin{proposition}[Infinitesimal shape of an ordinary spectral stack]
\label{prop:infinitesimal-shape-forgets-variation}
For an ordinary spectral Zariski stack \(X\), regarded in infinitesimal
semantics through \(i_!X\), there is a canonical equivalence
\[
    \Im(i_!X)\simeq i_*X.
\]
At a finite infinitesimal context \(j:U\to T\),
\[
    (i_!X)(j)\simeq X(T),
    \qquad
    (\Im i_!X)(j)\simeq X(U),
\]
and the unit
\[
    i_!X\longrightarrow\Im(i_!X)
\]
evaluates as restriction
\[
    X(T)\longrightarrow X(U).
\]
\end{proposition}

\begin{proof}
\leavevmode
\begin{enumerate}[label=\textup{(\arabic*)}]
    \item \textbf{Use the stated property.}
    Since \(i^*i_!\simeq\id\), one has
    \[
        \Im(i_!X)=i_*i^*i_!X\simeq i_*X.
    \]

    \item \textbf{pointwise formulas.}
    The pointwise formulas are \Ref{prop:endpoint-kan-extensions}.
\end{enumerate}
\end{proof}
The modal triple supplies more than endofunctors: the left-exact reflector
\(\Im\) has a relative form in every slice. Its defining geometry is the
pullback square
\[
\begin{tikzcd}[column sep=huge,row sep=large]
\Im_B E \arrow[r] \arrow[d]
    & \Im E \arrow[d]
\\
B \arrow[r]
    & \Im B ,
\end{tikzcd}
\]
which will be used as the formally \(\acute{e}\)tale reflection of a map over
\(B\).

\subsection{Formal \texorpdfstring{\'etalification}{etalification} and the lifting comparison}

\begin{definition}[Formally \texorpdfstring{\'etale}{etale} and formally smooth morphisms]
\label{def:infinitesimal-shape-etale}
Let \(g:E\to B\) be a morphism in \(\mathcal H_{\inf}\). It is
\textbf{formally \(\acute{e}\)tale} if the naturality square of the unit
\(\id\to\Im\) is Cartesian, equivalently if the canonical comparison
\[
    E\longrightarrow B\times_{\Im B}\Im E
\]
is an equivalence. It is \textbf{formally smooth} if the same comparison is an
effective epimorphism.
\end{definition}

\begin{construction}[Relative infinitesimal shape and formal etalification]
\label{con:formal-etalification}
For \(g:E\to B\) in \(\mathcal H_{\inf}\), define its \textbf{relative
infinitesimal shape}, or \textbf{formal etalification}, by
\[
    \Im_B E
    :=
    B\times_{\Im B}\Im E,
\]
with its projection \(\Im_Bg:\Im_BE\to B\). The unit and \(g\) give the
canonical factorization
\[
    E\xrightarrow{\theta_g}\Im_BE\xrightarrow{\Im_Bg}B.
\]
\end{construction}

\begin{proposition}[Relative infinitesimal-shape factorization]
\label{prop:infinitesimal-shape-factorization}
The morphism \(\theta_g\) becomes an equivalence after applying \(\Im\), while
\(\Im_Bg\) is formally \(\acute{e}\)tale. The construction is stable under
base change in the slice over \(B\). Moreover, \(\Im_Bg\) is the formally
\(\acute{e}\)tale reflection of \(g\): for every formally \(\acute{e}\)tale
\(Z\to B\), precomposition with \(\theta_g\) induces a natural equivalence
\[
    \Map_{\mathcal H_{\inf}/B}(\Im_BE,Z)
    \xrightarrow{\simeq}
    \Map_{\mathcal H_{\inf}/B}(E,Z).
\]
\end{proposition}

\begin{proof}
\leavevmode
\begin{enumerate}[label=\textup{(\arabic*)}]
    \item \textbf{Use exactness.}
    Left exactness and idempotence of \(\Im\) give
    \[
    \begin{aligned}
        \Im(\Im_BE)
        &\simeq
        \Im B\times_{\Im\Im B}\Im\Im E
        \\
        &\simeq
        \Im E,
    \end{aligned}
    \]
    so \(\Im(\theta_g)\) is an equivalence.

    \item \textbf{Verify naturality.}
    The defining square of \(\Im_BE\) is exactly the unit-naturality square for \(\Im_Bg\),
    proving formal \(\acute{e}\)taleness; left exactness gives base-change stability.

    \item \textbf{Apply the universal property.}
    For the universal property, let \(Z\to B\) be formally \(\acute{e}\)tale.

    \item \textbf{Establish the next comparison.}
    Then
    \[
        Z\simeq B\times_{\Im B}\Im Z.
    \]

    \item \textbf{map over is therefore equivalently a map.}
    A map \(E\to Z\) over \(B\) is therefore equivalently a map \(E\to\Im Z\) whose
    composite to \(\Im B\) is the prescribed one.

    \item \textbf{Use the stated property.}
    Since \(\Im Z\) is coreduced, the reflector gives
    \[
        \Map(E,\Im Z)\simeq\Map(\Im E,\Im Z),
    \]
    compatibly over \(\Im B\).

    \item \textbf{Conclude the step.}
    Hence
    \[
        \Map_B(E,Z)
        \simeq
        \Map_{\Im B}(\Im E,\Im Z).
    \]

    \item \textbf{Compute the mapping object.}
    The same calculation for \(\Im_BE\), using \(\Im(\Im_BE)\simeq\Im E\), gives the identical mapping space.

    \item \textbf{Transport the structure.}
    Under these identifications precomposition with \(\theta_g\) is the identity.
\end{enumerate}
\end{proof}

\begin{construction}[The formal lifting comparison]
\label{con:formal-lifting-comparison}
Let \(f:X\to S\) be a morphism in \(\mathcal H_{\Sp}\). Define
\[
    \Theta_f
    :=
    \theta_{i_!f}:
    i_!X
    \longrightarrow
    \Im_{i_!S}(i_!X).
\]
Using \(\Im(i_!X)\simeq i_*X\) and the corresponding equivalence for \(S\),
there is a canonical identification
\[
    \Im_{i_!S}(i_!X)
    \simeq
    i_!S\times_{i_*S}i_*X,
\]
so \(\Theta_f\) is the canonical comparison induced by \(i_!\to i_*\).
\end{construction}

\begin{proposition}[The formal lifting comparison is the etalification unit]
\label{prop:theta-is-etalification-unit}
The identification above exhibits \(\Theta_f\) as the unit of the relative
formally \(\acute{e}\)tale reflection of \(i_!f\).
\end{proposition}

\begin{proof}
\leavevmode
\begin{enumerate}[label=\textup{(\arabic*)}]
    \item \textbf{Identify the resulting comparison.}
    This is \Ref{prop:infinitesimal-shape-forgets-variation} substituted into the definition of
    \(\Im_{i_!S}(i_!X)\).
\end{enumerate}
\end{proof}

\begin{proposition}[Pointwise interpretation of the lifting comparison]
\label{prop:formal-lifting-pointwise}
For every object
\[
    j:U\longrightarrow T
\]
of \(\Aff_{\Sp}^{\inf}\),
\[
    \Theta_f(j):
    X(T)
    \longrightarrow
    S(T)\times_{S(U)}X(U).
\]
Its homotopy fibre over a compatible pair
\[
    T\longrightarrow S,
    \qquad
    U\longrightarrow X
\]
is the space of extensions \(T\to X\) together with the specified
compatibility homotopies.
\end{proposition}

\begin{proof}
\leavevmode
\begin{enumerate}[label=\textup{(\arabic*)}]
    \item \textbf{Identify the map.}
    The pointwise formulas for \(i_!\) and \(i_*\) give the displayed map, and limits of
    sheaves are computed objectwise.

    \item \textbf{Apply the universal property.}
    The fibre statement is the universal property of the pullback of spaces.
\end{enumerate}
\end{proof}

\begin{definition}[Formal smoothness and formal \texorpdfstring{\'etaleness}{etaleness} over ordinary spectral stacks]
\label{def:formal-smooth-etale}
A morphism \(f:X\to S\) in \(\mathcal H_{\Sp}\) is \textbf{formally
\(\acute{e}\)tale} if \(i_!f\) is formally \(\acute{e}\)tale in
\(\mathcal H_{\inf}\), equivalently if \(\Theta_f\) is an equivalence. It
is \textbf{formally smooth} if \(i_!f\) is formally smooth, equivalently if
\(\Theta_f\) is an effective epimorphism.
\end{definition}

\begin{theorem}[Formal etaleness is etaleness for infinitesimal shape]
\label{thm:formal-etale-infinitesimal-shape}
For a morphism \(f:X\to S\) of spectral Zariski stacks,
\[
    f\text{ is formally }\acute{e}\text{tale}
    \quad\Longleftrightarrow\quad
    i_!f\text{ is }\Im\text{-closed}.
\]
Equivalently, the unit-naturality square
\[
\begin{tikzcd}[column sep=large,row sep=large]
    i_!X \arrow[r] \arrow[d,"i_!f"']
        & \Im(i_!X) \arrow[d,"\Im(i_!f)"]
    \\
    i_!S \arrow[r]
        & \Im(i_!S)
\end{tikzcd}
\]
is Cartesian.
\end{theorem}

\begin{proof}
\leavevmode
\begin{enumerate}[label=\textup{(\arabic*)}]
    \item \textbf{Identify the resulting comparison.}
    This is the definition of a formally \(\acute{e}\)tale morphism in \(\mathcal H_{\inf}\), together with
    \Ref{prop:theta-is-etalification-unit}.
\end{enumerate}
\end{proof}

\begin{proposition}[Base change of formal lifting properties]
\label{prop:formal-lifting-base-change}
Formally \(\acute{e}\)tale and formally smooth morphisms in
\(\mathcal H_{\inf}\) are stable under base change. Consequently the same is
true for formally \(\acute{e}\)tale and formally smooth morphisms in
\(\mathcal H_{\Sp}\).
\end{proposition}

\begin{proof}
\leavevmode
\begin{enumerate}[label=\textup{(\arabic*)}]
    \item \textbf{Use exactness.}
    Left exactness of \(\Im\) identifies the relative etalification comparison of a pullback
    with the pullback of the original comparison.

    \item \textbf{Equivalences and effective epimorphisms in an -topos.}
    Equivalences and effective epimorphisms in an \(\infty\)-topos are stable under pullback.

    \item \textbf{Use preservation.}
    For morphisms in \(\mathcal H_{\Sp}\), the functor \(i_!\) preserves finite limits, so the same
    argument applies after the fully faithful inclusion.
\end{enumerate}
\end{proof}
Relative infinitesimal shape provides the canonical comparison from a map to
its formally \(\acute{e}\)tale reflection. Taking its kernel pair turns that
comparison into an intrinsic relation of infinitesimal equality.

\subsection{Relative infinitesimal disks and infinitesimal equality}

\begin{definition}[Relative infinitesimal disk bundle and groupoid]
\label{def:first-infinitesimal-neighbour-relation}
For \(f:X\to S\) in \(\mathcal H_{\Sp}\), define the \textbf{relative
infinitesimal disk bundle}
\[
    T_{\inf}(X/S)
    :=
    i_!X
    \times_{\Im_{i_!S}(i_!X)}
    i_!X.
\]
Its full \textbf{relative infinitesimal disk groupoid} is the \v{C}ech nerve
\[
    T_{\inf}^{[\bullet]}(X/S)
    :=
    \check C(\Theta_f).
\]
Thus \(T_{\inf}(X/S)\) is the degree-one relation object of the \v{C}ech
groupoid generated by the constant infinitesimal-path inclusion relative to
\(S\).
\end{definition}

\begin{proposition}[Pointwise infinitesimal equality]
\label{prop:pointwise-infinitesimal-equality}
For a finite infinitesimal context \(j:U\to T\),
\[
    T_{\inf}(X/S)(j)
    \simeq
    X(T)
    \times_{S(T)\times_{S(U)}X(U)}
    X(T).
\]
Thus two \(S\)-points of \(X\) over \(T\) are infinitesimally neighbouring
exactly when their restrictions to the reduction \(U\) agree, with the
specified compatibility over \(S\).
\end{proposition}

\begin{proof}
\leavevmode
\begin{enumerate}[label=\textup{(\arabic*)}]
    \item \textbf{Evaluate the defining kernel pair objectwise.}
    Evaluate the defining kernel pair objectwise and apply \Ref{prop:formal-lifting-pointwise}.
\end{enumerate}
\end{proof}

\begin{remark}[Infinitesimal equality is generated by infinitesimal shape]
\label{rem:infinitesimal-equality-reduction}
The groupoid \(T_{\inf}^{[\bullet]}(X/S)\) is the \v{C}ech groupoid of
the canonical map to relative infinitesimal shape. Its objects, disk pairs,
and higher compositions are therefore generated functorially by the unit of
\(\Im\).
\end{remark}
The modal formulas are pointwise on affine finite-infinitesimal arrows. Later
lifting arguments also use represented finite infinitesimal arrows that need
not themselves be affine objects of the test category. We therefore extend
Yoneda only as far as needed to evaluate the same comparison on those
represented arrows.

\subsection{Infinitesimal Yoneda realization of represented finite arrows}

\begin{construction}[Infinitesimal Yoneda realization]
\label{con:infinitesimal-yoneda-realization}
Let
\[
    j:U\longrightarrow T
\]
be a finite spectral infinitesimal substitution of spectral schemes, without
an affineness requirement on \(U\) or \(T\). Define a presheaf on
\(\Aff_{\Sp}^{\inf}\) by
\[
    \mathbf h^{\inf}_j(k:V\to W)
    :=
    \Map_{\Fun(\Delta^1,\operatorname{Sch}^{\mathrm{sp}})}(k,j).
\]
Since \(\Aff_{\Sp}^{\inf}\) is a full subcategory of the arrow category, this
is the restriction of the ordinary arrow-category Yoneda functor along the
affine finite-infinitesimal test category.
\end{construction}

\begin{proposition}[Descent and affine recovery for infinitesimal Yoneda realization]
\label{prop:infinitesimal-yoneda-realization-sheaf}
The presheaf \(\mathbf h^{\inf}_j\) is a
\(\tau^{\inf}_{\mathrm{Zar}}\)-sheaf. Hence it determines an object of
\(\mathcal H_{\inf}\). If \(j\) is an object of \(\Aff_{\Sp}^{\inf}\),
then there is a canonical equivalence
\[
    \mathbf h^{\inf}_j\simeq h_j
\]
with the representable infinitesimal sheaf.
\end{proposition}

\begin{proof}
\leavevmode
\begin{enumerate}[label=\textup{(\arabic*)}]
    \item \textbf{Compute the mapping object.}
    For an affine finite-infinitesimal test arrow \(k:V\to W\), the arrow-category mapping space
    is the pullback
    \[
    \begin{aligned}
        \mathbf h^{\inf}_j(k)
        \simeq
        \Map_{\operatorname{Sch}^{\mathrm{sp}}}(V,U)
        \times_{\Map_{\operatorname{Sch}^{\mathrm{sp}}}(V,T)}
        \Map_{\operatorname{Sch}^{\mathrm{sp}}}(W,T).
    \end{aligned}
    \]
    For a codomain-pullback Zariski basis cover of \(k\), the targets form a spectral
    Zariski cover of \(W\) and the sources form the corresponding spectral Zariski cover
    of \(V\).

    \item \textbf{Apply Zariski descent.}
    Represented spectral schemes satisfy spectral Zariski descent, and limits commute with
    limits in spaces.

    \item \textbf{Identify the map.}
    The preceding pullback of mapping functors therefore satisfies descent.

    \item \textbf{Use the stated property.}
    Since the basis covers generate \(\tau^{\inf}_{\mathrm{Zar}}\), the presheaf is a sheaf.

    \item \textbf{Compute the mapping object.}
    If \(j\) is an object of \(\Aff_{\Sp}^{\inf}\), both source and target lie in the affine
    infinitesimal arrow category and the displayed mapping space is exactly \(\Map_{\Aff_{\Sp}^{\inf}}(k,j)\).

    \item \textbf{Identify the resulting object.}
    Subcanonicity then identifies the sheaf with the representable \(h_j\).
\end{enumerate}
\end{proof}

\begin{construction}[Represented formal lifting comparison]
\label{con:represented-formal-lifting-comparison}
Let
\[
    f:X\longrightarrow S
\]
be a morphism in \(\mathcal H_{\Sp}\), and let
\[
    j:U\longrightarrow T
\]
be a finite spectral infinitesimal substitution of spectral schemes, without
an affineness requirement on \(U\) or \(T\). Define the
\textbf{represented formal lifting comparison} of \(f\) at \(j\) by
\[
\begin{aligned}
    \Theta_f^{\mathrm{rep}}(j):
    \Map_{\mathcal H_{\Sp}}(T,X)
    \longrightarrow{}&
    \Map_{\mathcal H_{\Sp}}(T,S)
    \\
    &\times_{\Map_{\mathcal H_{\Sp}}(U,S)}
    \Map_{\mathcal H_{\Sp}}(U,X),
\end{aligned}
\]
where the map sends
\[
    b:T\longrightarrow X
\]
to the pair
\[
    f\circ b:T\longrightarrow S,
    \qquad
    b\circ j:U\longrightarrow X,
\]
together with the canonical equality of their restrictions to \(U\).

If \(j\) is an object of \(\Aff_{\Sp}^{\inf}\), then the endpoint formulas of
\Ref{prop:endpoint-kan-extensions} and the pointwise computation of
\Ref{prop:formal-lifting-pointwise} canonically identify
\[
    \Theta_f^{\mathrm{rep}}(j)
    \simeq
    \Theta_f(j).
\]
Thus \(\Theta_f^{\mathrm{rep}}\) extends the already constructed affine
pointwise lifting comparison to represented finite infinitesimal arrows. It
does not declare a nonaffine arrow to be an object of
\(\Aff_{\Sp}^{\inf}\), nor does it introduce an additional infinitesimal-site
structure.
\end{construction}
This gives the endpoint interface needed for formal geometry: relative
infinitesimal shape supplies the lifting comparison, infinitesimal disks its
kernel-pair relation, and represented finite arrows the same pointwise
interpretation. The remaining problem is to compute the fibres of the lifting
comparison.

\section{Cotangent analysis of infinitesimal lifting}
\label{sec:lifting-judgments}

For every morphism \(f:X\to S\), differential cohesion has already constructed
the relative formal-etalification unit
\[
    \Theta_f:i_!X\longrightarrow\Im_{i_!S}(i_!X).
\]
Its fibres are nonlinear lifting spaces. On a structured square-zero stage,
the cotangent classifier turns each such fibre into a nullhomotopy problem in a
stable coefficient category. The geometric and linear descriptions fit into
the pullback picture
\[
\begin{tikzcd}[column sep=huge,row sep=large]
\operatorname{Lift}_S(a;i) \arrow[r] \arrow[d]
    & X(T) \arrow[d,"{\Theta_f^{\mathrm{rep}}(i)}"]
\\
* \arrow[r,"{(T\to S,a)}"']
    & S(T)\times_{S(U)}X(U).
\end{tikzcd}
\]
This section first computes that fibre and then uses one-stage detection to
recognize formal \(\acute{e}\)taleness and formal smoothness from cotangent
data.

\begin{proposition}[Formal lifting across a fixed affine square-zero extension]
\label{prop:formal-lifting-fixed-square-zero}
Let
\[
    A\longrightarrow B,
    \qquad
    A\longrightarrow C
\]
be morphisms of commutative ring spectra. Let
\[
    q:C'\longrightarrow C
\]
be a structured square-zero extension in commutative \(A\)-algebras by a
\(C\)-module \(J\), classified by
\[
    \delta:L_{C/A}\longrightarrow J[1].
\]
For an \(A\)-algebra morphism \(\phi:B\to C\), define the obstruction term
\[
    o(\phi):
    L_{B/A}\otimes_BC
    \longrightarrow
    L_{C/A}
    \xrightarrow{\delta}
    J[1],
\]
where the first map is induced by cotangent functoriality. Then the space of
lifts of \(\phi\) to an \(A\)-algebra morphism \(B\to C'\) is naturally
equivalent to the space of nullhomotopies of \(o(\phi)\).

Consequently:
\begin{enumerate}[label=(\roman*)]
    \item a lift exists if and only if \([o(\phi)]\) is zero in
    \[
        \pi_0\map_{\Mod_C}
        (L_{B/A}\otimes_BC,J[1]);
    \]
    \item when nonempty, the lift space is a torsor under
    \[
        \Omega^\infty
        \map_{\Mod_C}(L_{B/A}\otimes_BC,J);
    \]
    \item if \(L_{B/A}\simeq0\), the lift space is contractible.
\end{enumerate}
\end{proposition}

\begin{proof}
\leavevmode
\begin{enumerate}[label=\textup{(\arabic*)}]
    \item \textbf{Fix the notation.}
    Put
    \[
        \mathcal A:=\CAlg(\Sp)_{A/}.
    \]

    \item \textbf{Identify the map.}
    The structured extension \(q:C'\to C\) is, by \Ref{con:structured-square-zero-extension}, the
    pullback in \(\mathcal A_{/C}\)
    \[
        q
        \simeq
        \id_C\times_p\id_C,
    \]
    where
    \[
        p:C\oplus J[1]\longrightarrow C
    \]
    is the split augmentation and the two maps \(\id_C\rightrightarrows p\) are \(d_\delta\) and \(d_0\).

    \item \textbf{Identify the map.}
    The morphism \(\phi:B\to C\) is an object of \(\mathcal A_{/C}\).

    \item \textbf{space of lifts of through , including the specified homotopy over.}
    The space of lifts of \(\phi\) through \(q\), including the specified homotopy over
    \(C\), is therefore
    \[
        \Map_{\mathcal A_{/C}}(\phi,q).
    \]

    \item \textbf{Compute the mapping object.}
    Mapping out of \(\phi\) preserves the defining pullback, so
    \[
    \begin{aligned}
        \Map_{\mathcal A_{/C}}(\phi,q)
        \simeq{}&
        \Map_{\mathcal A_{/C}}(\phi,\id_C)
        \\
        &\times_{\Map_{\mathcal A_{/C}}(\phi,p)}
        \Map_{\mathcal A_{/C}}(\phi,\id_C).
    \end{aligned}
    \]

    \item \textbf{Compute the mapping object.}
    The object \(\id_C\) is terminal in \(\mathcal A_{/C}\), hence each outer mapping space is
    contractible.

    \item \textbf{Their two images in.}
    Their two images in \(\Map_{\mathcal A_{/C}}(\phi,p)\) are the sections
    \[
        d_\delta\phi
        \qquad\text{and}\qquad
        d_0\phi.
    \]

    \item \textbf{Conclude the step.}
    Thus the lift space is canonically the path space between these two derivations.

    \item \textbf{Transport the structure.}
    Under the relative cotangent classifier, their difference is classified by the composite
    \[
        L_{B/A}\otimes_BC
        \longrightarrow
        L_{C/A}
        \xrightarrow{\delta}
        J[1]
        =o(\phi).
    \]

    \item \textbf{Conclude the step.}
    Consequently
    \[
        \operatorname{Lift}_A(\phi;q)
        \simeq
        \operatorname{Null}(o(\phi)).
    \]

    \item \textbf{Compute the mapping object.}
    A nullhomotopy exists exactly when the class of \(o(\phi)\) vanishes in \(\pi_0\) of the
    indicated mapping spectrum.

    \item \textbf{Identify the resulting object.}
    When the nullhomotopy space is nonempty, translation by a chosen nullhomotopy identifies it
    as a torsor under the loop space
    \[
        \Omega\Omega^\infty
        \map_{\Mod_C}(L_{B/A}\otimes_BC,J[1])
        \simeq
        \Omega^\infty
        \map_{\Mod_C}(L_{B/A}\otimes_BC,J).
    \]

    \item \textbf{Compute the mapping object.}
    If \(L_{B/A}\simeq0\), the relevant mapping spectrum is zero and the lift space is contractible.
\end{enumerate}
\end{proof}

\begin{remark}[Lifting as a judgment]
\label{rem:lifting-judgment}
The proposition may be read as the equivalence
\[
    \left\{
    \text{extensions of a term across an infinitesimal context}
    \right\}
    \simeq
    \left\{
    \text{nullhomotopies of its canonical tangent boundary}
    \right\}.
\]
The obstruction is the boundary term constructed from cotangent functoriality
and the specified infinitesimal extension.
\end{remark}

\begin{theorem}[Global lifting judgment]
\label{thm:global-lifting-judgment}
Let \(f:X\to S\) be a morphism of spectral schemes and let
\[
    i:U\longrightarrow T
\]
be a structured infinitesimal extension over \(S\) by
\(J\in\QCoh(U)_{\geq0}\), classified by
\[
    \delta:L_{U/S}\longrightarrow J[1].
\]
For an \(S\)-morphism \(a:U\to X\), define
\[
    o(a):
    a^*L_{X/S}
    \longrightarrow
    L_{U/S}
    \xrightarrow{\delta}
    J[1].
\]
The compatible pair
\[
    T\longrightarrow S,
    \qquad
    a:U\longrightarrow X
\]
defines a point of the target of the represented formal lifting comparison
\[
    \Theta_f^{\mathrm{rep}}(i):
    X(T)
    \longrightarrow
    S(T)\times_{S(U)}X(U)
\]
of \Ref{con:represented-formal-lifting-comparison}. Its homotopy fibre over
that point is the space of \(S\)-morphisms
\(\widetilde a:T\to X\) equipped with a specified homotopy
\[
    \widetilde a\circ i\simeq a,
\]
and there is a natural equivalence
\[
\begin{aligned}
    \operatorname{fib}_{(T\to S,a)}
    \bigl(\Theta_f^{\mathrm{rep}}(i)\bigr)
    &\simeq
    \operatorname{Lift}_S(a;i)
    \\
    &\simeq
    \operatorname{Null}_{\QCoh(U)}\bigl(o(a)\bigr).
\end{aligned}
\]
In particular, when nonempty, the lift space is a torsor under
\[
    \Omega^\infty
    \map_{\QCoh(U)}(a^*L_{X/S},J).
\]
When \(i\) is affine, the canonical comparison
\[
    \Theta_f^{\mathrm{rep}}(i)\simeq\Theta_f(i)
\]
recovers the pointwise differential-cohesive lifting comparison of
\Ref{prop:formal-lifting-pointwise}.
\end{theorem}

\begin{proof}
\leavevmode
\begin{enumerate}[label=\textup{(\arabic*)}]
    \item \textbf{Choose the local data.}
    Choose an affine Zariski cover of \(T\).

    \item \textbf{Invoke the cited result.}
    By \Ref{prop:global-square-zero-geometric-properties}, its pullback to \(U\) is the
    corresponding affine cover, and the square-zero extension on every member is represented by
    a structured affine square-zero extension.

    \item \textbf{Refine the image of each affine source.}
    Refine the image of each affine source in \(X\) and \(S\) by compatible affine
    opens.

    \item \textbf{Restrict the construction.}
    On every resulting affine piece, \Ref{prop:formal-lifting-fixed-square-zero} identifies the
    lift space with the nullhomotopy space of the restriction of \(o(a)\).

    \item \textbf{Verify naturality.}
    On finite intersections these equivalences agree by cotangent base change and naturality of
    the classifying derivation.

    \item \textbf{Apply Zariski descent.}
    Mapping descent for the represented spectral scheme \(X\) identifies the global lift
    space with the totalization of the local lift spaces.

    \item \textbf{Apply Zariski descent.}
    Quasi-coherent descent identifies
    \[
        \Map_{\QCoh(U)}(a^*L_{X/S},J[1])
    \]
    with the totalization of the corresponding affine mapping spaces, and limits in spaces
    preserve homotopy fibres over the zero point.

    \item \textbf{Identify the equivalence.}
    The affine nullhomotopy equivalences therefore totalize to the displayed global equivalence.

    \item \textbf{Invoke the cited result.}
    By definition,
    \[
        \Theta_f^{\mathrm{rep}}(i):
        X(T)\longrightarrow S(T)\times_{S(U)}X(U)
    \]
    has as its fibre over the compatible pair \((T\to S,a)\) precisely the space of maps
    \[
        \widetilde a:T\longrightarrow X
    \]
    over \(S\), equipped with the specified homotopy \(\widetilde a i\simeq a\).

    \item \textbf{Conclude the step.}
    Hence this fibre is the same global lift space computed above.

    \item \textbf{Compute the mapping object.}
    The torsor statement follows by looping the mapping spectrum.

    \item \textbf{If is affine, the final assertion follows.}
    If \(i\) is affine, the final assertion follows from
    \Ref{con:represented-formal-lifting-comparison} and \Ref{prop:formal-lifting-pointwise}.
\end{enumerate}
\end{proof}
The structured one-stage calculation is the local building block. We now show
that testing only square-zero stages already detects the lifting behavior on
arbitrary finite infinitesimal substitutions, after which the cotangent
criteria become intrinsic statements about represented spectral contexts.

\subsection{One-stage detection and cotangent recognition}

\begin{proposition}[One-stage detection]
\label{prop:formal-lifting-one-stage-detection}
Let \(f:X\to S\) be a morphism in \(\mathcal H_{\Sp}\).

Formal \(\acute{e}\)taleness is equivalent to contractible lifting against
every affine square-zero infinitesimal substitution, equivalently against
every one-stage square-zero object
\[
    j:U\longrightarrow T
\]
of \(\Aff_{\Sp}^{\inf}\).

Formal smoothness is equivalent to Zariski-local existence of lifts against
every such affine square-zero infinitesimal substitution.
\end{proposition}

\begin{proof}
\leavevmode
\begin{enumerate}[label=\textup{(\arabic*)}]
    \item \textbf{forward implications.}
    The forward implications follow by evaluating \(\Theta_f\) on one-stage objects of
    \(\Aff_{\Sp}^{\inf}\).

    \item \textbf{Construct the factorization.}
    Conversely, let
    \[
        j:U\longrightarrow T
    \]
    be an arbitrary object of \(\Aff_{\Sp}^{\inf}\) and choose, only inside the proof, a finite
    factorization
    \[
        U=U_0
        \longrightarrow
        U_1
        \longrightarrow
        \cdots
        \longrightarrow
        U_r=T
    \]
    into underlying square-zero substitutions.

    \item \textbf{Each intermediate.}
    Each intermediate \(U_k\) is affine: work backward from the affine target \(T\),
    using affineness of every square-zero substitution.

    \item \textbf{Identify contractibility.}
    For formal \(\acute{e}\)taleness, the lift space across the first stage is contractible by the
    one-stage condition.

    \item \textbf{Compute the fibre.}
    Applying the same statement to the resulting lift at the next stage and iterating expresses
    the lift space for the composite as a finite iterated homotopy fibre of contractible spaces;
    it is therefore contractible.

    \item \textbf{Identify contractibility.}
    Hence every fibre of \(\Theta_f(j)\) is contractible, so \(\Theta_f\) is an equivalence.

    \item \textbf{For formal smoothness, construct a lift successively.}
    For formal smoothness, construct a lift successively after Zariski refinement.

    \item \textbf{Identify the map.}
    The square-zero map
    \[
        U_k\longrightarrow U_{k+1}
    \]
    identifies their small Zariski sites by \Ref{prop:square-zero-substitution-permanence}, so a
    cover on one side transports uniquely to the other.

    \item \textbf{Apply the one-stage local lifting property, refine finitely at each.}
    Apply the one-stage local lifting property, refine finitely at each step, and compose the
    resulting covers.

    \item \textbf{After the finite number of stages a.}
    After the finite number of stages a lift exists over a Zariski cover of \(T\).

    \item \textbf{Verify surjectivity.}
    Thus \(\Theta_f\) is locally surjective on sections and hence an effective epimorphism.

    \item \textbf{Construct the factorization.}
    The factorization and intermediate covers are proof witnesses only.
\end{enumerate}
\end{proof}

\begin{theorem}[Cotangent recognition of formal \texorpdfstring{\'etaleness}{etaleness}]
\label{thm:cotangent-recognition-formal-etale}
Let \(f:X\to S\) be represented by spectral schemes. Then
\[
    f\text{ is formally }\acute{e}\text{tale}
    \quad\Longleftrightarrow\quad
    L_{X/S}\simeq0.
\]
\end{theorem}

\begin{proof}
\leavevmode
\begin{enumerate}[label=\textup{(\arabic*)}]
    \item \textbf{Suppose first that .}
    Suppose first that \(L_{X/S}\simeq0\).

    \item \textbf{Invoke the cited result.}
    By \Ref{prop:formal-lifting-one-stage-detection}, it suffices to solve a one-stage
    underlying square-zero lifting problem.

    \item \textbf{Fix the notation.}
    Let
    \[
            j:U\longrightarrow T
        \]
    be an affine square-zero infinitesimal substitution over \(S\), together with
    \(a:U\to X\).

    \item \textbf{Choose the local data.}
    Choose, only inside the proof, a structured presentation of \(j\).

    \item \textbf{Construct the factorization.}
    Refine \(T\) by a finite affine Zariski cover whose composites with \(S\) factor
    through affine opens of \(S\).

    \item \textbf{On each member the chosen presentation may.}
    On each member the chosen presentation may be refined relative to that affine base by
    \Ref{lem:underlying-square-zero-relative-refinement}.

    \item \textbf{Identify contractibility.}
    The global lifting calculation \Ref{thm:global-lifting-judgment} then identifies the local
    lift space with the nullhomotopy space of a map out of the pullback of \(L_{X/S}\), hence
    with a contractible space.

    \item \textbf{Apply Zariski descent.}
    The same is true on all finite intersections, and mapping descent totalizes these
    contractible local lift spaces to a contractible global lift space.

    \item \textbf{Identify contractibility.}
    Thus every one-stage lifting problem has a contractible solution space, and \(f\) is
    formally \(\acute{e}\)tale.

    \item \textbf{Prove the converse.}
    Conversely suppose \(f\) is formally \(\acute{e}\)tale.

    \item \textbf{Construct the factorization.}
    Let
    \[
            U=\Spec(B)\hookrightarrow X
        \]
    be an affine open whose composite with \(S\) factors through an affine open
    \[
            V=\Spec(A)\hookrightarrow S.
        \]

    \item \textbf{Verify connectivity.}
    For a connective \(B\)-module \(M\), consider the split square-zero test
    \[
            j:U=\Spec(B)\longrightarrow T:=\Spec(B\oplus M).
        \]

    \item \textbf{Identify the map.}
    The structural map \(T\to V\to S\) and \(U\to X\) define a one-stage lifting problem.

    \item \textbf{Identify contractibility.}
    Formal \(\acute{e}\)taleness makes its lift space contractible.

    \item \textbf{Construct the factorization.}
    Every lift factors through \(U\subseteq X\): the inverse image of \(U\) is an open
    subobject of \(T\) containing the centre, and \(U\to T\) is a homeomorphism on
    underlying topological spaces.

    \item \textbf{Identify contractibility.}
    Hence that open is all of \(T\), and the factorization space is contractible because
    an open immersion is a monomorphism.

    \item \textbf{lift space.}
    The lift space is therefore the space of \(A\)-algebra sections
    \[
            B\longrightarrow B\oplus M,
        \]
    namely
    \[
            \Omega^\infty\map_{\Mod_B}(L_{B/A},M).
        \]

    \item \textbf{Verify connectivity.}
    Take \(M=L_{B/A}\), which is connective.

    \item \textbf{Establish the next comparison.}
    Contractibility makes the identity of \(L_{B/A}\) null, so \(L_{B/A}\simeq0\).

    \item \textbf{Glue the local data.}
    These affine vanishings are compatible under restriction and glue by
    \Ref{thm:global-cotangent-classifier}, giving \(L_{X/S}\simeq0\).
\end{enumerate}
\end{proof}

\begin{proposition}[Finite locally free cotangent complexes give formal smoothness]
\label{prop:finite-locally-free-formal-smooth}
Let \(f:X\to S\) be a morphism of spectral schemes. If \(L_{X/S}\) is
finite locally free, then \(f\) is formally smooth.
\end{proposition}

\begin{proof}
\leavevmode
\begin{enumerate}[label=\textup{(\arabic*)}]
    \item \textbf{Invoke the cited result.}
    By \Ref{prop:formal-lifting-one-stage-detection}, it suffices to solve a one-stage
    underlying square-zero lifting problem Zariski-locally.

    \item \textbf{Fix the notation.}
    Let \(U\to T\) be an affine square-zero infinitesimal substitution over \(S\) and let
    \(U\to X\) be the given map.

    \item \textbf{Cover its image in by affine opens.}
    Cover its image in \(X\) by affine opens, pull the cover back to the affine
    \(U\), and choose a finite principal refinement.

    \item \textbf{Zariski-site invariance across transports this to a finite.}
    Zariski-site invariance across \(U\to T\) transports this to a finite affine cover
    \(\{T_\alpha\to T\}\), with \(U_\alpha\simeq U\times_TT_\alpha\), such that each restricted map lands in an affine open of
    \(X\).

    \item \textbf{On each member choose, only inside the.}
    On each member choose, only inside the proof, a structured presentation of the underlying
    square-zero arrow and refine it relative to an affine open of the base \(S\) using
    \Ref{lem:underlying-square-zero-relative-refinement}.

    \item \textbf{Verify connectivity.}
    The obstruction of \Ref{prop:formal-lifting-fixed-square-zero} then lies in
    \[
        \pi_0\map_C(P,J[1]),
    \]
    where \(P\) is the pullback of \(L_{X/S}\) and \(J\) is the connective
    coefficient of that proof presentation.

    \item \textbf{Finite local freeness.}
    Finite local freeness makes \(P\) finite projective, hence
    \[
        \map_C(P,J[1])
        \simeq
        P^\vee\otimes_CJ[1].
    \]

    \item \textbf{Verify connectivity.}
    The right side is \(1\)-connective, so its \(\pi_0\) vanishes.

    \item \textbf{Derive the comparison.}
    The obstruction is therefore zero and the fixed-extension lifting theorem gives a lift over
    each \(T_\alpha\).

    \item \textbf{These local lifts.}
    These local lifts are precisely the effective epimorphism condition for \(\Theta_f\).
\end{enumerate}
\end{proof}
The lifting comparison has now been identified cotangentially and its formal
\(\acute{e}\)tale and smooth regimes have been recognized. The relative
infinitesimal disk groupoid contains more than these lifting fibres, however:
its diagonal relation has a canonical first-order piece whose difference map
is itself universal.

\subsection{The canonical first infinitesimal neighbourhood of the diagonal}
\label{sec:first-diagonal-neighbourhood}

We isolate that first-order piece as a structured thickening of the diagonal.
Its two retractions to \(X\) agree on the centre, and their first-order
difference recovers the universal relative derivation. This identifies the
cotangent classifier as the linearization of infinitesimal equality without
identifying all finite infinitesimal structure with first-order data.

\begin{construction}[First structured infinitesimal neighbourhood]
\label{con:first-diagonal-neighbourhood}
Let \(p:X\to S\) be a morphism of spectral schemes and put
\[
    Y:=X\times_SX.
\]
Let \(\Delta:X\to Y\) be the diagonal and regard \(Y\) over \(X\) through
\(\operatorname{pr}_1\). Global cotangent base change gives
\[
    L_{Y/X}
    \simeq
    \operatorname{pr}_2^*L_{X/S},
    \qquad
    \Delta^*L_{Y/X}
    \simeq
    L_{X/S}.
\]
Since
\[
    X\xrightarrow{\Delta}Y\xrightarrow{\operatorname{pr}_1}X
\]
is the identity, transitivity supplies a canonical connecting equivalence
\[
    \partial_\Delta:
    L_{X/Y}
    \xrightarrow{\simeq}
    L_{X/S}[1].
\]
Apply the global structured infinitesimal context-extension construction to \(\partial_\Delta\). This gives a
canonical structured square-zero thickening
\[
    X
    \xrightarrow{i_1}
    D^{(1)}_{X/S}
    \xrightarrow{\nu_1}
    X\times_SX
\]
with coefficient \(L_{X/S}\). Define
\[
    h_0:=\operatorname{pr}_1\nu_1,
    \qquad
    h_1:=\operatorname{pr}_2\nu_1
    :D^{(1)}_{X/S}\rightrightarrows X.
\]
Both are retractions of the centre:
\[
    h_0i_1\simeq\id_X,
    \qquad
    h_1i_1\simeq\id_X.
\]
\end{construction}

\begin{proposition}[Base change of the first neighbourhood]
\label{prop:first-diagonal-base-change}
The construction \(D^{(1)}_{X/S}\), its centre, the two retractions, and the
coefficient \(L_{X/S}\) commute with arbitrary Cartesian base change of
represented spectral contexts.
\end{proposition}

\begin{proof}
\leavevmode
\begin{enumerate}[label=\textup{(\arabic*)}]
    \item \textbf{Apply cotangent base change.}
    The diagonal, global cotangent base change, transitivity, the connecting morphism
    \(\partial_\Delta\), and global structured infinitesimal context extension all commute with
    represented Cartesian base change.

    \item \textbf{two retractions.}
    The two retractions are obtained by composing with the two product projections.
\end{enumerate}
\end{proof}

\begin{theorem}[The two diagonal retractions differ by the universal derivation]
\label{thm:first-neighbourhood-universal-derivation}
Let \(p:X\to S\). Split the first structured neighbourhood by the retraction
\(h_0\). Under the resulting identification of its square-zero coefficient
with \(L_{X/S}\), the second retraction \(h_1\) is the graph of the universal
relative \(S\)-derivation. Equivalently, on every affine base test
\(S=\Spec(R)\) and affine open \(X=\Spec(A)\), the induced derivation is
\[
    d_{A/R}:A\longrightarrow L_{A/R}.
\]
Thus the first-order difference of \(h_1\) and \(h_0\) is the universal
relative derivation of the structure algebra; it is not an
\(\mathcal O_X\)-linear map from \(\mathcal O_X\) to \(L_{X/S}\).
\end{theorem}

\begin{proof}
\leavevmode
\begin{enumerate}[label=\textup{(\arabic*)}]
    \item \textbf{Set up the argument.}
    The assertion is local on \(X\), so take \(S=\Spec(R)\) and \(X=\Spec(A)\).

    \item \textbf{Fix the notation.}
    Put
    \[
            M:=L_{A/R},
            \qquad
            B:=A\otimes_RA.
        \]

    \item \textbf{Fix the notation.}
    Let
    \[
            j_0,j_1:A\rightrightarrows B
        \]
    be the two tensor-factor maps and let \(\mu:B\to A\) be multiplication.

    \item \textbf{Regard as an -algebra through .}
    Regard \(B\) as an \(A\)-algebra through \(j_0\).

    \item \textbf{Apply cotangent base change.}
    Cotangent base change gives
    \[
            L_{B/A}
            \simeq
            B\otimes_{A,j_1}M,
        \]
    and hence
    \[
            P:=A\otimes_BL_{B/A}
            \simeq
            M.
        \]

    \item \textbf{Apply cotangent transitivity.}
    The composite \(A\xrightarrow{j_0}B\xrightarrow{\mu}A\) is the identity, so transitivity gives
    \[
            P\longrightarrow0\longrightarrow L_{A/B}
        \]
    and therefore a connecting equivalence
    \[
            \partial:L_{A/B}
            \xrightarrow{\simeq}
            P[1]
            \simeq
            M[1].
        \]

    \item \textbf{Restrict the construction.}
    This is the affine restriction of \(\partial_\Delta\).

    \item \textbf{Fix the notation.}
    Let \(q:A'\to A\) be the structured square-zero \(B\)-algebra extension classified by
    \(\partial\), and let \(\beta:B\to A'\) be its \(B\)-algebra structure.

    \item \textbf{Identify the map.}
    The two tensor-factor maps give two sections of \(q\),
    \[
            \alpha_i:=\beta j_i:A\longrightarrow A',
            \qquad i=0,1,
        \]
    because \(q\alpha_i=\mu j_i=\id_A\).

    \item \textbf{These algebra sections.}
    These algebra sections are contravariantly the geometric retractions \(h_i\).

    \item \textbf{Restrict the construction.}
    Restrict the structured extension along \(j_0:A\to B\).

    \item \textbf{Identify contractibility.}
    Its relative cotangent complex is \(L_{A/A}\simeq0\), so square-zero classification identifies it,
    through a contractible space, with the split extension
    \[
            A'\simeq A\oplus M
        \]
    in such a way that
    \[
            \alpha_0(a)=(a,0).
        \]

    \item \textbf{There is therefore a unique -derivation.}
    There is therefore a unique \(R\)-derivation
    \[
            \delta:A\longrightarrow M
        \]
    with
    \[
            \alpha_1(a)=(a,\delta(a)).
        \]

    \item \textbf{Transport the structure.}
    Under the same splitting, write
    \[
            \beta(b)=(\mu(b),D(b)),
        \]
    where \(D:B\to M\) is a derivation relative to \(j_0\).

    \item \textbf{Its classifying -linear morphism.}
    Its classifying \(A\)-linear morphism is
    \[
            \bar D:P=A\otimes_BL_{B/A}\longrightarrow M.
        \]

    \item \textbf{Apply cotangent transitivity.}
    Applying \(\map_{\Mod_A}(-,M[1])\) to the transitivity sequence \(P\to0\to L_{A/B}\) identifies maps \(P\to M\)
    with maps \(L_{A/B}\to M[1]\).

    \item \textbf{Transport the structure.}
    Under this connecting equivalence, the classifying map \(\partial\) corresponds to the
    identity of \(P\simeq M\).

    \item \textbf{Conclude the step.}
    Hence
    \[
            \bar D=\id_M.
        \]

    \item \textbf{Restrict the construction.}
    The derivation \(\delta\) is the restriction of \(D\) along \(j_1\).

    \item \textbf{Apply base change.}
    Under the base-change identification \(P\simeq L_{A/R}=M\), its classifying morphism is therefore
    \(\id_M\).

    \item \textbf{Apply the universal property.}
    The universal property of the cotangent complex gives
    \[
            \delta=d_{A/R}.
        \]

    \item \textbf{Restrict the construction.}
    The affine identifications commute with restriction and therefore globalize.
\end{enumerate}
\end{proof}

\begin{remark}[First order is not all infinitesimal order]
\label{rem:first-order-not-all-infinitesimal}
The theorem identifies the first structured neighbourhood of the diagonal with
the universal first-order variation. Higher finite infinitesimal relations may
contain data beyond the classifier \(L_{X/S}\), and higher infinitesimal
functions may carry structure beyond exterior powers of the cotangent
classifier. The first neighbourhood is the canonical linear quotient of the
infinitesimal geometry.
\end{remark}

\begin{proposition}[The first diagonal neighbourhood is a represented first-order infinitesimal disk]
\label{prop:first-diagonal-neighbour-pair}
Let \(p:X\to S\), and let
\[
    i_1:X\longrightarrow D^{(1)}_{X/S}
\]
be the canonical first structured infinitesimal neighbourhood. The two
retractions
\[
    h_0,h_1:D^{(1)}_{X/S}\rightrightarrows X
\]
determine a canonical morphism of infinitesimal sheaves
\[
    (h_0,h_1):
    \mathbf h^{\inf}_{i_1}
    \longrightarrow
    T_{\inf}(X/S).
\]
If \(X\) is affine---equivalently, by
\Ref{lem:finite-infinitesimal-affineness-invariance}, if
\(D^{(1)}_{X/S}\) is affine---then \(i_1\) is an object of
\(\Aff_{\Sp}^{\inf}\), and the affine recovery of \Ref{prop:infinitesimal-yoneda-realization-sheaf} identifies
this morphism by Yoneda with a canonical point
\[
    (h_0,h_1)
    \in
    T_{\inf}(X/S)(i_1).
\]
Under the splitting by \(h_0\), the first-order difference of this represented
neighbour is the universal relative derivation.
\end{proposition}

\begin{proof}
\leavevmode
\begin{enumerate}[label=\textup{(\arabic*)}]
    \item \textbf{Fix the notation.}
    Let \(k:V\to W\) be an affine finite-infinitesimal test arrow and let
    \[
        k\longrightarrow i_1
    \]
    be a point of \(\mathbf h^{\inf}_{i_1}(k)\).

    \item \textbf{Verify compatibility.}
    Thus we have a commutative square
    \[
    \begin{tikzcd}[column sep=large,row sep=large]
        V \arrow[r] \arrow[d,"k"']
            & X \arrow[d,"i_1"]
        \\
        W \arrow[r]
            & D^{(1)}_{X/S}.
    \end{tikzcd}
    \]

    \item \textbf{Compose the lower.}
    Compose the lower map with \(h_0\) and \(h_1\).

    \item \textbf{Identify the resulting comparison.}
    This produces two maps
    \[
        W\rightrightarrows X
    \]
    over the same map \(W\to S\).

    \item \textbf{Restrict the construction.}
    Their restrictions to \(V\) agree with the given map \(V\to X\), because
    \[
        h_0i_1\simeq\id_X\simeq h_1i_1.
    \]

    \item \textbf{Invoke the cited result.}
    By the pointwise formula of \Ref{prop:pointwise-infinitesimal-equality}, this is a point of
    \(T_{\inf}(X/S)(k)\).

    \item \textbf{Identify the map.}
    The construction is natural in both \(k\) and the map \(k\to i_1\), hence defines the
    displayed morphism of sheaves.

    \item \textbf{When is affine, the cited result makes.}
    When \(X\) is affine, \Ref{lem:finite-infinitesimal-affineness-invariance} makes
    \(D^{(1)}_{X/S}\) affine as well, so \(\mathbf h^{\inf}_{i_1}\simeq h_{i_1}\).

    \item \textbf{Apply Yoneda.}
    Ordinary Yoneda in the infinitesimal topos gives the stated point.

    \item \textbf{final assertion.}
    The final assertion is exactly \Ref{thm:first-neighbourhood-universal-derivation}.
\end{enumerate}
\end{proof}

\begin{corollary}[Cotangent classification linearizes infinitesimal equality]
\label{cor:cotangent-linearizes-infinitesimal-equality}
For every represented relative spectral context \(X\to S\), the cotangent
classifier \(L_{X/S}\) is the universal linear coefficient measuring the
first-order difference of infinitesimally neighbouring points along the
canonical first neighbourhood of the diagonal.
\end{corollary}

\begin{proof}
\leavevmode
\begin{enumerate}[label=\textup{(\arabic*)}]
    \item \textbf{Handle the first component.}
    The first diagonal disk pair is universal by
    \Ref{thm:first-neighbourhood-universal-derivation}; derivation judgments are represented by
    \(L_{X/S}\) by \Ref{def:global-derivation-term} and its affine classifier.

    \item \textbf{Conclude the step.}
    Thus every first-order linear difference is obtained by applying a coefficient map out of
    \(L_{X/S}\) to the universal one.
\end{enumerate}
\end{proof}
Two facts remain active from the lifting analysis: one-stage square-zero tests
control finite infinitesimal lifting, and the first diagonal neighbourhood
realizes the universal first-order difference geometrically. Applying the same
mechanism to the base context gives the first-order deformation calculus.

\section{First-order deformation spaces and the cotangent boundary calculus}
\label{sec:first-order-deformation-boundary}

A first-order deformation is a lifting problem in which the base is replaced
by one structured square-zero thickening. After identifying the reduction,
the geometry has the form
\[
\begin{tikzcd}[column sep=huge,row sep=large]
X \arrow[r,hook] \arrow[d,"f"']
    & X' \arrow[d]
\\
S \arrow[r,hook]
    & S'.
\end{tikzcd}
\]
Cotangent transitivity turns the existence of such a lift, its variation, and
its automorphisms into one exact boundary calculus. Throughout this section
\textbf{deformation} refers to this first-order square-zero situation.

\subsection{Affine deformation spaces}

\begin{definition}[Affine deformation over an infinitesimal base]
\label{def:affine-deformation-square-zero}
Let
\[
    A'\longrightarrow A
\]
be a structured square-zero extension of connective commutative ring spectra
by a connective \(A\)-module \(I\). Let \(B\) be a connective commutative
\(A\)-algebra. A \textbf{deformation of \(B\) over \(A'\)} is a pair
\((B',\alpha)\) consisting of a connective commutative \(A'\)-algebra \(B'\)
and an equivalence
\[
    \alpha:B'\otimes_{A'}A\xrightarrow{\simeq}B
\]
in commutative \(A\)-algebras. Morphisms are equivalences over \(A'\)
compatible with the specified reduction equivalences. The resulting
\(\infty\)-groupoid is denoted
\[
    \operatorname{Def}_{A'/A}(B).
\]
\end{definition}

\begin{theorem}[Affine deformation as a homotopy fibre]
\label{thm:affine-deformation-homotopy-fibre}
Suppose \(A'\to A\) is classified by
\[
    \eta_A:L_A\longrightarrow I[1],
\]
and put
\[
    J:=B\otimes_AI.
\]
Base change gives
\[
    \eta_B:
    L_A\otimes_AB
    \longrightarrow
    J[1].
\]
Then there is a natural equivalence
\[
\operatorname{Def}_{A'/A}(B)
\simeq
\operatorname{fib}_{\eta_B}
\left(
    \Map_{\Mod_B}(L_B,J[1])
    \longrightarrow
    \Map_{\Mod_B}(L_A\otimes_AB,J[1])
\right).
\]
Equivalently, a deformation is a morphism
\[
    \eta_e:L_B\longrightarrow J[1]
\]
together with a specified homotopy between its restriction to
\(L_A\otimes_AB\) and \(\eta_B\).
\end{theorem}

\begin{proof}
\leavevmode
\begin{enumerate}[label=\textup{(\arabic*)}]
    \item \textbf{Verify connectivity.}
    The class \(\eta_A:L_A\to I[1]\) is an object of the connective derivation category \(\operatorname{Der}^+\) of
    \cite{LurieDAGIV}, because \(A\) and \(I\) are connective.

    \item \textbf{Its associated algebra arrow.}
    Its associated algebra arrow is the structured square-zero extension \(A'\to A\).

    \item \textbf{Verify connectivity.}
    The connective deformation equivalence \cite[Proposition~3.22 and Remark~3.23]{LurieDAGIV}
    identifies connective \(A'\)-algebras with connective derivations over \(\eta_A\).

    \item \textbf{Verify connectivity.}
    Restrict this equivalence to those objects whose reduction is the fixed connective
    \(A\)-algebra \(B\).

    \item \textbf{Verify compatibility.}
    The resulting fibre consists precisely of derivation classes
    \[
        \eta_e:L_B\longrightarrow J[1]
    \]
    together with a specified homotopy making the square
    \[
    \begin{tikzcd}[column sep=large,row sep=large]
        L_A\otimes_AB \arrow[r,"\eta_B"] \arrow[d]
            & J[1] \arrow[d,equal]
        \\
        L_B \arrow[r,"\eta_e"']
            & J[1]
    \end{tikzcd}
    \]
    commute.

    \item \textbf{Apply the universal property.}
    By the universal property of a homotopy fibre, this is exactly
    \[
    \operatorname{fib}_{\eta_B}
    \left(
        \Map_{\Mod_B}(L_B,J[1])
        \longrightarrow
        \Map_{\Mod_B}(L_A\otimes_AB,J[1])
    \right).
    \]

    \item \textbf{Identify the equivalence.}
    The cited equivalence is an equivalence of \(\infty\)-categories, so taking the indicated
    fibre gives an equivalence of the full deformation \(\infty\)-groupoids, including all
    morphisms and higher homotopies.
\end{enumerate}
\end{proof}

\begin{theorem}[Affine cotangent boundary calculus]
\label{thm:affine-cotangent-obstruction}
In the situation of
\Ref{thm:affine-deformation-homotopy-fibre}, the transitivity cofibre sequence
\[
    L_A\otimes_AB
    \longrightarrow
    L_B
    \longrightarrow
    L_{B/A}
\]
determines a canonical boundary element
\[
    \operatorname{ob}(B/A')
    \in
    \pi_0\map_{\Mod_B}(L_{B/A},J[2]).
\]
It satisfies:
\begin{enumerate}[label=(\roman*)]
    \item \(\operatorname{ob}(B/A')=0\) if and only if
    \(\operatorname{Def}_{A'/A}(B)\) is nonempty;
    \item when nonempty, the deformation space is a torsor under
    \[
        \Omega^\infty\map_{\Mod_B}(L_{B/A},J[1]);
    \]
    \item for a fixed deformation \(\xi\), its automorphism space is naturally
    equivalent to
    \[
        \Omega^\infty\map_{\Mod_B}(L_{B/A},J).
    \]
\end{enumerate}
\end{theorem}

\begin{proof}
\leavevmode
\begin{enumerate}[label=\textup{(\arabic*)}]
    \item \textbf{Apply cotangent transitivity.}
    Apply the exact contravariant functor \(\map_{\Mod_B}(-,J[1])\) to the transitivity sequence.

    \item \textbf{Compute the fibre.}
    This gives a fibre sequence of spectra
    \[
    \begin{tikzcd}[column sep=small]
        \map(L_{B/A},J[1]) \arrow[r]
            & \map(L_B,J[1]) \arrow[r]
            & \map(L_A\otimes_AB,J[1]) \arrow[r,"\partial"]
            & \map(L_{B/A},J[2]).
    \end{tikzcd}
    \]

    \item \textbf{Establish the next comparison.}
    Define
    \[
        \operatorname{ob}(B/A'):=\partial(\eta_B).
    \]

    \item \textbf{Invoke the cited result.}
    By \Ref{thm:affine-deformation-homotopy-fibre}, a deformation is exactly a lift of the point
    \(\eta_B\) through the middle map.

    \item \textbf{Such a lift exists exactly when its boundary.}
    Such a lift exists exactly when its boundary is null.

    \item \textbf{Compute the fibre.}
    If the fibre is inhabited, translation by a chosen point identifies it with the homotopy
    fibre over zero, which is \(\Omega^\infty\map(L_{B/A},J[1])\).

    \item \textbf{Identify the map.}
    The automorphism space is the based loop space of the deformation \(\infty\)-groupoid at
    \(\xi\), giving the final formula.
\end{enumerate}
\end{proof}
The affine deformation groupoid is now expressed as a homotopy fibre and its
obstruction theory is the boundary map of cotangent transitivity. We next
globalize this description using the invariance of the small Zariski site
across the base thickening.

\subsection{Global deformation spaces}

\begin{definition}[Deformation of a spectral context]
\label{def:global-deformation-square-zero}
Let
\[
    i:S\longrightarrow S'
\]
be a structured infinitesimal extension of spectral schemes and let
\(f:X\to S\) be a spectral scheme. A \textbf{deformation of \(X\) over
\(S'\)} is a pair \((X',\alpha)\) consisting of a spectral scheme
\(X'\to S'\) and an equivalence
\[
    \alpha:X'\times_{S'}S\xrightarrow{\simeq}X
\]
over \(S\). Morphisms are equivalences over \(S'\) with the specified
coherent compatibility with \(\alpha\). Write
\[
    \operatorname{Def}_{S'/S}(X)
\]
for the resulting \(\infty\)-groupoid.
\end{definition}

\begin{lemma}[Fixed Zariski-site description of a deformation]
\label{lem:fixed-zariski-site-deformation}
Let \(i:S\to S'\) be a structured infinitesimal extension and let
\((X',\alpha)\in\operatorname{Def}_{S'/S}(X)\). Then the base-change map
\[
    j:X\simeq X'\times_{S'}S\longrightarrow X'
\]
is an underlying square-zero infinitesimal extension and induces an equivalence
of small spectral Zariski sites. Through this equivalence, the deformation is
equivalently a square-zero deformation of the native structure coefficients on
the fixed spectral Zariski site of \(X\), compatible with the corresponding
base deformation. More precisely, restriction to affine opens identifies
\(\operatorname{Def}_{S'/S}(X)\) with the descent groupoid of compatible
affine deformations over the corresponding affine restrictions of
\(S\to S'\).
\end{lemma}

\begin{proof}
\leavevmode
\begin{enumerate}[label=\textup{(\arabic*)}]
    \item \textbf{Invoke the cited result.}
    By \Ref{prop:global-square-zero-geometric-properties}\textup{(ii),(v)}, the structured base
    extension \(i:S\to S'\) is an underlying square-zero infinitesimal substitution.

    \item \textbf{Apply base change.}
    The morphism \(j\) is its geometric Cartesian base change.

    \item \textbf{Conclude the step.}
    Hence \Ref{prop:square-zero-substitution-permanence} makes \(j\) an underlying
    square-zero infinitesimal substitution and identifies the small spectral Zariski sites of
    \(X\) and \(X'\).

    \item \textbf{is affine.}
    In particular, \(j\) is affine.

    \item \textbf{We first pass from a global deformation.}
    We first pass from a global deformation to affine coefficient data.

    \item \textbf{Construct the factorization.}
    Choose an affine Zariski atlas
    \[
        \{T'_\alpha=\Spec(C'_\alpha)\to X'\}_\alpha
    \]
    and refine it so that each composite \(T'_\alpha\to S'\) factors through an affine open
    \[
        V'_\alpha=\Spec(A'_\alpha)\longrightarrow S'.
    \]

    \item \textbf{Fix the notation.}
    Put
    \[
        U_\alpha:=T'_\alpha\times_{X'}X,
        \qquad
        V_\alpha:=V'_\alpha\times_{S'}S.
    \]

    \item \textbf{Use the stated property.}
    Because both \(j:X\to X'\) and \(i:S\to S'\) are affine square-zero substitutions, \(U_\alpha\)
    and \(V_\alpha\) are affine.

    \item \textbf{Fix the notation.}
    Write
    \[
        U_\alpha=\Spec(C_\alpha),
        \qquad
        V_\alpha=\Spec(A_\alpha).
    \]

    \item \textbf{Verify connectivity.}
    The defining Cartesian squares give a commutative diagram of connective commutative ring
    spectra
    \[
    \begin{tikzcd}[column sep=large,row sep=large]
        A'_\alpha \arrow[r] \arrow[d]
            & C'_\alpha \arrow[d]
        \\
        A_\alpha \arrow[r]
            & C_\alpha
    \end{tikzcd}
    \]
    with a canonical equivalence
    \[
        C'_\alpha\otimes_{A'_\alpha}A_\alpha
        \simeq
        C_\alpha.
    \]

    \item \textbf{Conclude the step.}
    Thus every affine member is an affine deformation in the sense of
    \Ref{def:affine-deformation-square-zero}.

    \item \textbf{On overlaps and iterated overlaps, refine by.}
    On overlaps and iterated overlaps, refine by affine opens.

    \item \textbf{Apply base change.}
    The resulting affine deformation data agree under restriction because they are obtained by
    Cartesian base change from the single global object \(X'\).

    \item \textbf{Transport the structure.}
    Under the small-site equivalence, the opens \(U_\alpha\) form the corresponding affine basis
    of \(X\).

    \item \textbf{Prove the converse.}
    Conversely, suppose compatible affine deformation data are given on an affine basis of the
    fixed small Zariski site of \(X\), subordinate to compatible affine pairs for
    \(X\to S\).

    \item \textbf{Derive the comparison.}
    For a pair
    \[
        U=\Spec(B)\longrightarrow V=\Spec(A),
    \]
    the site equivalence for the structured base extension \(S\to S'\) gives the corresponding
    affine open
    \[
        V'=\Spec(A')\longrightarrow S'.
    \]

    \item \textbf{Verify connectivity.}
    The affine deformation datum is a connective commutative \(A'\)-algebra \(B'\)
    with
    \[
        B'\otimes_{A'}A\simeq B.
    \]

    \item \textbf{Apply base change.}
    The map \(B'\to B\) is the coordinate pushout corresponding to the geometric Cartesian base
    change of the structured square-zero map \(\Spec(A)\to\Spec(A')\).

    \item \textbf{Conclude the step.}
    Hence it is itself a structured square-zero extension by
    \Ref{thm:structured-square-zero-geometric-base-change}.

    \item \textbf{Restrict the construction.}
    If \((U_1,V_1)\to(U,V)\) is an affine refinement, the square-zero Zariski-restriction theorem
    \Ref{lem:square-zero-compatible-pair-restriction} identifies the corresponding \(\Spec(B'_1)\)
    with the affine open restriction of \(\Spec(B')\).

    \item \textbf{Transport the same comparison.}
    The same statement on all iterated refinements supplies the canonical higher pasting
    compatibilities.

    \item \textbf{Glue the local data.}
    Hence the affine deformations form a spectral Zariski gluing diagram and glue to a spectral
    scheme
    \[
        X'\longrightarrow S'.
    \]

    \item \textbf{Glue the local data.}
    Their specified reductions glue to an equivalence
    \[
        X'\times_{S'}S\xrightarrow{\simeq}X.
    \]

    \item \textbf{Use full faithfulness.}
    Full faithfulness of spectral Zariski descent shows that the space of such gluings is
    contractible and that morphisms and higher homotopies are obtained by the same descent.

    \item \textbf{Identify the equivalence.}
    The two constructions are therefore inverse equivalences of \(\infty\)-groupoids.

    \item \textbf{No affine atlas or coefficient presentation.}
    No affine atlas or coefficient presentation is retained after the construction.
\end{enumerate}
\end{proof}

\begin{theorem}[Global deformation as a homotopy fibre]
\label{thm:global-deformation-homotopy-fibre}
Let \(i:S\to S'\) be a structured infinitesimal extension by
\(I\in\QCoh(S)_{\geq0}\), classified by
\[
    \eta_S:L_S\longrightarrow I[1].
\]
Let \(f:X\to S\), put
\[
    J:=f^*I,
\]
and let
\[
    \eta_X:
    f^*L_S
    \longrightarrow
    J[1]
\]
be the pullback of \(\eta_S\). Then
\[
\operatorname{Def}_{S'/S}(X)
\simeq
\operatorname{fib}_{\eta_X}
\left(
    \Map_{\QCoh(X)}(L_X,J[1])
    \longrightarrow
    \Map_{\QCoh(X)}(f^*L_S,J[1])
\right).
\]
Equivalently, a deformation is a morphism
\[
    \eta_e:L_X\longrightarrow J[1]
\]
together with a specified homotopy
\[
    \eta_e|_{f^*L_S}\simeq\eta_X.
\]
\end{theorem}

\begin{proof}
\leavevmode
\begin{enumerate}[label=\textup{(\arabic*)}]
    \item \textbf{Choose the local data.}
    Choose an affine basis of compatible pairs
    \[
        U_\alpha=\Spec(B_\alpha)
        \longrightarrow
        V_\alpha=\Spec(A_\alpha)
    \]
    for \(X\to S\), closed under passage to a chosen affine basis on finite intersections.

    \item \textbf{structured extension transports each to its corresponding.}
    The structured extension \(S\to S'\) transports each \(V_\alpha\) to its corresponding affine
    square-zero thickening
    \[
        V'_\alpha=\Spec(A'_\alpha)\longrightarrow S'.
    \]

    \item \textbf{Apply Zariski descent.}
    By \Ref{lem:fixed-zariski-site-deformation}, restriction to this fixed basis identifies the
    global deformation groupoid with the descent limit of the affine deformation groupoids
    \[
        \operatorname{Def}_{A'_\alpha/A_\alpha}(B_\alpha)
    \]
    and their corresponding values on iterated affine intersections.

    \item \textbf{Compute the fibre.}
    On every affine member, \Ref{thm:affine-deformation-homotopy-fibre} identifies this
    deformation groupoid with the homotopy fibre over the restricted base class \(\eta_{X,\alpha}\) of
    \[
        \Map_{\Mod_{B_\alpha}}
        (L_{B_\alpha},J_\alpha[1])
        \longrightarrow
        \Map_{\Mod_{B_\alpha}}
        (L_{A_\alpha}\otimes_{A_\alpha}B_\alpha,J_\alpha[1]).
    \]

    \item \textbf{Apply cotangent base change.}
    Cotangent base change makes these affine equivalences compatible with every restriction and
    every iterated intersection.

    \item \textbf{Restrict the construction.}
    The affine characterization of the global cotangent complexes identifies the two resulting
    cosimplicial mapping diagrams with the restrictions of
    \[
        \Map_{\QCoh(X)}(L_X,J[1])
        \quad\text{and}\quad
        \Map_{\QCoh(X)}(f^*L_S,J[1]).
    \]

    \item \textbf{Apply Zariski descent.}
    Quasi-coherent Zariski descent computes each global mapping space as the limit of its affine
    mapping diagram.

    \item \textbf{Compute the fibre.}
    Since limits of spaces commute with homotopy fibres over compatible points, totalizing the
    affine deformation equivalences gives
    \[
    \operatorname{Def}_{S'/S}(X)
    \simeq
    \operatorname{fib}_{\eta_X}
    \left(
        \Map_{\QCoh(X)}(L_X,J[1])
        \longrightarrow
        \Map_{\QCoh(X)}(f^*L_S,J[1])
    \right).
    \]

    \item \textbf{Transport the same comparison.}
    The same descent calculation identifies morphisms and all higher homotopies, so this is an
    equivalence of \(\infty\)-groupoids.

    \item \textbf{Glue the local data.}
    Conversely, a point of the global homotopy fibre restricts to the compatible affine
    deformation data above, and the converse construction in
    \Ref{lem:fixed-zariski-site-deformation} glues them to the corresponding spectral scheme
    over \(S'\).

    \item \textbf{Glue the local data.}
    Thus no independent global deformation or gluing datum is assumed.
\end{enumerate}
\end{proof}

\begin{theorem}[Global cotangent boundary calculus]
\label{thm:global-cotangent-obstruction}
In the situation of
\Ref{thm:global-deformation-homotopy-fibre}, the global transitivity sequence
\[
    f^*L_S
    \longrightarrow
    L_X
    \longrightarrow
    L_{X/S}
\]
determines a canonical obstruction
\[
    \operatorname{ob}(X/S')
    \in
    \pi_0\map_{\QCoh(X)}(L_{X/S},J[2]).
\]
It satisfies:
\begin{enumerate}[label=(\roman*)]
    \item the obstruction vanishes if and only if \(X\) admits a deformation
    over \(S'\);
    \item when it vanishes, \(\operatorname{Def}_{S'/S}(X)\) is a torsor
    under
    \[
        \Omega^\infty\map_{\QCoh(X)}(L_{X/S},J[1]);
    \]
    \item for a chosen deformation \(\xi\),
    \[
        \operatorname{Aut}(\xi)
        \simeq
        \Omega^\infty\map_{\QCoh(X)}(L_{X/S},J).
    \]
\end{enumerate}
\end{theorem}

\begin{proof}
\leavevmode
\begin{enumerate}[label=\textup{(\arabic*)}]
    \item \textbf{Apply to the cited result.}
    Apply \(\map_{\QCoh(X)}(-,J[1])\) to \Ref{thm:global-cotangent-transitivity}.

    \item \textbf{Compute the fibre.}
    The resulting fibre sequence has boundary
    \[
        \partial:
        \map(f^*L_S,J[1])
        \longrightarrow
        \map(L_{X/S},J[2]).
    \]

    \item \textbf{Establish the next comparison.}
    Define
    \[
        \operatorname{ob}(X/S'):=\partial(\eta_X).
    \]

    \item \textbf{Now apply the homotopy-fibre description of the.}
    Now apply the homotopy-fibre description of \Ref{thm:global-deformation-homotopy-fibre}
    exactly as in the affine proof.
\end{enumerate}
\end{proof}

\begin{corollary}[Split first-order deformations]
\label{cor:split-first-order-deformations}
If the base infinitesimal extension is split, so that its classifying map
\(L_S\to I[1]\) is zero, then
\[
    \operatorname{Def}_{S'/S}(X)
    \simeq
    \Map_{\QCoh(X)}(L_{X/S},f^*I[1]).
\]
\end{corollary}

\begin{proof}
\leavevmode
\begin{enumerate}[label=\textup{(\arabic*)}]
    \item \textbf{deformation space.}
    The deformation space is the fibre over zero in \Ref{thm:global-deformation-homotopy-fibre}.

    \item \textbf{Apply cotangent transitivity.}
    Apply the mapping-spectrum fibre sequence associated to transitivity.
\end{enumerate}
\end{proof}
The deformation calculus now gives, both affinely and globally, the obstruction,
the deformation torsor, and the automorphism space from one cotangent exact
triangle. A second geometric use of the same differential structure is to
analyze contexts defined by equations.

\section{Equations, conormal classifiers, and first-order variation}
\label{sec:equations-first-order-variation}

Linear comprehension from Chapter~\Ref{chap:spectral-algebraic-geometry}
turns a coefficient into a context of linear functionals. Equalizing such a
term with the zero term gives a derived solution context. The relative
cotangent complex of that closed context records the linear type of the imposed
equations through its shifted conormal classifier.

\subsection{Conormal classifiers and quasi-smooth closed contexts}

\begin{definition}[Local finite presentation for spectral schemes]
\label{def:global-locally-finite-presentation}
A morphism \(f:Y\to X\) of spectral schemes is \textbf{locally of finite
presentation} if, for every affine open \(U=\Spec(A)\subseteq X\) and every
affine open \(V=\Spec(B)\subseteq Y\times_XU\), the connective commutative
ring spectrum \(B\) is locally of finite presentation over \(A\) in the sense
of Chapter~\Ref{chap:spectral-algebraic-geometry}. This is a property of the
morphism; no affine presentation is retained.
\end{definition}

\begin{proposition}[Locality and base change of local finite presentation]
\label{prop:global-lfp-locality}
Local finite presentation is local on source and target for the spectral
Zariski topology and is stable under arbitrary base change and composition.
For a closed immersion \(Z\hookrightarrow X\), it is enough to test the
property after pullback to affine opens of \(X\).
\end{proposition}

\begin{proof}
\leavevmode
\begin{enumerate}[label=\textup{(\arabic*)}]
    \item \textbf{Apply base change.}
    The affine notion is stable under base change and composition and is local on affine
    spectral Zariski covers.

    \item \textbf{Two affine computations admit common affine refinements by.}
    Two affine computations admit common affine refinements by the principal-open basis of
    Chapter~\Ref{chap:spectral-algebraic-geometry}; hence the definition is independent of
    charts.

    \item \textbf{Apply base change.}
    The locality and base-change assertions follow after passing to such refinements.

    \item \textbf{closed immersion.}
    A closed immersion is affine, so its pullback over an affine target is affine, giving the
    final assertion.
\end{enumerate}
\end{proof}

\begin{definition}[Conormal classifier]
\label{def:conormal-classifier}
Let
\[
    i:Z\hookrightarrow X
\]
be a closed immersion of spectral schemes. Its \textbf{conormal classifier} is
\[
    N^*_{Z/X}
    :=
    L_{Z/X}[-1]
    \in
    \QCoh(Z).
\]
\end{definition}

\begin{definition}[Quasi-smooth closed context extension]
\label{def:quasi-smooth-closed-context}
A closed immersion \(i:Z\hookrightarrow X\) is \textbf{quasi-smooth} if it is
locally of finite presentation and its conormal classifier
\(N^*_{Z/X}\) is finite locally free. For such an immersion define the normal
coefficient
\[
    N_{Z/X}:=(N^*_{Z/X})^\vee.
\]
Its associated normal linear context is the relative linear space of
Chapter~\Ref{chap:spectral-algebraic-geometry} determined by this finite
locally free coefficient.
\end{definition}

\begin{proposition}[Base change of quasi-smooth closed contexts]
\label{prop:quasi-smooth-base-change}
Quasi-smooth closed immersions are stable under arbitrary base change. If
\(Y\to X\) is any morphism and \(Z_Y:=Z\times_XY\), then
\[
    N^*_{Z_Y/Y}
    \simeq
    N^*_{Z/X}|_{Z_Y}.
\]
\end{proposition}

\begin{proof}
\leavevmode
\begin{enumerate}[label=\textup{(\arabic*)}]
    \item \textbf{Apply base change.}
    Closed immersions and local finite presentation are stable under base change.

    \item \textbf{cotangent comparison.}
    The cotangent comparison is \Ref{thm:global-cotangent-base-change}; shift by \([-1]\).

    \item \textbf{Use preservation.}
    Finite local freeness is preserved by substitution in the coefficient doctrine of
    Chapter~\Ref{chap:spectral-algebraic-geometry}.
\end{enumerate}
\end{proof}

\begin{lemma}[Ambient lifting of principal opens along a closed immersion]
\label{lem:closed-immersion-principal-open-lifting}
Let \(A\to B\) be a morphism of connective commutative ring spectra which is
surjective on \(\pi_0\). For every
\[
    g\in\pi_0(B)
\]
and every lift \(f\in\pi_0(A)\), there is a canonical equivalence
\[
    B\otimes_AA[f^{-1}]
    \simeq
    B[g^{-1}].
\]
Consequently the principal spectral open \(D_B^{\mathrm{sp}}(g)\) is the
Cartesian pullback of \(D_A^{\mathrm{sp}}(f)\). In particular, every
principal open neighbourhood in a closed affine subcontext is obtained by
restricting an ambient principal open neighbourhood.
\end{lemma}

\begin{proof}
\leavevmode
\begin{enumerate}[label=\textup{(\arabic*)}]
    \item \textbf{image of in.}
    The image of \(f\) in \(\pi_0(B)\) is \(g\).

    \item \textbf{Verify invertibility.}
    Both displayed \(B\)-algebras represent the universal connective commutative
    \(B\)-algebra in which this class is invertible.

    \item \textbf{Apply the universal property.}
    The universal property of spectral localization therefore supplies the canonical
    equivalence.

    \item \textbf{Derive the comparison.}
    Passing to opposite affine contexts gives the Cartesian statement.

    \item \textbf{No lift of.}
    No lift of \(g\) is retained: two lifts differ by an element of the kernel of
    \(\pi_0(A)\to\pi_0(B)\), and the resulting principal opens have the same pullback to \(\Spec(B)\).
\end{enumerate}
\end{proof}
The conormal classifier provides the intrinsic first-order coefficient of a
closed context. We now show that finite locally free linear equations produce
exactly the quasi-smooth closed contexts locally, so that the equation
presentation is a theorem rather than part of the structure.

\subsection{Derived zero loci as equational context formation}

\begin{construction}[Derived solution context of a linear equation]
\label{con:derived-zero-locus}
Let \(\mathcal E\in\QCoh(X)\) be finite locally free and let
\[
    s:\mathcal E\longrightarrow\mathcal O_X
\]
be a coefficient term. By the linear-functional representing property of
\Ref{rem:linear-coordinate-section-conventions}, \(s\) determines a section
\[
    \sigma_s:X\longrightarrow\mathbb L_X(\mathcal E).
\]
The zero morphism determines the zero section
\[
    0:X\longrightarrow\mathbb L_X(\mathcal E).
\]
Define the \textbf{derived solution context} or \textbf{derived zero locus} of
\(s\) by the Cartesian square
\[
\begin{tikzcd}[column sep=large, row sep=large]
    Z(s) \arrow[r] \arrow[d]
        & X \arrow[d,"\sigma_s"]
    \\
    X \arrow[r,"0"']
        & \mathbb L_X(\mathcal E).
\end{tikzcd}
\]
Thus \(Z(s)\) is the context obtained by imposing the equation \(s=0\).
\end{construction}

\begin{theorem}[Derived linear equations are quasi-smooth]
\label{thm:derived-zero-locus-quasi-smooth}
In \Ref{con:derived-zero-locus}, the projection
\[
    Z(s)\longrightarrow X
\]
is a quasi-smooth closed immersion, and there is a canonical equivalence
\[
    N^*_{Z(s)/X}
    \simeq
    \mathcal E|_{Z(s)}.
\]
\end{theorem}

\begin{proof}
\leavevmode
\begin{enumerate}[label=\textup{(\arabic*)}]
    \item \textbf{Set up the argument.}
    The assertion is Zariski-local on \(X\).

    \item \textbf{Take an affine open on which.}
    Take an affine open \(X=\Spec(A)\) on which
    \[
        \mathcal E\simeq\mathcal O_X^{\oplus n}.
    \]

    \item \textbf{Fix the notation.}
    Put
    \[
        R:=\operatorname{Sym}_A(A^{\oplus n}).
    \]

    \item \textbf{Identify the map.}
    The section \(s\) and the zero section correspond contravariantly to two
    \(A\)-algebra maps
    \[
        \epsilon_s,\epsilon_0:R\rightrightarrows A.
    \]

    \item \textbf{Fix the notation.}
    Write \(A_s\) and \(A_0\) for the two copies of \(A\) equipped with these
    \(R\)-algebra structures.

    \item \textbf{coordinate algebra of the zero locus.}
    The coordinate algebra of the zero locus is
    \[
        C:=A_s\otimes_RA_0.
    \]

    \item \textbf{Apply cotangent transitivity.}
    Apply transitivity to
    \[
        A\longrightarrow R\xrightarrow{\epsilon_0}A_0.
    \]

    \item \textbf{Derive the comparison.}
    The composite is the identity, hence \(L_{A_0/A}\simeq0\), while \Ref{prop:cotangent-free-algebra}
    gives
    \[
        L_{R/A}\simeq R\otimes_AA^{\oplus n}.
    \]

    \item \textbf{Conclude the step.}
    Therefore
    \[
        L_{A_0/R}
        \simeq
        A_0^{\oplus n}[1].
    \]

    \item \textbf{zero section.}
    The zero section is a closed immersion by \Ref{prop:affine-space-sections-closed}.

    \item \textbf{It is locally of finite presentation by.}
    It is locally of finite presentation by \Ref{thm:cotangent-finite-presentation}: on degree
    zero its coordinate map is the quotient by the ideal generated by the \(n\) universal
    variables, hence is finitely presented, and its relative cotangent complex is finite free.

    \item \textbf{Apply base change.}
    The projection \(\Spec(C)\to\Spec(A_s)\) is its base change along \(\sigma_s\), so it is a closed immersion
    locally of finite presentation.

    \item \textbf{Apply cotangent base change.}
    Cotangent base change identifies
    \[
        L_{C/A_s}
        \simeq
        L_{A_0/R}\otimes_{A_0}C
        \simeq
        C^{\oplus n}[1].
    \]

    \item \textbf{Restrict the construction.}
    Thus
    \[
        L_{C/A_s}[-1]\simeq C^{\oplus n},
    \]
    which is the restriction of \(\mathcal E\).

    \item \textbf{Glue the local data.}
    The affine identifications are natural under further Zariski localization and therefore glue
    globally.
\end{enumerate}
\end{proof}

\begin{lemma}[Universal equations and the conormal basis]
\label{lem:universal-equations-conormal-basis}
Let \(A\) be a connective commutative ring spectrum and put
\[
    R:=\operatorname{Sym}_A(A^{\oplus n}).
\]
Let
\[
    \epsilon_0,\epsilon_f:R\rightrightarrows A
\]
be respectively the zero augmentation and the map classified by points
\(f_1,\ldots,f_n\in\Omega^\infty A\). Write
\(A_{\epsilon_f}\) and \(A_{\epsilon_0}\) for the two copies of \(A\) with
these \(R\)-algebra structures and put
\[
    C:=A_{\epsilon_f}\otimes_R A_{\epsilon_0}.
\]
Regard \(C\) as an \(A\)-algebra through
\(A=A_{\epsilon_f}\to C\), and put
\[
    K_C:=\operatorname{fib}(A\to C).
\]
For each \(1\leq i\leq n\), the \(i\)-th universal equation together with
its canonical nullhomotopy in \(C\) determines a class
\[
    \lambda_i\in\pi_0(K_C).
\]
Under the canonical equivalence
\[
    L_{C/A}[-1]\simeq C^{\oplus n},
\]
the derived Hurewicz--conormal morphism
\[
    C\otimes_AK_C\longrightarrow L_{C/A}[-1]
\]
carries the image of \(\lambda_i\) to the \(i\)-th standard basis vector of
\(\pi_0(C)^{\oplus n}\).
\end{lemma}

\begin{proof}
\leavevmode
\begin{enumerate}[label=\textup{(\arabic*)}]
    \item \textbf{Fix the notation.}
    Let \(x_1,\ldots,x_n\) denote the universal generators of \(R=\operatorname{Sym}_A(A^{\oplus n})\).

    \item \textbf{Derive the comparison.}
    The universal derivation gives
    \[
        L_{R/A}\simeq R\otimes_AA^{\oplus n},
    \]
    under which \(dx_i\) is the \(i\)-th standard basis vector.

    \item \textbf{Apply cotangent transitivity.}
    Apply transitivity to
    \[
        A\longrightarrow R\xrightarrow{\epsilon_0}A_{\epsilon_0}.
    \]

    \item \textbf{Apply cotangent transitivity.}
    The composite is the identity on the underlying \(A\)-algebra, so \(L_{A_{\epsilon_0}/A}\simeq0\), and
    transitivity gives
    \[
        L_{A_{\epsilon_0}/R}[-1]
        \simeq
        L_{R/A}\otimes_RA_{\epsilon_0}
        \simeq
        A_{\epsilon_0}^{\oplus n}.
    \]

    \item \textbf{Apply cotangent transitivity.}
    The class determined by \(x_i\) and its nullhomotopy under the zero augmentation maps to
    the \(i\)-th basis vector; this is the boundary interpretation of the transitivity
    morphism applied to the universal derivation.

    \item \textbf{Apply cotangent base change.}
    The pushout square defining \(C\) and cotangent base change identify
    \[
        L_{C/A}
        \simeq
        L_{A_{\epsilon_0}/R}\otimes_{A_{\epsilon_0}}C.
    \]

    \item \textbf{Derive the comparison.}
    The canonical homotopy between the images of \(x_i\) through the two legs of the pushout
    gives \(\lambda_i\in\pi_0(K_C)\).

    \item \textbf{Verify naturality.}
    Naturality of the relative-cofibre-to-cotangent morphism identifies its image with the base
    change of the universal class for \(\epsilon_0\), hence with the \(i\)-th standard basis
    vector.
\end{enumerate}
\end{proof}

\begin{theorem}[Local equational characterization of quasi-smooth closed contexts]
\label{thm:local-zero-locus-characterization}
Let \(i:Z\hookrightarrow X\) be a closed immersion of spectral schemes. The
following are equivalent.
\begin{enumerate}[label=(\roman*)]
    \item \(i\) is quasi-smooth.
    \item Zariski-locally on \(X\), the immersion is equivalent to the derived
    zero locus of a term
    \[
        \mathcal E\longrightarrow\mathcal O_X
    \]
    for a finite locally free coefficient \(\mathcal E\).
\end{enumerate}
Under such an identification, the restricted coefficient
\(\mathcal E|_Z\) identifies canonically with \(N^*_{Z/X}\).
\end{theorem}

\begin{proof}
\leavevmode
\begin{enumerate}[label=\textup{(\arabic*)}]
    \item \textbf{Set up the argument.}
    The implication \((\mathrm{ii})\Rightarrow(\mathrm{i})\) is \Ref{thm:derived-zero-locus-quasi-smooth}.

    \item \textbf{Assume .}
    Assume \((\mathrm{i})\).

    \item \textbf{Establish the next comparison.}
    The assertion is local on \(X\).

    \item \textbf{On the complement of , the immersion.}
    On the complement of \(Z\), the immersion is empty, and the empty immersion is the
    derived zero locus of the unit term \(\mathcal O_X\to\mathcal O_X\): affine-locally its coordinate pushout has
    zero degree-zero ring and hence is the zero ring spectrum.

    \item \textbf{It therefore suffices to work near a.}
    It therefore suffices to work near a point of \(Z\).

    \item \textbf{Verify connectivity.}
    Choose an affine neighbourhood and write
    \[
            Z=\Spec(B)\longrightarrow X=\Spec(A),
        \]
    where \(A\to B\) is connective, surjective on \(\pi_0\), and locally of finite
    presentation.

    \item \textbf{Shrink the ambient affine around the selected.}
    Shrink the ambient affine around the selected point so that
    \[
            L_{B/A}[-1]\simeq B^{\oplus n}.
        \]

    \item \textbf{Use spectral localization.}
    This shrinking may be performed in the ambient context by
    \Ref{lem:closed-immersion-principal-open-lifting}: choose a principal open of \(\Spec(B)\) on
    which the conormal classifier is free, lift its defining class from \(\pi_0(B)\) to
    \(\pi_0(A)\), and replace the ambient affine by the corresponding principal localization.

    \item \textbf{Fix the notation.}
    Put
    \[
            K:=\operatorname{fib}(A\to B),
            \qquad
            R:=\pi_0(A),
            \qquad
            S:=\pi_0(B),
            \qquad
            I:=\ker(R\to S).
        \]

    \item \textbf{Derive the comparison.}
    The relative Hurewicz--conormal comparison gives
    \[
            \pi_0(B\otimes_AK)
            \xrightarrow{\cong}
            \pi_0(L_{B/A}[-1]).
        \]

    \item \textbf{Verify connectivity.}
    Because \(B\) and \(K\) are connective,
    \[
            \pi_0(B\otimes_AK)
            \simeq
            S\otimes_R\pi_0(K)
            \simeq
            \pi_0(K)/I\pi_0(K).
        \]

    \item \textbf{Verify surjectivity.}
    Hence \(\pi_0(K)\to\pi_0(L_{B/A}[-1])\) is surjective.

    \item \textbf{Fix the notation.}
    Let \(e_1,\ldots,e_n\) be the basis of \(\pi_0(L_{B/A}[-1])\simeq S^{\oplus n}\) determined by the chosen local trivialization.

    \item \textbf{Choose the local data.}
    Choose points
    \[
            \widetilde\kappa_i\in\Omega^\infty K
        \]
    whose components \(\kappa_i\in\pi_0(K)\) map to the \(e_i\).

    \item \textbf{Their images in.}
    Their images in \(\Omega^\infty A\) determine elements
    \[
            f_1,\ldots,f_n\in R
        \]
    together with specified nullhomotopies of their images in \(B\), so \(f_i\in I\).

    \item \textbf{Invoke the cited result.}
    By \Ref{lem:derived-conormal-classical-quotient}, the image of \(\kappa_i\) in \(I/I^2\) is
    the class of \(f_i\).

    \item \textbf{Use the stated property.}
    Since the \(\kappa_i\) map to a basis of the derived conormal module, the classes of
    \(f_1,\ldots,f_n\) generate \(I/I^2\).

    \item \textbf{We next prove that is finitely generated.}
    We next prove that \(I\) is finitely generated.

    \item \textbf{Verify connectivity.}
    Local finite presentation of \(A\to B\) means that \(B\) is compact in connective
    commutative \(A\)-algebras.

    \item \textbf{Use the adjunction.}
    If \(\{T_\alpha\}\) is a filtered diagram of ordinary commutative \(R\)-algebras,
    Eilenberg--Mac Lane formation commutes with the filtered colimit and the \(\pi_0\dashv H\)
    adjunction gives
    \[
        \begin{aligned}
            \Map_{\operatorname{CRing}_{R/}}
            \left(S,\operatorname*{colim}_\alpha T_\alpha\right)
            &\simeq
            \Map_{\CAlg(\Sp)_{A/}}
            \left(B,H\operatorname*{colim}_\alpha T_\alpha\right)
            \\
            &\simeq
            \operatorname*{colim}_\alpha
            \Map_{\CAlg(\Sp)_{A/}}(B,HT_\alpha)
            \\
            &\simeq
            \operatorname*{colim}_\alpha
            \Map_{\operatorname{CRing}_{R/}}(S,T_\alpha).
        \end{aligned}
        \]

    \item \textbf{Conclude the step.}
    Thus \(S\) is finitely presented over \(R\).

    \item \textbf{Use the stated property.}
    Since \(R\twoheadrightarrow
    S\), its kernel \(I\) is finitely generated.

    \item \textbf{Establish the next comparison.}
    Set
    \[
            Q:=I/(f_1,\ldots,f_n).
        \]

    \item \textbf{Derive the comparison.}
    The module \(Q\) is finitely generated, and generation of \(I/I^2\) gives
    \[
            Q=IQ.
        \]

    \item \textbf{Use spectral localization.}
    At the prime corresponding to the chosen point, \(I\) lies in the maximal ideal, so
    Nakayama's lemma gives \(Q=0\) after localization.

    \item \textbf{Use the stated property.}
    Since \(Q\) is finitely generated, shrink once more so that
    \[
            I=(f_1,\ldots,f_n).
        \]

    \item \textbf{Fix the notation.}
    Put
    \[
            P:=\operatorname{Sym}_A(A^{\oplus n}).
        \]

    \item \textbf{Fix the notation.}
    Let \(\epsilon_0:P\to A\) be the zero augmentation and \(\epsilon_f:P\to A\) the map classified by the selected
    points lifting \(f_1,\ldots,f_n\).

    \item \textbf{Distinguish the two -algebra structures by writing and ,.}
    Distinguish the two \(P\)-algebra structures by writing \(A_{\epsilon_f}\) and \(A_{\epsilon_0}\),
    and define
    \[
            C:=A_{\epsilon_f}\otimes_P A_{\epsilon_0}.
        \]

    \item \textbf{Apply the universal property.}
    The specified nullhomotopies carried by the points \(\widetilde\kappa_i\) identify the two composites
    \(P\rightrightarrows B\), so the pushout universal property gives an \(A\)-algebra morphism
    \[
            C\longrightarrow B.
        \]

    \item \textbf{Verify connectivity.}
    Connective pushouts give
    \[
            \pi_0(C)
            \simeq
            R/(f_1,\ldots,f_n)
            \simeq
            R/I
            \simeq
            S.
        \]

    \item \textbf{Conclude the step.}
    Hence \(C\to B\) is an isomorphism on \(\pi_0\).

    \item \textbf{Invoke the cited result.}
    By \Ref{thm:derived-zero-locus-quasi-smooth},
    \[
            L_{C/A}[-1]\simeq C^{\oplus n}.
        \]

    \item \textbf{Fix the notation.}
    Put
    \[
            K_C:=\operatorname{fib}(A\to C),
            \qquad
            K_B:=\operatorname{fib}(A\to B)=K.
        \]

    \item \textbf{For each universal equation, its canonical nullhomotopy in.}
    For each universal equation, its canonical nullhomotopy in \(C\) determines
    \(\lambda_i\in\pi_0(K_C)\).

    \item \textbf{Invoke the cited result.}
    By \Ref{lem:universal-equations-conormal-basis}, the derived conormal map sends \(\lambda_i\)
    to the \(i\)-th standard basis vector.

    \item \textbf{Transport the structure.}
    Under \(K_C\to K_B\), the class \(\lambda_i\) maps to \(\kappa_i\), because \(C\to B\) was
    constructed from precisely those nullhomotopies.

    \item \textbf{Verify naturality.}
    Naturality of the Hurewicz--conormal comparison gives a commutative square
    \[
        \begin{tikzcd}[column sep=huge, row sep=large]
            B\otimes_AK_C \arrow[r] \arrow[d]
                & (L_{C/A}\otimes_CB)[-1] \arrow[d]
            \\
            B\otimes_AK_B \arrow[r]
                & L_{B/A}[-1].
        \end{tikzcd}
        \]

    \item \textbf{Identify the map.}
    The right vertical map carries the standard basis to \(e_1,\ldots,e_n\) and therefore induces an
    isomorphism on \(\pi_0\).

    \item \textbf{Compare the two constructions.}
    Both its source and target are finite projective \(B\)-modules.

    \item \textbf{Use finite projectivity.}
    A morphism between finite projective \(B\)-modules which is an isomorphism on
    \(\pi_0\) is an equivalence, by the flat homotopy-group formula for finite projectives.

    \item \textbf{Conclude the step.}
    Therefore
    \[
            L_{C/A}\otimes_CB
            \xrightarrow{\simeq}
            L_{B/A}.
        \]

    \item \textbf{Apply cotangent transitivity.}
    Transitivity for \(A\to C\to B\) gives \(L_{B/C}\simeq0\).

    \item \textbf{Derive the comparison.}
    Together with \(\pi_0(C)\cong\pi_0(B)\), cotangent rigidity yields \(C\simeq B\).

    \item \textbf{Conclude the step.}
    Thus the original closed immersion is, near the selected point, the derived zero locus of
    the constructed finite tuple of equations.
\end{enumerate}
\end{proof}

\begin{remark}[No equation presentation is retained]
\label{rem:zero-locus-presentation-witness}
The local tuple \((f_1,\ldots,f_n)\), its chosen lifts from the relative fibre,
and the trivialization of the conormal classifier occur only in the proof of
\Ref{thm:local-zero-locus-characterization}. A quasi-smooth closed context is
not equipped with local equations or a chosen zero-locus presentation.
\end{remark}
The geometric branch has now reached its principal applications: cotangent
data controls infinitesimal lifting, first-order deformations, and local
equational presentations of quasi-smooth closed contexts. The final
substantive construction is a stable consequence of the endpoint semantics
itself.

\section{Stable reconstruction from exact endpoint semantics}
\label{sec:stable-bnv-differential-cohomology}

The geometric branch has produced the endpoint adjunction string before any
stable differential-cohomological terminology was introduced. We now
stabilize that string and isolate the exact categorical mechanism responsible
for its fracture pattern. The order is deliberate: first prove an abstract
fracture theorem for an exact adjoint triple with fully faithful middle
functor, then specialize it to the stabilized infinitesimal semantics, and
only afterward compare the resulting structure with established
differential-cohomological fracture patterns.

\subsection{Stabilization of the endpoint adjunction string}

\begin{definition}[Stable native and infinitesimal coefficient categories]
\label{def:stable-native-infinitesimal-categories}
Define
\[
    \mathcal D_{\Sp}
    :=
    \operatorname{Sp}\bigl((\mathcal H_{\Sp})_*\bigr),
    \qquad
    \mathcal D_{\inf}
    :=
    \operatorname{Sp}\bigl((\mathcal H_{\inf})_*\bigr).
\]
These are the stabilizations of the pointed native spectral Zariski topos and
the pointed finite-infinitesimal spectral Zariski topos. We call objects of
\(\mathcal D_{\inf}\) \textbf{stable infinitesimal coefficients}.
\end{definition}

\begin{theorem}[Stable endpoint adjunction string]
\label{thm:stable-bnv-adjoint-string}
The endpoint adjunction string stabilizes canonically to an exact adjoint
string
\[
\boxed{
    i_!^{\mathrm{st}}
    \dashv
    i^{*,\mathrm{st}}
    \dashv
    i_*^{\mathrm{st}}
    \dashv
    i^{!,\mathrm{st}}
    :
    \mathcal D_{\Sp}
    \rightleftarrows
    \mathcal D_{\inf}.
}
\]
The functors \(i_!^{\mathrm{st}}\) and \(i_*^{\mathrm{st}}\) are fully
faithful. Moreover \(i_*^{\mathrm{st}}\) preserves all small limits and
colimits. In particular,
\[
\boxed{
    i^{*,\mathrm{st}}
    \dashv
    i_*^{\mathrm{st}}
    \dashv
    i^{!,\mathrm{st}}
}
\]
is an exact adjoint triple with fully faithful middle functor.
\end{theorem}

\begin{proof}
\leavevmode
\begin{enumerate}[label=\textup{(\arabic*)}]
    \item \textbf{Invoke the cited result.}
    By \Ref{thm:infinitesimal-adjoint-string}, the unstable endpoint semantics is
    \[
        i_!\dashv i^*\dashv i_*\dashv i^!.
    \]

    \item \textbf{Use the adjunction.}
    The functor \(i_!=\operatorname{cod}^*\) preserves all limits and, being a left adjoint, all colimits.

    \item \textbf{Use the adjunction.}
    Restriction \(i^*\) preserves all limits and, as the left adjoint of \(i_*\), all
    colimits.

    \item \textbf{Use preservation.}
    The functor \(i_*=\operatorname{dom}^*\) preserves all limits and preserves all colimits by
    \Ref{thm:source-pullback-sheafification}.

    \item \textbf{Conclude the argument.}
    Finally \(i^!\) preserves all limits.

    \item \textbf{Use the adjunction.}
    Thus the left adjoint in each adjacent pair is left exact.

    \item \textbf{Use the adjunction.}
    Passage to pointed objects and then to spectrum objects is therefore functorial for each
    adjacent adjunction and yields the displayed string.

    \item \textbf{Apply the statement objectwise.}
    Every induced functor between the stable categories is exact.

    \item \textbf{Use full faithfulness.}
    Stabilization preserves the unit and counit equivalences expressing full faithfulness of
    \(i_!\) and \(i_*\), so the two stabilized embeddings remain fully faithful.

    \item \textbf{Use the adjunction.}
    Since \(i_*^{\mathrm{st}}\) has both a left and a right adjoint, it preserves all small colimits and
    limits.
\end{enumerate}
\end{proof}

\begin{definition}[Stable differential-cohesive modalities]
\label{def:stable-differential-cohesive-modalities}
Define exact endofunctors of \(\mathcal D_{\inf}\) by
\[
    \Re^{\mathrm{st}}
    :=
    i_!^{\mathrm{st}}i^{*,\mathrm{st}},
    \qquad
    \Im^{\mathrm{st}}
    :=
    i_*^{\mathrm{st}}i^{*,\mathrm{st}},
    \qquad
    \&^{\mathrm{st}}
    :=
    i_*^{\mathrm{st}}i^{!,\mathrm{st}}.
\]
\end{definition}

\begin{proposition}[Stable modal triple]
\label{prop:stable-differential-cohesive-modal-triple}
There is a canonical exact adjoint triple
\[
    \Re^{\mathrm{st}}
    \dashv
    \Im^{\mathrm{st}}
    \dashv
    \&^{\mathrm{st}}.
\]
All three endofunctors are idempotent. There are canonical transformations
\[
    \Re^{\mathrm{st}}E\longrightarrow E\longrightarrow\Im^{\mathrm{st}}E,
\]
and the essential images of \(\Im^{\mathrm{st}}\) and
\(\&^{\mathrm{st}}\) are the fully faithful copy of
\(\mathcal D_{\Sp}\) determined by \(i_*^{\mathrm{st}}\).
\end{proposition}

\begin{proof}
\leavevmode
\begin{enumerate}[label=\textup{(\arabic*)}]
    \item \textbf{Identify the resulting comparison.}
    This is the stabilized form of \Ref{thm:endpoint-cohesive-triple}.

    \item \textbf{Use full faithfulness.}
    The adjunctions follow by composing the stabilized endpoint adjunctions, and idempotence
    follows from full faithfulness.

    \item \textbf{common essential image claim uses.}
    The common essential image claim uses
    \[
        i^{*,\mathrm{st}}i_*^{\mathrm{st}}\simeq\id,
        \qquad
        i^{!,\mathrm{st}}i_*^{\mathrm{st}}\simeq\id.
    \]
\end{enumerate}
\end{proof}
Stabilization produces the exact adjoint triple needed for the final
reconstruction. Before naming any of its pieces, we isolate the general
fracture mechanism that follows solely from exactness, adjunction, and full
faithfulness.

\subsection{An exact-adjoint fracture theorem}

\begin{theorem}[Exact-adjoint fracture from a fully faithful middle functor]
\label{thm:exact-adjoint-fracture}
Let
\[
    L\dashv j\dashv R:
    \mathcal A\rightleftarrows\mathcal B
\]
be an exact adjoint triple of stable \(\infty\)-categories with \(j\) fully
faithful. For \(E\in\mathcal B\), let
\[
    \eta_E:E\longrightarrow jLE,
    \qquad
    \epsilon_E:jRE\longrightarrow E
\]
be the unit and counit, and define
\[
    \mathsf A(E):=\fib(\eta_E),
    \qquad
    \mathsf Z(E):=\cofib(\epsilon_E).
\]
Put
\[
    \mathcal B^{\mathrm{pure}}:=\ker(R).
\]
Then the following structures and equivalences are canonical.
\begin{enumerate}[label=(\roman*)]
    \item Under the equivalences \(Lj\simeq\id_{\mathcal A}\) and
    \(Rj\simeq\id_{\mathcal A}\), the two adjacent-mates morphisms
    \[
        L(\epsilon_E),\quad R(\eta_E):RE\longrightarrow LE
    \]
    are canonically equivalent. Write their common value as
    \(\chi_E:RE\to LE\).

    \item \(\mathsf Z(E)\in\mathcal B^{\mathrm{pure}}\), and
    \[
        \mathsf Z:\mathcal B\longrightarrow\mathcal B^{\mathrm{pure}}
    \]
    is an exact reflector left adjoint to the inclusion. Moreover
    \[
        L\mathsf A(E)\simeq0,
        \qquad
        \mathsf Z\mathsf A(E)\simeq\mathsf Z(E),
        \qquad
        \mathsf Z(jLE)\simeq0.
    \]

    \item Applying \(L\) to the cycle triangle gives the characteristic exact
    triangle
    \[
        RE\xrightarrow{\chi_E}LE
        \xrightarrow{\phi_E}L\mathsf Z(E).
    \]
    Applying the cycle reflector to \(\mathsf A(E)\) gives the upper exact
    triangle
    \[
        j\Omega L\mathsf Z(E)
        \longrightarrow
        \mathsf A(E)
        \longrightarrow
        \mathsf Z(E).
    \]

    \item The canonical square
    \[
    \begin{tikzcd}[column sep=large,row sep=large]
        E \arrow[r] \arrow[d,"\eta_E"']
            & \mathsf Z(E) \arrow[d,"\eta_{\mathsf Z(E)}"]
        \\
        jLE \arrow[r,"j\phi_E"']
            & jL\mathsf Z(E)
    \end{tikzcd}
    \]
    is Cartesian. Equivalently,
    \[
        E\simeq
        \mathsf Z(E)\times_{jL\mathsf Z(E)}jLE.
    \]

    \item Let \(\mathcal T_j\) be the pullback \(\infty\)-category whose
    objects are triples
    \[
        (U,P,\phi),
        \qquad
        U\in\mathcal A,
        \quad
        P\in\mathcal B^{\mathrm{pure}},
        \quad
        \phi:U\to LP.
    \]
    Then
    \[
        \mathcal B\xrightarrow{\simeq}\mathcal T_j,
        \qquad
        E\longmapsto(LE,\mathsf Z(E),\phi_E)
    \]
    is an equivalence of stable \(\infty\)-categories. An inverse is
    \[
        \operatorname{Diff}_j(U,P,\phi)
        :=
        P\times_{jLP}jU.
    \]
\end{enumerate}
If \(\mathcal A\) and \(\mathcal B\) are presentable and the adjoints are
accessible, then \(\mathcal B^{\mathrm{pure}}\) is presentable and
\(\mathsf Z\) is an accessible exact localization.
\end{theorem}

\begin{proof}
\leavevmode
\begin{enumerate}[label=\textup{(\arabic*)}]
    \item \textbf{Use full faithfulness.}
    Full faithfulness of \(j\) gives canonical equivalences
    \[
            Lj\simeq\id_{\mathcal A},
            \qquad
            Rj\simeq\id_{\mathcal A}.
        \]

    \item \textbf{Apply to the two candidate maps in (i).}
    Apply \(j\) to the two candidate maps in \textup{(i)}.

    \item \textbf{Verify naturality.}
    Naturality of the unit with respect to \(\epsilon_E:jRE\to E\) gives
    \[
            jL(\epsilon_E)\,\eta_{jRE}
            \simeq
            \eta_E\epsilon_E,
        \]
    while naturality of the counit with respect to \(\eta_E:E\to jLE\) gives
    \[
            \epsilon_{jLE}\,jR(\eta_E)
            \simeq
            \eta_E\epsilon_E.
        \]

    \item \textbf{Identify the equivalence.}
    The maps \(\eta_{jRE}\) and \(\epsilon_{jLE}\) are equivalences because their objects lie in the
    essential image of \(j\).

    \item \textbf{Conclude the step.}
    Thus the two candidate maps have canonically equivalent images under the fully faithful
    functor \(j\), proving \textup{(i)}.

    \item \textbf{Apply to.}
    Apply \(R\) to
    \[
            jRE\xrightarrow{\epsilon_E}E\longrightarrow\mathsf Z(E).
        \]

    \item \textbf{Handle the first component.}
    The first induced map is an equivalence, hence \(R\mathsf Z(E)\simeq0\).

    \item \textbf{If , then.}
    If \(P\in\ker(R)\), then
    \[
            \map_{\mathcal B}(jRE,P)
            \simeq
            \map_{\mathcal A}(RE,RP)
            \simeq0.
        \]

    \item \textbf{Compute the mapping object.}
    Mapping the cycle triangle into \(P\) gives
    \[
            \map_{\mathcal B}(\mathsf Z(E),P)
            \xrightarrow{\simeq}
            \map_{\mathcal B}(E,P),
        \]
    so \(\mathsf Z\) is the reflector onto \(\ker(R)\).

    \item \textbf{Use the stated property.}
    Since \(jR\) and the identity functor are exact, their functorial cofiber \(\mathsf Z\)
    is exact.

    \item \textbf{Applying to.}
    Applying \(L\) to
    \[
            \mathsf A(E)\longrightarrow E\xrightarrow{\eta_E}jLE
        \]
    shows \(L\mathsf A(E)\simeq0\).

    \item \textbf{Use exactness.}
    Since \(jLE\) lies in the image of \(j\), its cycle quotient vanishes; exactness
    of \(\mathsf Z\) applied to the same triangle therefore gives \(\mathsf Z\mathsf A(E)\simeq\mathsf Z(E)\).

    \item \textbf{Identify the resulting comparison.}
    This proves \textup{(ii)}.

    \item \textbf{Derive the comparison.}
    Applying \(L\) to the cycle triangle and using \(Lj\simeq\id\) yields
    \[
            RE\xrightarrow{\chi_E}LE\xrightarrow{\phi_E}L\mathsf Z(E).
        \]

    \item \textbf{Use exactness.}
    Applying \(R\) to the deformation triangle gives
    \[
            R\mathsf A(E)\longrightarrow RE\xrightarrow{\chi_E}LE,
        \]
    so exactness identifies
    \[
            R\mathsf A(E)\simeq\Omega L\mathsf Z(E).
        \]

    \item \textbf{cycle triangle of , together.}
    The cycle triangle of \(\mathsf A(E)\), together with \(\mathsf Z\mathsf A(E)\simeq\mathsf Z(E)\), is therefore the upper
    triangle of \textup{(iii)}.

    \item \textbf{Verify naturality.}
    For \textup{(iv)}, naturality of \(\eta\) makes the displayed square commute.

    \item \textbf{Identify the map.}
    The fibre of its upper horizontal map is \(jRE\).

    \item \textbf{Use exactness.}
    By the characteristic triangle and exactness of \(j\), the fibre of the lower
    horizontal map is also canonically \(jRE\).

    \item \textbf{Transport the structure.}
    Under these identifications the induced map of fibres is the unit equivalence on an object
    in the image of \(j\).

    \item \textbf{square in a stable -category.}
    A square in a stable \(\infty\)-category is Cartesian exactly when the induced map of
    horizontal fibres is an equivalence, proving the fracture formula.

    \item \textbf{For (v), let.}
    For \textup{(v)}, let
    \[
            E:=P\times_{jLP}jU.
        \]

    \item \textbf{Use exactness.}
    Exactness and \(Lj\simeq\id\) give
    \[
            LE\simeq LP\times_{LP}U\simeq U.
        \]

    \item \textbf{Use the stated property.}
    Since \(RP\simeq0\) and \(Rj\simeq\id\), applying \(R\) gives
    \[
            RE\simeq0\times_{LP}U\simeq\fib(\phi).
        \]

    \item \textbf{projection therefore.}
    The projection \(E\to P\) therefore has fibre
    \[
            j\fib(\phi)\simeq jRE.
        \]

    \item \textbf{Use full faithfulness.}
    Its composite with the counit \(jRE\to E\) becomes null after mapping to \(P\) because
    \(RP=0\); the resulting map from \(jRE\) to the fibre is an equivalence after
    applying \(R\), hence is an equivalence by full faithfulness on the image of
    \(j\).

    \item \textbf{Conclude the step.}
    Thus \(E\to P\) is the cycle quotient, so \(\mathsf Z(E)\simeq P\).

    \item \textbf{Identify the resulting object.}
    Applying \(L\) to this quotient identifies the characteristic map with the given
    \(\phi\).

    \item \textbf{Prove the converse.}
    Conversely, \textup{(iv)} reconstructs every \(E\) from its characteristic triple.

    \item \textbf{Identify the equivalence.}
    The two assignments are functorial and exact, hence give the asserted stable equivalence.

    \item \textbf{and are accessible and so is the.}
    In the presentable case, \(R\) and \(j\) are accessible and so is the functorial
    cofiber \(\mathsf Z\).

    \item \textbf{Use spectral localization.}
    An accessible reflective localization of a presentable \(\infty\)-category has presentable
    target, proving the last statement.
\end{enumerate}
\end{proof}
The abstract reconstruction mechanism is now internal to the chapter: it uses
only exactness, the adjacent adjunctions, and full faithfulness of the middle
functor. We may therefore specialize it to the stabilized endpoint semantics.
Only at this point do we attach BNV-type terminology to the resulting sectors
and characteristic data; that terminology records a later comparison and is
not part of their construction.

\subsection{The stabilized spectral instance and BNV-type sectors}

Put
\[
    L:=i^{*,\mathrm{st}},
    \qquad
    j:=i_*^{\mathrm{st}},
    \qquad
    R:=i^{!,\mathrm{st}}.
\]

\begin{lemma}[Adjacent-mates comparison]
\label{lem:stable-bnv-adjacent-mates}
For \(E\in\mathcal D_{\inf}\), the two maps
\[
    i^{*,\mathrm{st}}(\epsilon_E),
    \qquad
    i^{!,\mathrm{st}}(\eta_E):
    i^{!,\mathrm{st}}E\longrightarrow i^{*,\mathrm{st}}E
\]
are canonically equivalent. Their common value
\[
    \chi_E:i^{!,\mathrm{st}}E\longrightarrow i^{*,\mathrm{st}}E
\]
is the adjacent-mates morphism of
\Ref{thm:exact-adjoint-fracture}.
\end{lemma}

\begin{proof}
\leavevmode
\begin{enumerate}[label=\textup{(\arabic*)}]
    \item \textbf{Apply the cited result(i) to the stabilized endpoint triple.}
    Apply \Ref{thm:exact-adjoint-fracture}\textup{(i)} to the stabilized endpoint triple.
\end{enumerate}
\end{proof}

\begin{definition}[Stable BNV-type sectors]
\label{def:stable-bnv-sectors}
For \(E\in\mathcal D_{\inf}\), define
\[
    \mathcal A(E)
    :=
    \fib\bigl(E\xrightarrow{\eta_E}\Im^{\mathrm{st}}E\bigr),
    \qquad
    \mathcal Z(E)
    :=
    \cofib\bigl(\&^{\mathrm{st}}E\xrightarrow{\epsilon_E}E\bigr).
\]
The endpoint values
\[
    i^{*,\mathrm{st}}E,
    \qquad
    i^{!,\mathrm{st}}E
\]
are respectively the underlying and secondary sectors. The map
\[
    \phi_E:
    i^{*,\mathrm{st}}E
    \longrightarrow
    i^{*,\mathrm{st}}\mathcal Z(E)
\]
is the characteristic morphism obtained from the exact triangle
\[
    i^{!,\mathrm{st}}E
    \xrightarrow{\chi_E}
    i^{*,\mathrm{st}}E
    \xrightarrow{\phi_E}
    i^{*,\mathrm{st}}\mathcal Z(E).
\]
\end{definition}

\begin{definition}[BNV-pure cycle coefficient]
\label{def:stable-bnv-pure-cycle}
Define
\[
    \mathcal D_{\inf}^{\mathrm{BNV}}
    :=
    \ker(i^{!,\mathrm{st}})
    =
    \left\{
        P\in\mathcal D_{\inf}
        \ \middle|\
        i^{!,\mathrm{st}}P\simeq0
    \right\}.
\]
Objects of this full subcategory are called \textbf{BNV-pure cycle
coefficients}. The terminology records the comparison made below; the
subcategory itself is constructed solely from the endpoint right adjoint.
\end{definition}

\begin{proposition}[Cycle reflector and deformation identities]
\label{prop:stable-bnv-cycle-reflector}
The category \(\mathcal D_{\inf}^{\mathrm{BNV}}\) is stable and
\[
    \mathcal Z:
    \mathcal D_{\inf}\longrightarrow\mathcal D_{\inf}^{\mathrm{BNV}}
\]
is an exact reflector. Moreover
\[
    i^{!,\mathrm{st}}\mathcal Z(E)\simeq0,
    \qquad
    \mathcal Z^2(E)\simeq\mathcal Z(E),
\]
and
\[
\boxed{
    i^{*,\mathrm{st}}\mathcal A(E)\simeq0,
    \qquad
    \mathcal Z\mathcal A(E)\simeq\mathcal Z(E),
    \qquad
    \mathcal Z\Im^{\mathrm{st}}(E)\simeq0.
}
\]
\end{proposition}

\begin{proof}
\leavevmode
\begin{enumerate}[label=\textup{(\arabic*)}]
    \item \textbf{Identify the resulting comparison.}
    This is \Ref{thm:exact-adjoint-fracture}\textup{(ii)} for the stabilized endpoint triple.
\end{enumerate}
\end{proof}

\begin{proposition}[Presentability of the BNV-pure sector]
\label{prop:stable-bnv-pure-presentable}
The stable category \(\mathcal D_{\inf}^{\mathrm{BNV}}\) is presentable and
\(\mathcal Z\) is an accessible exact localization. For every small diagram
\(k\mapsto P_k\) of BNV-pure coefficients,
\[
    \operatorname*{colim}^{\mathcal D_{\inf}^{\mathrm{BNV}}}_kP_k
    \simeq
    \mathcal Z\left(
        \operatorname*{colim}^{\mathcal D_{\inf}}_kP_k
    \right).
\]
\end{proposition}

\begin{proof}
\leavevmode
\begin{enumerate}[label=\textup{(\arabic*)}]
    \item \textbf{presentability statement.}
    The presentability statement is the last clause of \Ref{thm:exact-adjoint-fracture}.

    \item \textbf{colimit formula.}
    The colimit formula is the standard formula for colimits in the target of a reflector.
\end{enumerate}
\end{proof}

\begin{theorem}[Upper BNV-type triangle]
\label{thm:stable-bnv-upper-triangle}
For every \(E\in\mathcal D_{\inf}\), there is a canonical exact triangle
\[
\boxed{
    i_*^{\mathrm{st}}\Omega i^{*,\mathrm{st}}\mathcal Z(E)
    \longrightarrow
    \mathcal A(E)
    \longrightarrow
    \mathcal Z(E).
}
\]
\end{theorem}

\begin{proof}
\leavevmode
\begin{enumerate}[label=\textup{(\arabic*)}]
    \item \textbf{Identify the resulting comparison.}
    This is \Ref{thm:exact-adjoint-fracture}\textup{(iii)}.
\end{enumerate}
\end{proof}
The stable sectors and characteristic morphism are now defined. The remaining
task is reconstruction: the fracture square must recover an arbitrary stable
infinitesimal coefficient from its underlying value and pure cycle data.

\subsection{Reconstruction and characteristic classification}

\begin{theorem}[Stable BNV-type reconstruction square]
\label{thm:stable-bnv-reconstruction}
For every \(E\in\mathcal D_{\inf}\), the canonical square
\[
\begin{tikzcd}[column sep=large,row sep=large]
    E \arrow[r] \arrow[d,"\eta_E"']
        & \mathcal Z(E) \arrow[d,"\eta_{\mathcal Z(E)}"]
    \\
    i_*^{\mathrm{st}}i^{*,\mathrm{st}}E
        \arrow[r,"i_*^{\mathrm{st}}\phi_E"']
        & i_*^{\mathrm{st}}i^{*,\mathrm{st}}\mathcal Z(E)
\end{tikzcd}
\]
is a pullback. Equivalently,
\[
\boxed{
    E
    \simeq
    \mathcal Z(E)
    \times_{i_*^{\mathrm{st}}i^{*,\mathrm{st}}\mathcal Z(E)}
    i_*^{\mathrm{st}}i^{*,\mathrm{st}}E.
}
\]
\end{theorem}

\begin{proof}
\leavevmode
\begin{enumerate}[label=\textup{(\arabic*)}]
    \item \textbf{Apply the cited result(iv).}
    Apply \Ref{thm:exact-adjoint-fracture}\textup{(iv)}.
\end{enumerate}
\end{proof}

\begin{corollary}[BNV-type differential-cohomological pullback of mapping spectra]
\label{cor:stable-bnv-differential-cohomology-pullback}
For \(K,E\in\mathcal D_{\inf}\), mapping out of \(K\) gives a canonical
homotopy pullback
\[
\begin{aligned}
    \map_{\mathcal D_{\inf}}(K,E)
    \simeq{}&
    \map_{\mathcal D_{\inf}}(K,\mathcal Z(E))
    \\
    &\times_{
        \map_{\mathcal D_{\inf}}
        (K,i_*^{\mathrm{st}}i^{*,\mathrm{st}}\mathcal Z(E))
    }
    \map_{\mathcal D_{\inf}}
        (K,i_*^{\mathrm{st}}i^{*,\mathrm{st}}E).
\end{aligned}
\]
Thus the graded groups
\[
    \widehat E^{\,n}(K)
    :=
    \pi_{-n}\map_{\mathcal D_{\inf}}(K,E)
\]
carry the corresponding BNV-type differential-cohomological fracture. No
identification with a geometric differential-cohomology theory is asserted
without further cycle or curvature data.
\end{corollary}

\begin{proof}
\leavevmode
\begin{enumerate}[label=\textup{(\arabic*)}]
    \item \textbf{Compute the mapping object.}
    Mapping spectra preserve limits in the second variable.

    \item \textbf{Apply them to the cited result.}
    Apply them to \Ref{thm:stable-bnv-reconstruction}.
\end{enumerate}
\end{proof}

\begin{definition}[Characteristic pullback category]
\label{def:stable-bnv-characteristic-category}
Let \(\mathcal T_{\inf}\) be the pullback \(\infty\)-category whose objects
are triples
\[
    (U,P,\phi),
    \qquad
    U\in\mathcal D_{\Sp},
    \quad
    P\in\mathcal D_{\inf}^{\mathrm{BNV}},
    \quad
    \phi:U\longrightarrow i^{*,\mathrm{st}}P.
\]
Equivalently it is the pullback
\[
\begin{tikzcd}[column sep=large,row sep=large]
    \mathcal T_{\inf} \arrow[r] \arrow[d]
        & \Fun(\Delta^1,\mathcal D_{\Sp})
          \arrow[d,"{(\operatorname{ev}_0,\operatorname{ev}_1)}"]
    \\
    \mathcal D_{\Sp}\times\mathcal D_{\inf}^{\mathrm{BNV}}
        \arrow[r,"{(\id,i^{*,\mathrm{st}})}"']
        & \mathcal D_{\Sp}\times\mathcal D_{\Sp}.
\end{tikzcd}
\]
All functors in this pullback are exact, so \(\mathcal T_{\inf}\) carries its
canonical stable \(\infty\)-category structure.
\end{definition}

\begin{construction}[Differential function object of characteristic data]
\label{con:stable-bnv-differential-function-object}
For \((U,P,\phi)\in\mathcal T_{\inf}\), define
\[
\boxed{
    \operatorname{Diff}(U,P,\phi)
    :=
    P\times_{\Im^{\mathrm{st}}P}i_*^{\mathrm{st}}U,
}
\]
where \(P\to\Im^{\mathrm{st}}P\) is the stable infinitesimal-shape unit and
\[
    i_*^{\mathrm{st}}U
    \xrightarrow{i_*^{\mathrm{st}}(\phi)}
    i_*^{\mathrm{st}}i^{*,\mathrm{st}}P
    =
    \Im^{\mathrm{st}}P
\]
is induced by \(\phi\).
\end{construction}

\begin{theorem}[BNV-type characteristic classification]
\label{thm:stable-bnv-characteristic-classification}
There is a canonical equivalence of stable \(\infty\)-categories
\[
\boxed{
    \mathcal D_{\inf}\simeq\mathcal T_{\inf}.
}
\]
It sends \(E\) to
\[
    \bigl(i^{*,\mathrm{st}}E,\mathcal Z(E),\phi_E\bigr),
\]
and its inverse is \(\operatorname{Diff}\). In particular,
\[
    i^{*,\mathrm{st}}\operatorname{Diff}(U,P,\phi)\simeq U,
    \qquad
    \mathcal Z\operatorname{Diff}(U,P,\phi)\simeq P,
\]
and the recovered characteristic morphism is \(\phi\).
\end{theorem}

\begin{proof}
\leavevmode
\begin{enumerate}[label=\textup{(\arabic*)}]
    \item \textbf{Identify the resulting comparison.}
    This is \Ref{thm:exact-adjoint-fracture}\textup{(v)} for
    \[
        i^{*,\mathrm{st}}\dashv i_*^{\mathrm{st}}\dashv i^{!,\mathrm{st}}.
    \]
\end{enumerate}
\end{proof}

\begin{proposition}[Colimits of characteristic triples]
\label{prop:stable-bnv-characteristic-colimits}
Let \(k\mapsto(U_k,P_k,\phi_k)\) be a small diagram in
\(\mathcal T_{\inf}\). Put
\[
    U:=\operatorname*{colim}^{\mathcal D_{\Sp}}_kU_k,
    \qquad
    P:=\operatorname*{colim}^{\mathcal D_{\inf}^{\mathrm{BNV}}}_kP_k.
\]
The cocone maps induce
\[
    \lambda_P:
    \operatorname*{colim}_k i^{*,\mathrm{st}}P_k
    \longrightarrow
    i^{*,\mathrm{st}}P,
\]
and the \(\phi_k\) induce
\[
    \overline\phi:
    U\longrightarrow\operatorname*{colim}_k i^{*,\mathrm{st}}P_k.
\]
Then
\[
    \operatorname*{colim}^{\mathcal T_{\inf}}_k(U_k,P_k,\phi_k)
    \simeq
    \bigl(U,P,\lambda_P\overline\phi\bigr).
\]
\end{proposition}

\begin{proof}
\leavevmode
\begin{enumerate}[label=\textup{(\arabic*)}]
    \item \textbf{Transport the diagram through the cited result.}
    Transport the diagram through \Ref{thm:stable-bnv-characteristic-classification}.

    \item \textbf{Use the adjunction.}
    The functor \(i^{*,\mathrm{st}}\) preserves colimits because it is a left adjoint, and \(\mathcal Z\)
    preserves colimits into the pure subcategory because it is the reflector.

    \item \textbf{Verify naturality.}
    Naturality of the cycle quotient identifies the resulting characteristic map with
    \(\lambda_P\overline\phi\).
\end{enumerate}
\end{proof}

\begin{remark}[Comparison with BNV and differential function spectra]
\label{rem:stable-bnv-differential-function-spectra}
The fracture square and characteristic classification above have the same
formal exact-adjunction shape as the characteristic pullback decomposition of
\cite[Proposition~3.5]{BunkeNikolausVoelkl}, and as the homotopy-pullback pattern of differential
function spectra in \cite{HopkinsSinger2005}. The comparison is made only
after the fracture theorem has been proved internally. Consequently the name
\textbf{stable BNV reconstruction} below is recognition terminology for this
constructed pattern; it does not supply an additional cycle, form, curvature,
period, or transfer datum.
\end{remark}

\begin{corollary}[Stable BNV reconstruction is intrinsic]
\label{cor:stable-bnv-intrinsic}
The stable BNV-type sectors, characteristic morphism, reconstruction square,
and characteristic classification are functorial outputs of the stabilized
endpoint adjunction string. No independent BNV package or differential
function datum is part of the input of a stable infinitesimal coefficient.
\end{corollary}

\begin{proof}
\leavevmode
\begin{enumerate}[label=\textup{(\arabic*)}]
    \item \textbf{They are the specialization of the cited.}
    They are the specialization of \Ref{thm:exact-adjoint-fracture} to the stabilized endpoint
    triple, followed only by the comparison terminology of the preceding remark.
\end{enumerate}
\end{proof}
The stable endpoint semantics is therefore determined by its characteristic
data. We close the formal development by collecting the chapter's export
interface. Before listing the exported structures, we fix the bookkeeping
principle that distinguishes constructed structure from proof presentations.

\section{Exported interface of the differentiated spectral doctrine}
\label{sec:differentiated-spectral-doctrine-interface}
\label{sec:differentiated-doctrine-audit}

The first part of the interface distinguishes intrinsic constructed structure
from auxiliary proof presentations. Constructions supplied functorially by
universal properties, descent, adjunctions, and exact calculus are retained as
structure; auxiliary presentations chosen only to prove their properties
remain proof witnesses.

\begin{theorem}[No differential or cohesive axiom disguised as presentation data]
\label{thm:differentiated-no-hidden-axioms}
None of the following is retained as additional input structure on a spectral
context, an underlying infinitesimal substitution, or an infinitesimal sheaf:
\begin{enumerate}[label=(\roman*)]
    \item a chosen affine atlas, affine refinement, or independently retained
    descent datum used to construct \(L_{X/S}\) or \(X_\eta\);
    \item a coefficient or derivation presenting an already specified
    underlying square-zero arrow;
    \item a finite square-zero factorization, nilpotent filtration, or numerical
    infinitesimal order attached to a finite infinitesimal substitution;
    \item a tuple of local equations presenting a quasi-smooth closed context;
    \item a presentation of the cofree cocommutative coalgebra \(C_X(M)\) by
    symmetric powers, homotopy fixed points, or any other chosen model;
    \item an independently specified monoidal exponential map, Seely
    comparison, primitive-coalgebra coherence, infinitesimal augmentation,
    codereliction, deriving transformation, or differential operator in a
    stable coefficient fibre, or a separate ordinary cofree-comonoid structure after
    passage to the homotopy category;
    \item an independently specified compatibility between substitution and the
    differential exponential: the actual comparison \(\beta_{f,M}\) and its
    comonad, monoidal, coalgebra, codereliction, and deriving compatibilities are
    constructed from pullback and cofreedom;
    \item an independently specified cohesion structure, reduction operation,
    reduction modality, infinitesimal-shape modality, infinitesimal-flat
    modality, Kock--Lawvere principle, microlinearity datum, or synthetic
    infinitesimal-object package;
    \item an independently specified BNV underlying sector, secondary sector,
    deformation sector, cycle sector, characteristic morphism,
    differential-function reconstruction datum, curvature datum, period map,
    differential-form model, or transfer structure on a stable infinitesimal
    coefficient;
    \item a Schlessinger-style lifting witness attached to formal smoothness or
    formal \(\acute{e}\)taleness;
    \item a chosen representative for a relative infinitesimal-site object
    beyond the actual diagram \(U\hookrightarrow T\) over \(X\to S\).
\end{enumerate}
The structures replacing these presentations are the cotangent representing
property, Zariski descent, the cofree coalgebra adjunction and its canonical
monoidal maps, the primitive coalgebra construction, the oplax substitution
comparison, the endpoint adjunctions and their stabilizations, exact fibre and
cofiber calculus, and their canonical composites.
\end{theorem}

\begin{proof}
\leavevmode
\begin{enumerate}[label=\textup{(\arabic*)}]
    \item \textbf{Each item.}
    Each item is a direct consequence of the construction that introduced it.

    \item \textbf{Apply Zariski descent.}
    Global cotangent objects and structured thickenings are determined by affine restriction and
    Zariski descent (\Ref{thm:global-cotangent-classifier},
    \Ref{con:global-structured-infinitesimal-extension}); underlying square-zero and finite
    infinitesimal arrows are defined by essential image and finite closure rather than by
    retained presentations (\Ref{def:underlying-spectral-square-zero-extension},
    \Ref{def:finite-spectral-infinitesimal-substitution}); and local equations are used only in
    the recognition theorem \Ref{thm:local-zero-locus-characterization}.

    \item \textbf{Verify functoriality.}
    In the linear branch, \(!_X\), its monoidal and Seely maps, the primitive coalgebra,
    infinitesimal augmentation, codereliction, deriving transformation, and substitution
    comparison are all constructed from the cofree adjunction, product preservation, opposite
    split square-zero formation, and their functoriality
    (\Ref{thm:cofree-coalgebra-linear-exponential}-- \Ref{prop:cofree-substitution-coherence}).

    \item \textbf{In the geometric branch, the cohesive modalities.}
    In the geometric branch, the cohesive modalities are composites of the endpoint quadruple
    (\Ref{thm:infinitesimal-adjoint-string} and
    \Ref{def:differential-cohesive-endpoint-constructions}), while formal etalification and
    infinitesimal disks are obtained from the left-exact reflector \(\Im\).

    \item \textbf{Use the adjunction.}
    Finally, the stable sectors, characteristic morphism, reconstruction square, and
    characteristic classification are obtained by specializing the exact-adjoint fracture
    theorem \Ref{thm:exact-adjoint-fracture} to the stabilized endpoint triple.

    \item \textbf{listed presentations therefore remain either theorem-domain properties.}
    The listed presentations therefore remain either theorem-domain properties or temporary
    proof witnesses, exactly as stated.
\end{enumerate}
\end{proof}

\begin{remark}[Genuine structured inputs]
\label{rem:differentiated-genuine-structured-inputs}
The audit preserves the structure that genuinely belongs to an object. A
structured infinitesimal extension by \(M\) genuinely contains its derivation
judgment \(L_{X/S}\to M[1]\); a relative context genuinely contains the map
\(X\to S\); an object of the relative infinitesimal site genuinely contains
the maps \(U\hookrightarrow X\), \(U\hookrightarrow T\), and \(T\to S\).
Presentations of already specified intrinsic objects and compatibility
witnesses supplied functorially remain proof-level data.
\end{remark}
With this bookkeeping convention fixed, the second part of the interface
records the differential structures that later chapters may use. The theorem
is deliberately a collecting statement; the conceptual retrospective is
reserved for the chapter summary.

\begin{theorem}[Differentiated spectral doctrine interface]
\label{thm:differentiated-spectral-doctrine-interface}
The constructions of this chapter canonically provide the following.
\begin{enumerate}[label=(\roman*)]
    \item \textbf{Affine tangent linearization.}
    For every commutative ring spectrum \(B\), the tangent stabilization of the
    ambient category of commutative spectral algebras at \(B\) is
    \[
        \operatorname{Stab}(\CAlg(\Sp)_{/B})\simeq\Mod_B,
    \]
    together with the split square-zero zeroth-space construction. Globally,
    \(\QCoh(X)\) is the descended stable coefficient fibre; no additional
    global stabilization equivalence is asserted.

    \item \textbf{Cotangent-classified differentiation.}
    Relative affine derivations are represented by \(L_{B/A}\). Globally
    \((X\to S)\mapsto L_{X/S}\) is a connective Cartesian section under
    represented substitution and satisfies the global cotangent chain rule.
    On affine relative contexts, the cotangent classifier additionally
    satisfies the Hurewicz--conormal comparison, cotangent rigidity, and the
    stated criterion for local finite presentation.

    \item \textbf{Differential-linear stable coefficient structure.}
    Every stable coefficient fibre \(\QCoh(X)\) carries the cofree
    cocommutative-coalgebra exponential \(!_X\), its canonically constructed
    monoidal coalgebra-modality maps and Seely equivalences, together with the
    canonical identification showing that these induce the same ordinary
    storage-modality and additive-bialgebra structure, primitive infinitesimal
    coalgebras, coherent infinitesimal augmentation, and canonical
    codereliction and deriving maps. Their images in the homotopy
    category give the ordinary differential-storage codereliction and deriving
    transformation. The actual comparison
    \(\beta_{f,M}:f^*(!_XM)\to !_Y(f^*M)\) makes these structures oplaxly
    compatible with substitution, including the comonad, monoidal, coalgebra,
    codereliction, and deriving maps. On finite projective linear models the
    resulting convolution derivation is canonically compatible along \(j_M\)
    with the spectral universal derivation.

    \item \textbf{Structured square-zero context extension.}
    Every connective coefficient and degree-one derivation
    \[
        \eta:L_{X/S}\longrightarrow M[1]
    \]
    constructs a spectral thickening \(X\hookrightarrow X_\eta\), stable under
    represented relative-base change. Algebraic pullback, relative-base
    change, and geometric Cartesian base change are separately constructed,
    and geometric Cartesian base change preserves structured square-zero
    extensions.

    \item \textbf{Presentation-free finite infinitesimal contexts.}
    Underlying square-zero substitutions and their finite-composition closure
    are affine properties of arrows, stable under base change, nilpotent on
    degree-zero coordinate rings, and invariant on underlying Zariski support
    and small Zariski sites. A structured representative may be chosen only as
    a proof witness and can be refined relative to a displayed base when a
    lifting argument requires it.

    \item \textbf{Infinitesimal Zariski semantics.}
    The category \(\Aff_{\Sp}^{\inf}\), its codomain-pullback Zariski topology,
    its relative affine-basis sites, and its subcanonical sheaf category
    \(\mathcal H_{\inf}\), together with the endpoint adjunctions
    \[
        i_!\dashv i^*\dashv i_*\dashv i^!.
    \]

    \item \textbf{Differential cohesion.}
    The endpoint quadruple is recognized as differential cohesion and induces
    the standard modal triple
    \[
        \Re\dashv\Im\dashv\&,
    \]
    where \(\Re\) is reduction, \(\Im\) is infinitesimal shape, and \(\&\)
    is infinitesimal flat. The modality \(\Im\) is a left-exact reflector.

    \item \textbf{Formal etalification and infinitesimal disks.}
    For \(g:E\to B\) in \(\mathcal H_{\inf}\), relative infinitesimal shape is
    \[
        \Im_BE=B\times_{\Im B}\Im E.
    \]
    The unit \(E\to\Im_BE\) is the formally \(\acute{e}\)tale reflection in
    the slice. For \(f:X\to S\), its kernel-pair \v{C}ech nerve gives the
    relative infinitesimal disk groupoid \(T_{\inf}^{[\bullet]}(X/S)\).

    \item \textbf{Cotangent lifting and formal geometry.}
    On an affine structured square-zero stage, the fibre of the modal lifting
    comparison \(\Theta_f\) is the nullhomotopy space of the canonical
    cotangent boundary. For an arbitrary represented structured square-zero
    stage, the represented comparison
    \(\Theta_f^{\mathrm{rep}}\) has the same cotangent description.
    One-stage detection extends the affine modal criterion to finite
    infinitesimal contexts. For represented spectral contexts,
    \[
        f\text{ formally }\acute{e}\text{tale}
        \quad\Longleftrightarrow\quad
        L_{X/S}\simeq0,
    \]
    while finite local freeness of \(L_{X/S}\) implies formal smoothness.
    The first structured neighbourhood of the diagonal maps canonically into
    \(T_{\inf}(X/S)\), and the difference of its two retractions is the
    universal relative derivation.

    \item \textbf{First-order deformation boundary calculus.}
    First-order deformation spaces are homotopy fibres of restriction on
    derivations; cotangent transitivity constructs their obstruction,
    deformation torsor, and automorphism space.

    \item \textbf{Equational context formation.}
    Derived zero loci of finite locally free linear equations are quasi-smooth
    closed contexts, and every quasi-smooth closed immersion is locally of this
    form. The conormal classifier is the restricted coefficient of the local
    equation context.

    \item \textbf{Stable BNV differential-cohomological reconstruction.}
    Stabilization of the endpoint semantics gives stable presentable
    categories
    \[
        \mathcal D_{\Sp}
        =
        \operatorname{Sp}((\mathcal H_{\Sp})_*),
        \qquad
        \mathcal D_{\inf}
        =
        \operatorname{Sp}((\mathcal H_{\inf})_*),
    \]
    and an exact adjoint triple
    \[
        i^{*,\mathrm{st}}
        \dashv
        i_*^{\mathrm{st}}
        \dashv
        i^{!,\mathrm{st}}.
    \]
    The abstract exact-adjoint fracture theorem
    \Ref{thm:exact-adjoint-fracture} applies to this fully faithful middle
    functor. Its stabilized modalities
    \[
        \Re^{\mathrm{st}},
        \qquad
        \Im^{\mathrm{st}},
        \qquad
        \&^{\mathrm{st}}
    \]
    construct the deformation and cycle sectors
    \[
        \mathcal A(E)
        =
        \fib(E\to\Im^{\mathrm{st}}E),
        \qquad
        \mathcal Z(E)
        =
        \cofib(\&^{\mathrm{st}}E\to E),
    \]
    while the already defined endpoint values
    \[
        i^{*,\mathrm{st}}E,
        \qquad
        i^{!,\mathrm{st}}E,
        \qquad
        i^{*,\mathrm{st}}\mathcal Z(E)
    \]
    serve as the underlying, secondary, and underlying-cycle sectors without
    further renaming. The characteristic morphism is
    \[
        \phi_E:
        i^{*,\mathrm{st}}E
        \longrightarrow
        i^{*,\mathrm{st}}\mathcal Z(E),
    \]
    and the canonical reconstruction is
    \[
        E
        \simeq
        \mathcal Z(E)
        \times_{i_*^{\mathrm{st}}i^{*,\mathrm{st}}\mathcal Z(E)}
        i_*^{\mathrm{st}}i^{*,\mathrm{st}}E.
    \]
    Equivalently, stable infinitesimal coefficients are classified by
    characteristic triples
    \[
        (U,P,\phi),
        \qquad
        \phi:U\longrightarrow i^{*,\mathrm{st}}P,
    \]
    with \(P\) BNV-pure, meaning
    \[
        i^{!,\mathrm{st}}P\simeq0,
    \]
    and mapping spectra satisfy the corresponding BNV-type
    differential-cohomological pullback. No identification with a geometric
    differential-cohomology theory and no separate geometric interpretation of
    the pure cycle coefficient is included in this abstract reconstruction.

    \item \textbf{Logical audit.}
    The constructions above distinguish genuine structured inputs from
    functorial outputs, theorem-domain properties, and temporary proof
    witnesses as in \Ref{thm:differentiated-no-hidden-axioms}.
\end{enumerate}
No additional differential, infinitesimal, cohesive, differential-cohomology,
presentation, admissibility, or compatibility axiom is introduced by this
collecting theorem.
\end{theorem}

\begin{proof}
\leavevmode
\begin{enumerate}[label=\textup{(\arabic*)}]
    \item \textbf{Identify the resulting comparison.}
    This theorem only collects results already proved.

    \item \textbf{Use the adjunction.}
    Items \textup{(i)}--\textup{(iii)} are the tangent, cotangent, and differential-linear
    constructions; items \textup{(iv)}--\textup{(vi)} are the square-zero, finite-infinitesimal,
    and infinitesimal-site constructions; items \textup{(vii)}--\textup{(ix)} are the endpoint
    modalities and cotangent lifting calculus; items \textup{(x)} and \textup{(xi)} are the
    deformation and equation calculi; item \textup{(xii)} is the stabilized exact-adjoint
    fracture construction together with its subsequent BNV-type comparison; and item
    \textup{(xiii)} is the bookkeeping audit immediately above.

    \item \textbf{No new construction.}
    No new construction is made here.
\end{enumerate}
\end{proof}

\section{Chapter summary}
\label{sec:differentiated-spectral-doctrine-summary}

The chapter has exposed two complementary differential mechanisms already
carried by the spectral doctrine. Affinely, the stable coefficient fibre is
the tangent stabilization of commutative spectral algebra, and relative
derivations in that tangent semantics are represented by the cotangent
classifier with its substitution law and chain rule. Independently, the same
stable coefficient fibres carry a cofree cocommutative-coalgebra
differential-exponential calculus. On free finite-projective models the two
descriptions meet: the convolution derivative determined by the cofree
calculus is compatible with the universal spectral K\"ahler derivation.

The cotangent classifier then drives the nonlinear geometric branch.
Degree-one cotangent classes produce structured square-zero extensions, whose
presentation-free finite composites form the infinitesimal test geometry.
Their Zariski topos yields the endpoint adjunctions and hence differential
cohesion. The resulting formal lifting comparison is computed by the same
cotangent classifier, so infinitesimal lifting, first-order deformation, and
local equation formation appear as different instances of one cotangent
boundary calculus.

Finally, stabilization converts the already constructed endpoint semantics
into an exact adjoint triple, and the resulting fracture theorem reconstructs
stable infinitesimal coefficients from characteristic data. Its BNV-type
differential-cohomological form is recognized after this reconstruction has
been obtained internally. The
distinction inherited from Chapter~\Ref{chap:spectral-algebraic-geometry}
remains active throughout: geometry is connective, coefficients are fully
stable, and every passage between the nonlinear, infinitesimal, and stable
layers is supplied by a construction established in this chapter.

